\documentclass[a4paper, 10pt]{article}

\usepackage{packageList}
\usepackage{macros}

%

\graphicspath{{./Figures/}}
\pgfplotsset{compat=1.11}
\newlength\fwidth

%
\title{Sharp inverse statements for kernel interpolation}

%
%
%
%
%
%
%

\author[1,2]{Tizian Wenzel \thanks{wenzel@math.lmu.de}}
\affil[1]{Institute for Applied Analysis and Numerical Simulation, University of Stuttgart, Germany}
\affil[2]{Department of Mathematics, Ludwig-Maximilians-Universität München, Germany}


\DeclareMathOperator{\spn}{span}
\DeclareMathOperator{\vol}{vol}


\DeclareMathOperator{\VN}{\tilde{V}_{N}}
\DeclareMathOperator{\proj}{\Pi_{\tilde{V}_{N}}}
\newcommand{\Ht}{\mathcal{H}_{\vartheta}}
\newcommand{\HT}{\mathcal{H}_{\Theta}}
\newcommand{\Nt}{{\tilde{N}}}

\newtheorem{conj}{Conjecture}

%
\newcommand{\tw}[1]{\textcolor{magenta}{TW: #1}}
\newcommand{\gs}[1]{\textcolor{blue}{GS: #1}}

%
\usepackage{comment}
%

%
\newif\iflong			
%
\longfalse

\newif\ifappendix
%
\appendixfalse

\usepackage{comment}
%

\usepackage[normalem]{ulem}			%
\usepackage{thm-restate}

\makeatletter
\newcommand*\bigcdot{\mathpalette\bigcdot@{.5}}
\newcommand*\bigcdot@[2]{\mathbin{\vcenter{\hbox{\scalebox{#2}{$\m@th#1\bullet$}}}}}
\makeatother


\begin{document}

\maketitle %



  
\begin{abstract}
%
%
%
%
%
%

%
While direct statements for kernel based interpolation on regions $\Omega \subset \R^d$ are well researched, 
far less is known about corresponding inverse statements.
The available inverse statements for kernel based interpolation so far are not sharp.

In this paper,
we derive sharp inverse statements for interpolation using finitely smooth kernels,
such as popular radial basis function (RBF) kernels like the class of Matérn or Wendland kernels.
In particular, the results show that there is a one-to-one correspondence between the smoothness of a function and its approximation rate via kernel interpolation: 
If a function can be approximated with a given rate, it has a corresponding smoothness and vice versa. \newline

%
%
%

%
%
\textbf{MSC classification:} 65D05, 46E22, 41A27, 41A17

%
\textbf{Keywords:} Kernel interpolation $\cdot$ Inverse estimates $\cdot$ Radial basis functions
\end{abstract}

%
%
%
%
%
%
%
%
 

%
%
%
%
%
%
%
%

%
 
%
%
%
%
%
%
%
%
%
%
%
%
%
%
%
%
%
%
%
%
%
%
%
%
%
%
%
%
%
%
%
%
%
%
%
%
%
%
%
%
%
%

%

\section{Introduction} \label{sec:introduction}

In approximation theory, direct statements give error estimates for the approximation of an element (e.g.\ a function) %
using previously defined basis elements, e.g.\ polynomials.
The derived convergence rate usually depends on properties of the element to be approximated, e.g.\ the smoothness of the function.
Inverse statements address the contrary and derive properties of the element which is approximated (e.g.\ the smoothness of the function) based on the approximation rate.


This paper deals with the particular case of function approximation via kernel interpolation, 
see e.g.~\cite{wendland2005scattered, fasshauer2007meshfree, fasshauer2015kernel, steinwart2008support} for a general introduction to kernel methods:
Given a nonempty set $\Omega$, a kernel is defined as a symmetric function $k: \Omega \times \Omega \rightarrow \R$.
While the set $\Omega$ can be in general quite arbitrary, we are interested in the frequent case of compact regions $\Omega \subset \R^d$. \\
For a given set of $n \in \N$ points $X_n \subset \Omega$, the kernel matrix $A_{X_n}$ is defined as $(A_{X_n})_{ij} = (k(x_i, x_j))_{ij} \in \R^{n \times n}$.
If the kernel matrix is positive definite for any $n$ pairwise distinct points $X_n \subset \Omega$ and any $n \in \N$, then the kernel is called strictly positive definite.
For such strictly positive definite kernels, there exists a unique associated Hilbert space $\ns$ of real-valued functions, the so called \textit{reproducing kernel Hilbert space} (RKHS), 
which is also called \textit{native space}.
In this space the kernel $k$ acts as a reproducing kernel, i.e.\ 
\begin{enumerate}
\item $k(\cdot, x) \in \ns ~ \forall x \in \Omega$,
\item $f(x) = \langle f, k(\cdot, x) \rangle_{\ns} ~ \forall x \in \Omega,  \forall f \in \ns$ \hfill (reproducing property).
\end{enumerate}
Kernels can be used for various tasks, and here we are concerned with the approximation of continuous functions by interpolation using continuous strictly positive definite kernels:
Given a continuous function $f \in \mathcal{C}(\Omega)$ and pairwise distinct interpolation points $X_n \subset \Omega$, we seek an interpolating function $s_{f, X_n}$ such that
\begin{align}
\label{eq:interpolation_conditions}
s_{f, X_n}(x_i) = f(x_i) \quad \forall x_i \in X_n.
\end{align}
A well known representer theorem shows
%
that there exists a minimum norm interpolant in $\ns$ of the form
\begin{align*}
s_{f, X_n}(\cdot) = \sum_{j=1}^n \alpha_j^{(n)} k(\cdot, x_j) \in \ns.
\end{align*}
The coefficients $\alpha_j^{(n)} \in \R, j= 1, \dots,n$ can be calculated by solving the linear equation system arising from Eq.~\eqref{eq:interpolation_conditions}, 
which is always solvable due to the assumption on the strict positive definiteness of the kernel.
Furthermore,
based on the reproducing property, it is easy to see that it holds $\Vert s_{f, X_n} \Vert_{\ns}^2 = (\alpha^{(n)})^\top A_{X_n} \alpha^{(n)}$.

%
%
%
%
%

%
%
%
%
%
%
%
%
%
%
%
%
%
%
%
%
%
Under suitable conditions on the kernel $k$ and on the compact region $\Omega \subset \R^d$ (such as a Lipschitz boundary), 
the RKHS can be characterized in terms of Sobolev spaces, i.e.\ there is a norm-equivalence $\mathcal{H}_k(\Omega) \asymp H^\tau(\Omega)$ for some $\tau > d/2$.
This is frequently the case, e.g.\ for radial basis function (RBF) kernels such as the popular classes of Matérn or Wendland kernels,
and will be the focus of the current paper.

In this case of $\ns \asymp H^\tau(\Omega)$ for some $\tau > d/2$, 
the approximation error between the function $f$ and the interpolant $s_{f, X}$ is usually bounded with help of the fill distance $h_{X, \Omega}$ and uniformity constant $\rho_X$ (defined in Eq.~\eqref{eq:fill_sep_dist} and Eq.~\eqref{eq:uniformly_distributed_points} respectively).
For $f \in H^\beta(\Omega) \supseteq H^\tau(\Omega)$ with $d/2 < \beta \leq \tau$, 
these estimates (see \Cref{th:error_estimate} for a precise statement) read
\begin{align}
\label{eq:direct_statement}
\Vert f - s_{f,X} \Vert_{L^2(\Omega)} \leq C h_{X, \Omega}^\beta \, \rho_{X}^{\tau-\beta} \Vert f \Vert_{H^\beta(\Omega)},
\end{align}
providing a convergence rate of $\beta$ in the fill distance $h_{X, \Omega}$ for uniformly distributed points $X \subset \Omega$.
This convergence rate matches the smoothness of the function $f \in H^\beta(\Omega)$.

In the current work, we prove corresponding inverse statements, 
which allow to derive the smoothness of a function $f \in C(\Omega)$ based on the rate of approximation. 
To be precise, we have the following main result of this paper, which will be proven in \Cref{sec:inverse_statement}:

\begin{restatable}{theorem}{mainresult}[Main result]
\label{th:inverse_statement_generalized}
Consider a compact Lipschitz region $\Omega \subset \R^d$ and a continuous kernel $k$ such that $\ns \asymp H^\tau(\Omega)$ for some $\tau > d/2$.
Consider $f \in \mathcal{C}(\Omega)$ such that for some $c_f, \beta, h_0 > 0$ it holds
\begin{align}
\label{eq:inverse_statement_assumption}
\Vert f - s_{f, X} \Vert_{L^2(\Omega)} \leq c_f h_{X, \Omega}^\beta
\end{align}
for all $X \subset \Omega$ with $h_{X, \Omega} \leq h_0$ and $\rho_X \leq 44$ (quasi-uniformly distributed). \\
If $\beta \in (0, \tau]$,
then $f \in H^{\beta'}(\Omega)$ for all $\beta' \in (0, \beta)$.
If $\beta > \tau$, then $f \in \ns \asymp H^\tau(\Omega)$. 
\end{restatable}

This result significantly extends the current state of the art compared to previously known inverse statements for kernel interpolation, 
see \Cref{subsec:related_work} and \Cref{sec:discussion_remarks} for a detailed discussion.
In particular this inverse statement \Cref{th:inverse_statement_generalized} is sharp, 
because it provides the correct smoothness in view of the corresponding direct statement Eq.~\eqref{eq:direct_statement} (up to an arbitrary small $\varepsilon > 0$).
Thus we obtain for the first time a one-to-one correspondence between the smoothness of a function and the rate of approximation for kernel interpolation.



%
%
%
%
%
%
%
%



%
%
%
%
%

%
%
%
%
%
%
%
%
%
%
%
%
%
%
%
%
%
%
%
%
%
%
%








%
%
%
%
%
%
%
%
%
%
%


%


%
%
%
%
%
%
%
%
%
%
%
%
%

The remaining paper is structured as follows: 
\Cref{sec:background} collects necessary background information on the tools which will be used. 
\Cref{sec:utility_statements} provides utility statements, 
which will then be leveraged in \Cref{sec:inverse_statement} to prove the main result \Cref{th:inverse_statement_generalized}.
%
%
\Cref{sec:discussion_remarks} discusses the results and elaborates on open directions for future research,
while \Cref{sec:conclusion} concludes the paper.




%


%
%
%
%
%
%
%
%
%
%
%
%
%
%
%
%
%
%
%
%
%
%
%
%
%
%
%
%
%
%
%
%
%
%
%
%
%




\section{Background}
\label{sec:background}
%

In the following, we collect some further required background information on geometric properties of $\Omega$ (\Cref{subsec:geometric_properties}) and Sobolev spaces (\Cref{subsec:sobolev_spaces}), 
kernels and kernel interpolation (\Cref{subsec:kernel_approx}) 
and review related literature (\Cref{subsec:related_work}).

%
%


\subsection{Geometric properties related to $\Omega$} \label{subsec:geometric_properties}

A standard assumption on the regularity of the region $\Omega \subset \R^d$, 
which is also frequently used in the following, is given by an interior cone condition.
We note that a Lipschitz boundary already implies the interior cone condition. 
%
\begin{definition}
A set $\Omega \subset \R^d$ is said to satisfy an interior cone condition, if there exists an angle $\alpha \in (0, \pi/2)$ and a radius $r>0$ such that for every $x \in \Omega$ a unit vector $\xi(x)$ exists such that the cone 
\begin{align}
\label{eq:definition_cone}
C(x, \xi(x), \alpha, r) := \{ x + \lambda y: y \in \R^d, \Vert y \Vert_2 = 1, y^\top \xi(x) \geq \cos(\alpha), \lambda \in [0, r] \}
\end{align}
is contained in $\Omega$.
\end{definition}
%

In order to quantify the distribution of finitely many scattered points $X \subset \Omega$ in a region $\Omega \subset \R^d$, 
we recall the definition of separation distance $q_{X}$ and fill distance $h_{X, \Omega}$ as
\begin{align}
\label{eq:fill_sep_dist}
\begin{aligned}
q_{X} :=& \frac{1}{2} \min_{x_i \neq x_j \in X} \Vert x_i - x_j \Vert_2, \\
h_X := h_{X, \Omega} :=& \sup_{x \in \Omega} \min_{x_j \in X} \Vert x - x_j \Vert_2.
\end{aligned}
\end{align}
We call a sequence $(X_n)_{n \in \N}$ of sets of points $X_n \subset \Omega$ quasi-uniformly distributed, 
if the uniformity constant $\rho_{X_n} := \rho_{X_n, \Omega} := h_{X_n, \Omega} / q_{X_n}$ is uniformly bounded for all $n \in N$, i.e.\
\begin{align}
\label{eq:uniformly_distributed_points}
\exists C>0 ~ \forall n \in \N \quad \rho_{X_n} \equiv \frac{h_{X_n, \Omega}}{q_{X_n}} \leq C < \infty.
\end{align}
%
%
Also a sequence of points $( x_n )_{n \in \N}$ is called quasi-uniformly distributed, 
if the sequence $( \{x_j\}_{j=1}^n )_{n \in \N}$ is quasi-uniformly distributed.

Let $C_d$ denote the volume of the unit ball $B_1(0) \subset \R^d$ and let $C_{d, \alpha}$ be the volume of the unit cone $C(x, \xi(x), \alpha, 1)$ (which is independent of $x \in \Omega$), i.e.\
\begin{align}
\label{eq:constants_ball_cone}
\begin{aligned}
C_d &:= \vol(B_1(0)), \\
C_{d, \alpha} &:= \vol(C(x, \xi(x), \alpha, 1)).
\end{aligned}
\end{align}
Given a bounded region $\Omega \subset \R^d$ which satisfies an interior cone condition, 
and $n \in \N$ points $X_n \subset \Omega$, 
it holds
%
%
\begin{align}
\label{eq:geometric_bounds}
\begin{aligned}
h_{X_n, \Omega} &\geq c_{\Omega} \cdot n^{-1/d}, \\
q_{X_n} &\leq C_{\Omega} \cdot n^{-1/d},
\end{aligned}
\end{align}
with constants
\begin{align*}
c_\Omega &= \frac{1}{\pi^{1/2}} \cdot \left( \vol(\Omega) \Gamma \left( \frac{d}{2} + 1 \right) \right)^{1/d}, \\
C_\Omega &= \left( \frac{2\pi}{\alpha} \right)^{1/d} \cdot \frac{1}{\pi^{1/2}} \cdot \left( \vol(\Omega) \Gamma \left( \frac{d}{2} + 1 \right) \right)^{1/d},
\end{align*}
(see \cite[Satz 2.1.6]{mueller2009komplexitaet} and \cite[Satz 2.1.7]{mueller2009komplexitaet}), especially it holds $c_\Omega < C_\Omega$. 
In the following,
we use the shorthand notation $h_X := h_{X, \Omega}$, as the region $\Omega$ is usually fixed and clear from the context.



A convenient way to obtain well distributed points for bounded regions $\Omega \subset \R^d$ is to use the \textit{geometric greedy algorithm}, 
which was introduced and analyzed in \cite{marchi2005optimal}.
We waive to recall the geometric greedy algorithm in itself, we simply recall that the resulting sequence of points satisfies for all $n \geq 2$ the convenient property \cite[Lemma 5.1]{marchi2005optimal}
\begin{align}
\label{eq:geometric_uniformity}
\frac{1}{2} h_{X_n, \Omega} \leq \frac{1}{2} h_{X_{n-1}, \Omega} \leq q_{X_n} \leq h_{X_n, \Omega}.
\end{align}
%
%
The constructions within the next sections will frequently make use of these points:
As soon as a region $\Omega$ is fixed, 
we consider \textit{the} sequence of geometric greedy points, 
though it is clear that there can be several such sequences, e.g.\ depending on the choice of the first point.
However, this will be used as notion to enforce that we think all the time about the same sequence of points.







\subsection{Sobolev spaces}
\label{subsec:sobolev_spaces}

We use standard notation \cite{agranovich2015sobolev}: 
The (fractional) Sobolev spaces $H^\tau(\R^d)$ can be defined for also non-integer $\tau \geq 0$ with help of the decay of the Fourier transform $\hat{u}$ of a function $u \in L^2(\R^d)$ as
\begin{align}
\begin{aligned}
\label{eq:definition_sobolev_Rd}
H^\tau(\R^d) &= \{ u \in L^2(\R^d) ~ | ~ \hat{u}(\cdot)(1 + \Vert \cdot \Vert_2^2)^{\tau/2} \in L^2(\R^d) \}, \\
\Vert u \Vert_{H^\tau(\R^d)}^2 &= \int_{\R^d} (1+\Vert \omega \Vert_2^2)^{\tau} \cdot |\hat{u}(\omega)|^2 ~ \mathrm{d}\omega.
\end{aligned}
\end{align}
For regions $\Omega \subset \R^d$, the localized Sobolev spaces $H^\tau(\Omega)$ can be defined via restrictions as
\begin{align*}
H^\tau(\Omega) &= \{ u|_\Omega ~ | ~ u \in H^\tau(\R^d) \}, \\
\Vert u \Vert_{H^\tau(\Omega)} &= \inf \{\Vert v \Vert_{H^\tau(\R^d)} ~ | ~ v \in H^\tau(\R^d), v|_\Omega = u \}.
\end{align*}
We remark that there are other definitions of these Sobolev spaces, for example with help of weak derivatives.
These yield norm-equivalent spaces under suitable conditions on $\Omega$, such as Lipschitz boundary, 
see e.g.\ \cite{triebel2002function, rychkov1999restrictions}.
We recall the well known Sobolev embedding theorems, 
which states for $\tau > d/2$ that $H^\tau(\Omega) \hookrightarrow \mathcal{C}(\Omega)$,
i.e.\ for any equivalence class within $H^\tau(\Omega)$ there is a continuous representer.

The following result allows to estimate intermediate Sobolev norms:

\begin{prop}
\label{prop:hoelder_interpol}
Let $\Omega \subset \R^d$ be a bounded Lipschitz domain, $\tau > 0$ and $\theta \in (0, 1)$.
Then there exists a constant $C > 0$ independent of $f$, such that
\begin{align}
\label{eq:sobolev_interpol_inequality}
\Vert f \Vert_{H^{\theta \tau}(\Omega)} \leq C \cdot \Vert f \Vert_{L^2(\Omega)}^{1-\theta} \cdot \Vert f \Vert_{H^\tau(\Omega)}^\theta \qquad \forall f \in H^\tau(\Omega).
\end{align}
\end{prop}

\begin{proof}
This is the Gagliardo–Nirenberg inequality, a well known result within the theory of interpolation of Sobolev spaces, 
see e.g.\ \cite[Theorem 1]{brezis2018gagliardo}.
%
\end{proof}


\subsection{Kernels and kernel interpolation}
\label{subsec:kernel_approx}

Additional to the notation introduced in \Cref{sec:introduction}, 
we recall further standard results on kernel interpolation, see mostly \cite{wendland2005scattered}:

An important class of kernels is given by \textit{translational invariant kernels}, 
for which there exists a function $\Phi: \R^d \rightarrow \R$ such that it holds
\begin{align*}
k(x, z) = \Phi(x - z) ~ \text{for all} ~ x, z \in \R^d.
\end{align*}
A particular instance of this type of kernels is given by radial basis function kernels $k$ (RBF kernels).
These can be represented with help of a \textit{radial basis function} $\varphi: \R_{\geq 0} \rightarrow \R$ as
\begin{align*}
k(x, z) = \varphi(\Vert x - z \Vert_2).
\end{align*}
Such a radial basis function kernel $k$ is again a translational invariant kernel with $\Phi(x) = \varphi(\Vert x \Vert_2)$.
These kernels are popular in applications due to their easy implementation independent of the dimension $d$ and region $\Omega$.

Translational invariant kernels and in particular RBF kernels can be characterized in terms of the Fourier transform $\hat{\Phi}$.
A frequent assumption is that it holds
\begin{align}
\label{eq:fourier_decay}
c_\Phi (1+\Vert \omega \Vert_2^2)^{-\tau} \leq \hat{\Phi}(\omega) \leq C_\Phi (1 + \Vert \omega \Vert_2^2)^{-\tau}
\end{align}
for some constants $c_\Phi, C_\Phi > 0$.
In view of the definiton of Sobolev spaces in \Cref{subsec:sobolev_spaces}, 
it follows that the RKHS $\ns$ of such kernels is norm-equivalent to the Sobolev space $H^\tau(\Omega)$ for Lipschitz regions $\Omega$, i.e.\ $\ns \asymp H^\tau(\Omega)$.
This motivates the notion of a \textit{kernel of finite smoothness $\tau > d/2$}.
Our main result \Cref{th:inverse_statement_generalized} will be formulated for such kernels of finite smoothness $\tau > d/2$, 
thus including in particular RBF kernels satisfying the decay from Eq.~\eqref{eq:fourier_decay}.

%
%

%


We continue by providing a detailed statement of the error estimate stated in Eq.~\eqref{eq:direct_statement},
which allows to bound the interpolation error for functions that are possibly outside the RKHS $\ns \asymp H^\tau(\Omega)$
with help of the fill distance $h_X$ and the uniformity constant $\rho_X$:

\begin{theorem}
\label{th:error_estimate}
Consider a RBF kernel $k$ that satisfies Eq.\ \eqref{eq:fourier_decay} for some $\tau > d/2$ and a compact Lipschitz region $\Omega \subset \R^d$.
For some $\beta$ with $d/2 < \beta \leq \tau$ let $f \in H^\beta(\Omega) \supseteq H^\tau(\Omega)$.
Then it holds
\begin{align*}
\Vert f - s_{f,X} \Vert_{L^2(\Omega)} \leq C h_{X}^\beta \rho_{X}^{\tau-\beta} \Vert f \Vert_{H^\beta(\Omega)}.
\end{align*}
\end{theorem}

\begin{proof}
This is a special case of \cite[Theorem 4.2]{narcowich2006sobolev} using $\mu = 0$ (with notation adopted to the current paper),
while the improvement $\beta > d/2$ (instead of $\lfloor \beta \rfloor > d/2$) is due to \cite[Theorem 4.1]{arcangeli2007extension}.
\end{proof}

For the proof of our corresponding inverse statement, we require stability statements for kernel interpolation:

\begin{theorem}
\label{th:stability_statement}
%
Let $k$ be a continuous kernel such that $\ns \asymp H^\tau(\Omega)$ for some Lipschitz region $\Omega \subset \R^d$.
Then there exists a constant $c_0 > 0$ such that for any set of pairwise distinct points $X \subset \Omega$ the kernel matrix $A_X$ satisfies
\begin{align*}
\lambda_{\min}(A_{X}) \geq c_0 q_{X}^{2\tau - d}
\end{align*}
and thus for the matrix norm $\Vert \cdot \Vert_{2, 2}$ it holds
\begin{align*}
\Vert A_X^{-1} \Vert_{2,2} \leq c_0^{-1} q_{X}^{d - 2\tau}.
\end{align*}
\end{theorem}

\begin{proof}
For the case of RBF kernels satisfying Eq.~\eqref{eq:fourier_decay}, which means $\ns \asymp H^\tau(\Omega)$, this is a well known results, see e.g.\ \cite{schaback1995error} or \cite[Chapter 12.2]{wendland2005scattered}.
The general case (where the kernel is not necessarily an RBF kernel) then holds due to the norm-equivalence $\ns \asymp H^\tau(\Omega)$, see e.g.~\cite[Section F.2]{haas2023mind}.
%
\end{proof}


Additionally to the ``continuous" $L^2(\Omega)$ norms we consider the \textit{discrete} $L^2(X)$ norm for some $X \subset \Omega$, i.e.\ for $u \in \mathcal{C}(\Omega)$ %
%
\begin{align}
\label{eq:discrete_L2_norm}
\Vert u \Vert_{L^2(X)} := \left( \frac{1}{|X|} \sum_{x \in X} |u(x)|^2 \right)^{1/2}.
\end{align}

The following additional proposition allows to bound the ${\Vert \cdot \Vert_{\ns}}$ norm of an kernel interpolant via a matrix norm and a discrete $L^2(X)$ norm.
The proof is based on kernel interpolation properties and we include it here for convenience:

\begin{prop}
\label{prop:upper_bound_ns_norm_via_L2}
Let $X := \{x_1, ..., x_n\} \subset \Omega$ be pairwise distinct. 
The norm $\Vert s_X \Vert_{\ns}$ of $s_X := \sum_{j=1}^n \alpha_j k(\cdot, x_j)$ can be estimated via the discrete $L^2(X)$ norm as
\begin{align*}
\Vert s_X \Vert_{\ns}^2 &\leq \Vert A_X^{-1} \Vert_{2,2} \cdot |X| \cdot \Vert s_X \Vert_{L^2(X)}^2.
\end{align*}
\end{prop}

\begin{proof}
We can calculate using $\alpha = (\alpha_1, ..., \alpha_n)^\top \in \R^n$:
\begin{align*}
\Vert s_X \Vert_{\ns}^2 &= \alpha^\top A_{X} \alpha = \alpha^\top A_{X} A_{X}^{-1} A_{X} \alpha 
\leq \Vert A_X^{-1} \Vert_{2,2} \cdot \Vert A_X \alpha \Vert_{2}^2 \\
& = \Vert A_X^{-1} \Vert_{2,2} \cdot |X| \cdot \Vert \sum_{j=1}^n \alpha_j k(\cdot, x_j) \Vert_{L^2(X)}^2 = \Vert A_X^{-1} \Vert_{2,2} \cdot |X| \cdot \Vert s_X \Vert_{L^2(X)}^2
\end{align*}
\end{proof}




%

%
%
%
%
%
%
%
%
%
%
%
%
%
%
%
%
%
%
%





%
%
%
%
%
%
%


%
%
%
%
%
%
%
%
%
%
%
%
%
%
%
%
%
%
%
%
%
%
%
%
%
%
%
%
%
%
%
%
%
%
%
%
%
%
%
%
%
%
%
%
%
%
%
%
%
%
%
%
%
%
%
%
%
%
%
%
%
%
%
%
%
%
%
%
%
%
%
%
%


%
\subsection{Related work} \label{subsec:related_work}


%
%
%
%
%
%




Direct and inverse statements in approximation theory are often associated with Jackson and Bernstein theorems.
The direct statements provide rates of approximation under certain assumptions of the element to be approximated,
while the inverse statements aim at deriving properties of that element based on the rate of approximation.
For classical approximation methods like polynomial, Fourier or spline approximation,
both direct and inverse statements are well researched \cite{butzer1969fundamental, devore1973saturation}.
For kernel interpolation, direct statements are also well researched \cite{wendland2005approximate, narcowich2006sobolev}, 
%
though the theory for corresponding inverse statements is less complete:

The standard way to prove inverse statements %
is to make use of Bernstein inequalities, 
which usually bound strong norms with help of a weaker norm. %
A typical Bernstein inequality for kernel-based interpolation on a region $\Omega \subseteq \R^d$ for a translational invariant kernel $k(x, z) = \Phi(x-z)$ satisfying Eq.~\eqref{eq:fourier_decay} and a set of pairwise distinct points $X_n \subset \Omega \subset \R^d$ reads e.g.\
\begin{align}
\label{eq:bernstein_not_available}
\left \Vert \sum_{j=1}^n \alpha_j \Phi(\cdot - x_j) \right \Vert_{H^\beta(\Omega)} \leq C q_{X_n}^{-\beta} \cdot \left \Vert \sum_{j=1}^n \alpha_j  \Phi(\cdot - x_j) \right \Vert_{L^2(\Omega)}
\end{align}
for some or all $\beta \in [0, \tau]$.
Such a Bernstein inequality Eq.~\eqref{eq:bernstein_not_available} would be helpful to prove our main result \Cref{th:inverse_statement_generalized},
however it is \textit{not} available in the literature in this general form to the best of the author's knowledge:

The statement \cite[Theorem 5.1]{narcowich2006sobolev} formulates and proves a Bernstein inequality,
however it considers $\Omega = \R^d$ instead of compact regions $\Omega \subset \R^d$ and thus cannot be used to show inverse statements on compact regions $\Omega \subset \R^d$.
In \cite{ward2012lp}, these Bernstein inequalities are extended to $L^p(\R^d)$ norms, still working on $\R^d$.
%
The statements in \cite[Section 6]{rieger2008sampling} generalize the $\R^d$ case to star-shaped domains $\Omega \subset \R^d$, 
however only for $\beta = \tau$.
For this, they require restrictive assumptions on the boundary and either restrict the point distribution $X \subset \Omega$  or require $\sum_{j=1}^n \alpha_j \Phi(\cdot - x_j)$ to be the interpolant of a function $f \in H^\tau(\Omega)$
(see \cite[Theorem 6.4.1]{rieger2008sampling} and \cite[Corollary 6.5.3]{rieger2008sampling}).
The result \cite[Lemma 3.2]{cheung2018convergence} states another inverse inequality,
however the assumptions on the involved norms are not appropriate for our purpose of deriving inverse statements.
The paper \cite{hangelbroek2018inverse} focusses on the case of regions $\Omega \subset \R^d$, noting that these are substantially more difficult due to the occuring boundary.
By considering localized Lagrange functions, which are obtained by augmenting the set of interpolation points with centers outside the region, 
inverse statements for such localized Lagrange functions can be derived.
Subsequently, \cite{hangelbroek2018direct} extends these techniques and results to manifolds.

%
The spherical case $\Omega = \mathbb{S}^{d-1} \subset \R^d$ is better researched, 
as no boundary occurs and additional tools like spherical basis functions are available.
In fact, \cite[Theorem 6.2]{narcowich2007direct} is the corresponding statement to our main result \Cref{th:inverse_statement_generalized} for the spherical case.
Also \cite{mhaskar2010bernstein} deals with this particular case of the sphere and generalizes the Bernstein inequalities to the $L^p(\mathbb{S}^{d-1})$ case,
and \cite[Lemma 9.2]{mirzaei2018petrov} generalizes these statements further.

For the more challenging general case of $\Omega \subset \R^d$ being a Lipschitz region,
the only available result so far seems to be \cite[Theorem 6.1]{schaback2002inverse}.
It states an inverse statement based on a decay assumption on $\Vert f - s_{f, X} \Vert_{L^\infty(\Omega)}$,
which can directly be applied pointwise.
However this result is not sharp, as it does not give the correct smoothness: In view of corresponding direct statements, it leaves a ``gap of $d/2$''.


All in all, sharp inverse statements for approximation by kernel interpolation are not yet available, 
and neither are the corresponding standard tools to do so.
The main result \Cref{th:inverse_statement_generalized} proven in \Cref{sec:inverse_statement} provides such sharp inverse statements by leveraging the technical tools developed in \Cref{sec:utility_statements}.

%




%
%
%
%
%
%
%
%
%
%
%
%
%
%
%
%
%
%

%
%
%
%


%
%




%



%
%
%
%
%
%
%
%
%
%
%
%
%
%
%
%
%
%
%
%
%
%
%
%
%
%
%
%
%
%
%
%
%
%
%
%
%
%
%
%
%
%
%
%
%
%
%
%
%
%
%
%
%
%
%
%
%
%
%
%
%
%
%
%
%
%
%
%
%
%
%




%
%
%
%
%
%
%
%
%
%
%
%
%
%
%




\section{Utility statements: From continuous $L^2(\Omega)$ to discrete $L^2(X)$ estimates}
\label{sec:utility_statements}

%

In this section we show novel results on how to obtain estimates on the discrete norm $\Vert g \Vert_{L^2(Y_1)}$ (for a continuous bounded functions $g \in \mathcal{C}(\Omega)$ and some well distributed set $Y_1 \subset \Omega$) based on estimates on the norm $\Vert g \Vert_{L^2(\Omega)}$.
%
The basic connection between these two norms is given by the fact, that the integral occuring within the $\Vert \cdot \Vert_{L^2(\Omega)}$ norm can be approximated by a quadrature given sufficiently many well distributed points.

In \Cref{th:utility_theorem_new} this is formalized, 
and then \Cref{prop:disc_L2_from_cont_L2_new} applies \Cref{th:utility_theorem_new} iteratively to even convert convergence rates on $\Vert f - s_{f, X_n} \Vert_{L^2(\Omega)}$ for sequences of points $(X_n)_{n \in \N} \subset \Omega$ into convergence rates on $\Vert f - s_{f, X_n} \Vert_{L^2(X_{n+1})}$.
%
%

The proofs are quite constructive and leverage the points obtained from the geometric greedy algorithm, which was recalled in \Cref{subsec:geometric_properties} and also resembles ideas from related geometric constructions as in
\cite{duchon1978erreur, rieger2008sampling, rieger2010sampling}.
As these theorems are only used as a technical tool, the proofs are deferred to \Cref{sec:proofs}.

\subsection{Continuous $L^2(\Omega)$ estimates yield discrete $L^2(X)$ estimates}

%
The following \Cref{th:utility_theorem_new} shows that estimates on $\Vert g \Vert_{L^2(\Omega)}$ for $g \in \mathcal{C}(\Omega)$ imply estimates on $\Vert g \Vert_{L^2(Y_1)}$ for some well distributed set of points $Y_1 \subset \Omega$.
For reasons that will become clear later on when applying these results in \Cref{sec:inverse_statement}, we aim for sets $Y_1 \subset \Omega$ which are not ``too close'' to a given reference set $Y_0 \subset \Omega$.
The reason for this is that also the set $Y_0 \cup Y_1 \subset \Omega$ needs to be well distributed (in the sense of having a sufficiently large separation distance $q_{Y_0 \cup Y_1}$). 



\begin{theorem}
\label{th:utility_theorem_new}
Consider a bounded region $\Omega \subset \R^d$ which satisfies an interior cone condition with angle $\alpha \in (0, \pi/2)$ and radius $r > 0$. \\
For any bounded function $g \in \mathcal{C}(\Omega)$, any set finite reference set $Y_0 \subset \Omega$ (of at least two points) and any $q \leq \min \left( \frac{2^{-1/d} c_\Omega}{7 C_\Omega} \cdot q_{Y_0}, \frac{2}{5}r \right)$, 
there exists a finite set of points $Y_1 \subset \Omega$ with
\begin{align*}
\frac{1}{3} q \leq q_{Y_1} \leq h_{Y_1} \leq \frac{22}{3} \cdot q,
\end{align*}
and
\begin{align*}
\frac{1}{6}q \leq q_{Y_0 \cup Y_1} 
\end{align*}
such that it holds
\begin{align}
\label{eq:utility_theorem_new}
%
%
\Vert g \Vert_{L^2(Y_1)} &\leq \sqrt{\tilde{C}_{d,\alpha}} \cdot \Vert g \Vert_{L^2(\Omega)},
\end{align}
with $\tilde{C}_{d,\alpha} := 4 \cdot  \frac{16^d C_d C_\Omega^d}{c_\Omega^{2d} C_{d, \alpha}^2}$ using $C_d, C_{d, \alpha}$ from Eq.~\eqref{eq:constants_ball_cone}.
\end{theorem}

The proof can be found in \Cref{sec:proofs} and an exemplary visualization of the point sets $Y_0$ and $Y_1$ is provided in the bottom right plot of \Cref{fig:vis_proof}.
We have two comments on \Cref{th:utility_theorem_new}:

First, using the definition of the discrete norm $\Vert \cdot \Vert_{L^2(Y_1)}$ from Eq.~\eqref{eq:discrete_L2_norm},
we can rewrite Eq.~\eqref{eq:utility_theorem_new} as
\begin{align*}
\frac{1}{|Y_1|} \sum_{y \in Y_1} |g(y)|^2 &\leq  \tilde{C}_{d,\alpha} \cdot \Vert g \Vert_{L^2(\Omega)}^2,
%
\end{align*}
which is related to the upper bound within Marcinkiewicz-type discretization theorems, 
see e.g.\ \cite[Eq.\ (1.1)]{temlyakov2018marcinkiewicz}.
%
For arbitrary functions $g \in \mathcal{C}(\Omega)$, there does not seem to exist a suitable Marcinkiewicz-type discretization theorem.
However as we are interested in only a single given function $g \in \mathcal{C}(\Omega)$ instead of a whole subspace of functions,
\Cref{th:utility_theorem_new} still allows to derive such an estimate.
In fact, the set $Y_1$ may depend on the function $g$.

Second, as it can be easily seen from the proof of \Cref{th:utility_theorem_new}, 
analogous inequalities using $L^p$ norms with $1 \leq p < \infty$ in Eq.~\eqref{eq:utility_theorem_new} can be derived, 
as the proof was based on a discretization of the occuring integrals.
In order to keep it more simple, we formulated the theorem only in terms of $p=2$, which is the case which will be leveraged in the following.



%
%
%


%

%
%
%
%
%
%
%
%
%
%
%
%
%
%
%
%
%
%
%
%
%
%
%
%
%
%
%
%
%
%
%
%
%
%
%
%
%
%
%
%
%
%
%
%
%
%
%
%
%
%
%
%
%
%
%
%
%
%
%
%
%
%
%
%
%
%
%
%
%
%
%
%
%
%
%
%
%
%
%
%
%
%
%
%
%
%
%
%
%
%
%
%
%
%
%
%
%
%
%
%
%
%
%
%
%
%
%
%
%
%
%
%
%
%
%
%
%
%
%
%
%
%
%
%
%
%
%
%
%
%
%
%
%
%
%
%
%
%
%
%
%
%
%
%
%
%
%
%
%
%
%
%
%
%
%
%
%
%
%
%
%
%
%
%
%
%
%
%
%
%
%
%
%
%
%
%
%
%
%
%
%
%
%
%
%
%
%
%
%
%
%
%
%
%
%
%
%
%
%
%
%
%
%
%


%
%
%

%
%
%
%
%
%
%
%
%
%
%
%
%
%
%
%
%
%
%
%
%
%
%
%
%
%
%
%
%
%
%
%
%
%
%
%
%
%
%
%
%
%
%
%
%
%
%
%
%
%
%
%
%
%
%
%
%
%
%
%
%
%
%
%
%
%
%
%
%
%
%
%
%
%
%
%
%
%
%
%
%
%
%
%
%
%
%
%
%
%
%
%
%
%
%
%
%
%
%
%
%
%
%
%
%
%
%
%
%
%
%
%
%
%
%
%
%
%
%
%
%
%
%
%
%
%

%
%
%
%

\subsection{Continuous $L^2(\Omega)$ convergence rates yield discrete $L^2(X)$ convergence rates}

In the following,
we consider the interpolation of continuous functions via kernel interpolation and their corresponding convergence rates.
The corresponding convergence rates are usually given in the ``continuous'' $\Vert \cdot \Vert_{L^2(\Omega)}$ norm for well distributed points, see \Cref{th:error_estimate}.
%

We aim at using \Cref{th:utility_theorem_new} iteratively for $g := f - s_{f, X_n}$ in order to obtain convergence rates in a discrete $\Vert f - s_{f, X_n} \Vert_{L^2(X_{n+1})}$ norm 
for some sequence $(X_n)_{n \in \N} \subset \Omega$ based on convergence rates in the continuous $\Vert f - s_{f, X_n} \Vert_{L^2(\Omega)}$ norm.
%

\begin{prop}
\label{prop:disc_L2_from_cont_L2_new}
%
Consider a bounded region $\Omega \subset \R^d$ which satisfies an interior cone condition with angle $\alpha \in (0, \pi/2)$ and radius $r>0$. \\
Consider a function $f \in \mathcal{C}(\Omega)$ and a continuous kernel $k$ such that for some $c_f, \mu, h_0 > 0$ it holds
\begin{align*}
\Vert f - s_{f, X} \Vert_{L^2(\Omega)} \leq c_f h_X^{\mu}
\end{align*}
for all $X \subset \Omega$ with $h_X \leq h_0$ and $\rho_X \leq 44$ (quasi-uniform points). \\
%
%
%
%
Then there exist constants $a \in (0, \frac{3}{154})$ and $c_0 > 0$ and a sequence of nested sets of points $X_0 \subset X_1 \subset ... \subset X_m \subset .. \subset \Omega$ with 
\begin{align*}
\frac{1}{6} c_0 a^{m+1} \leq q_{X_{m+1}} \leq h_{X_{m+1}} \leq \frac{22}{3} c_0 a^{m+1}
\end{align*}
%
%
%
%
such that
\begin{align*}
\Vert f - s_{f, X_m} \Vert_{L^2(X_{m+1})} \leq c_f \sqrt{\tilde{C}_{d,\alpha}} h_{X_m}^{\mu}
\end{align*}
for all $m > 0$, using $\tilde{C}_{d,\alpha}$ defined in \Cref{th:utility_theorem_new}.
\end{prop}

The proof can be found in \Cref{sec:proofs}.
%
%
%




%
%
%
%
%
%
%
%
%
%







%
%
%


%
%
%
%
%
%
%
%
%
%
%
%
%
%
%
%
%
%
%
%
%
%
%
%
%
%
%
%
%
%
%
%
%
%
%
%
%
%
%
%
%
%
%
%
%
%
%
%
%
%
%
%
%
%
%
%
%
%
%
%
%
%
%
%
%
%
%






\section{Inverse statement: Continuous $L^2(\Omega)$ convergence rates yield smoothness} \label{sec:inverse_statement}


In the following we state and prove our main result \Cref{th:inverse_statement_generalized},
which is an inverse statement: 
Given a convergence rate on $\Vert f - s_{f, X} \Vert_{L^2(\Omega)}$ in the fill distance $h_X$ for the approximation of a function $f \in \mathcal{C}(\Omega)$,
we derive a corresponding Sobolev smoothness of the function $f$.
%
%

%
%
%
%

%
%
%

\mainresult*

%
%
%
%
%
%
%
%

%
The idea of the proof is to consider a sequence of interpolants on a sequence of nested and increasingly denser point sets.
Then one shows that this sequence is a Cauchy sequence in a Sobolev space of some smoothness,
and by completeness of the Sobolev spaces one then obtains a limiting element.



\begin{proof}[Proof of \Cref{th:inverse_statement_generalized}]
We start by using \Cref{prop:disc_L2_from_cont_L2_new} for $\mu := \beta$:
Like this we obtain constants $a \in (0, \frac{3}{154})$ and $c_0 > 0$ and a sequence of nested sets of points $(X_m)_{m \geq 0}$ such that
\begin{align}
\begin{aligned}
\label{eq:proof:properties_sequence}
\frac{1}{6} c_0 a^{m+1} \leq& ~ q_{X_{m+1}} \leq h_{X_{m+1}} \leq \frac{22}{3} c_0 a^{m+1}, \\
%
\Vert f - s_{f, X_m} \Vert_{L^2(X_{m+1})} \leq& ~ c_f \sqrt{\tilde{C}_{d,\alpha}} h_{X_m}^{\beta}
\end{aligned}
\end{align}
Observe that these sets of points $X_m$ have a uniformity constant $\rho_{X_m}$ bounded as $\rho_{X_m} \equiv \frac{h_{X_m, \Omega}}{q_{X_m}} \leq 44$.

%
We consider $\beta' < \beta$ for $\beta \in (0, \tau]$ respective $\beta'=\tau$ for $\beta > \tau$ and 
continue by applying \Cref{prop:hoelder_interpol} to $g := s_{f, X_{n+1}} - s_{f, X_n} \in \ns \asymp H^\tau(\Omega)$ using $\theta = \beta'/\tau$ for $\beta \in (0, \tau]$ or $\theta = 1$ for $\beta > \tau$:
%
\begin{align}
\begin{aligned}
\label{eq:application_hoelder}
&\Vert s_{f, X_{m+1}} - s_{f, X_{m}} \Vert_{H^{\beta'}(\Omega)}^2 \\
\leq& ~ C \cdot \Vert s_{f, X_{m+1}} - s_{f, X_{m}} \Vert_{L^2(\Omega)}^{2 - 2\theta} \cdot \Vert s_{f, X_{m+1}} - s_{f, X_{m}} \Vert_{H^\tau(\Omega)}^{2\theta} \\
\leq& ~ C' \cdot\Vert s_{f, X_{m+1}} - s_{f, X_{m}} \Vert_{L^2(\Omega)}^{2 - 2\theta} \cdot \Vert s_{f, X_{m+1}} - s_{f, X_{m}} \Vert_{\ns}^{2\theta} \\
\leq& ~ C' \cdot \Vert s_{f, X_{m+1}} - s_{f, X_{m}} \Vert_{L^2(\Omega)}^{2 - 2\theta} \\ 
&\qquad \cdot \Vert A_{X_{m+1}}^{-1} \Vert_{2,2}^{\theta} \cdot |X_{m+1}|^{\theta} \cdot \Vert s_{f, X_{m+1}} - s_{f, X_{m}} \Vert_{L^2(X_{m+1})}^{2\theta},
\end{aligned}
\end{align}
where we estimated the $\Vert s_{f, X_{m+1}} - s_{f, X_{m}} \Vert_{\ns}$ quantity by using \Cref{prop:upper_bound_ns_norm_via_L2}.
%
%
%
For these four factors we have the following estimates:
%
\begin{itemize}
\item Adding $0 = f - f$, leveraging the assumed bound on $\Vert f - s_{f, X} \Vert_{L^2(\Omega)}$ and then estimating the fill distances as in Eq.\ \eqref{eq:proof:properties_sequence} gives
\begin{align*}
&\Vert s_{f, X_{m+1}} - s_{f, X_{m}} \Vert_{L^2(\Omega)} \\
\leq \, &\Vert f - s_{f, X_{m+1}} \Vert_{L^2(\Omega)} + \Vert f - s_{f, X_{m}} \Vert_{L^2(\Omega)} \\
\leq \, &2c_f h_{X_m}^{\beta} = 2c_f \left( \frac{22c_0}{3} \right)^\mu a^{\beta m}.
\end{align*}
\item For the norm of the inverse matrix we use \Cref{th:stability_statement} and subsequently the bound from Eq.\ \eqref{eq:proof:properties_sequence} on the separation distance to obtain
\begin{align*}
\Vert A_{X_{m+1}}^{-1} \Vert_{2,2} \leq& ~ c_0^{-1} q_{X_{m+1}}^{d-2\tau} \leq c_0^{-1} \left(\frac{1}{6} c_0 a^{m+1} \right)^{d-2\tau} \\
=& ~ \frac{6^{2\tau-d}}{c_0^{2\tau - d + 1} a^{2\tau-d}} a^{-m (2\tau - d)}.
%
\end{align*}
\item The number of points $|X_{m+1}|$ can be upper estimated via Eq.\ \eqref{eq:geometric_bounds} as 
\begin{align*}
|X_{m+1}| \leq C_\Omega^d q_{X_{m+1}}^{-d} \leq \frac{6^d C_\Omega^d}{c_0^d} a^{-d(m+1)} = \frac{6^d C_\Omega^d}{c_0^d a^d} a^{-dm}.
\end{align*}
\item Using $s_{f, X_{m+1}}|_{X_{m+1}} = f|_{X_{m+1}}$, 
the bound on $\Vert f - s_{f, X_m} \Vert_{L^2(X_{m+1})}$ from Eq.\ \eqref{eq:proof:properties_sequence}, and subsequently estimating the fill distances via Eq.\ \eqref{eq:proof:properties_sequence} gives
\begin{align*}
&\Vert s_{f, X_{m+1}} - s_{f, X_{m}} \Vert_{L^2(X_{m+1})} = \Vert f - s_{f, X_{m}} \Vert_{L^2(X_{m+1})} \\
%
\leq \, &c_f \cdot \sqrt{\tilde{C}_{d, \alpha}} \cdot h_{X_m}^\beta \leq \, c_f \cdot \sqrt{\tilde{C}_{d, \alpha}} \cdot \left( \frac{22}{3} c_0 a^m \right)^\beta \\
\leq \, &c_f \cdot \sqrt{\tilde{C}_{d, \alpha}} \cdot \left( \frac{22c_0}{3} \right)^\beta a^{\beta m}.
\end{align*}
\end{itemize}

\noindent Putting everything together we obtain %
\begin{align*}
&\Vert s_{f, X_{m+1}} - s_{f, X_{m}} \Vert_{H^{\beta'}(\Omega)}^2 \\
\leq &\, C' \cdot \Vert s_{f, X_{m+1}} - s_{f, X_{m}} \Vert_{L^2(\Omega)}^{2 - 2\theta} \cdot \Vert A_{X_{m+1}}^{-1} \Vert_{2,2}^{\theta} \cdot |X_{m+1}|^{\theta} \cdot  \Vert s_{f, X_{m+1}} - s_{f, X_{m}} \Vert_{L^2(X_{m+1})}^{2\theta} \\
\leq &\,C'' \cdot a^{2\beta m (1-\theta)} \cdot a^{-m (2\tau - d) \theta} \cdot a^{-md \theta} \cdot a^{2\beta m \theta} \\
= &\,C'' \cdot a^{m(\theta \cdot [-2\beta - 2\tau + d - d + 2\beta] + 2\beta)} \\
= &\,C'' \cdot a^{2m(\beta-\beta')},
\end{align*}
with a generic constant $C'' > 0$ given by the collection of the previously occuring constants. %
%
%
%
As $\beta' < \beta$ we obtain $\beta - \beta' =: \delta > 0$.
Then we can show that $(s_{f, X_m})_{m \geq 1}$ is a Cauchy series, i.e. for $n > m \geq m_0$:
\begin{align}
\label{eq:cauchy_sequence_calculation}
\Vert s_{f, X_{n}} - s_{f, X_m} \Vert_{H^{\beta'}(\Omega)} &= \left \Vert \sum_{\ell=m}^{n-1} s_{f, X_{\ell+1}} - s_{f, X_{\ell}} \right \Vert_{H^{\beta'}(\Omega)} \notag \\
&\leq \sum_{\ell=m}^{n-1} \Vert s_{f, X_{\ell+1}} - s_{f, X_{\ell}} \Vert_{H^{\beta'}(\Omega)} \notag \\
&\leq \sqrt{C''} \cdot \sum_{\ell=m}^{\infty} a^{\ell\delta} 
\leq \sqrt{C''} \cdot \frac{a^{\delta m_0}}{1 - a^{\delta}} \stackrel{m_0 \rightarrow \infty}{\longrightarrow} 0.
\end{align}
As the Sobolev space $H^{\beta'}(\Omega)$ is a complete space and $(s_{f, X_m})_{m \geq m_0}$ is a Cauchy sequence in $H^{\beta'}(\Omega)$, 
there exists a unique limit element $\tilde{f} \in H^{\beta'}(\Omega)$. \\
It remains to show that $\tilde{f}$ and $f$ coincide on $\Omega$. 
For this we consider
\begin{align*}
\Vert f - \tilde{f} \Vert_{L^2(\Omega)} %
&\leq \Vert f - s_{f, X_n} \Vert_{L^2(\Omega)} + \Vert \tilde{f} - s_{f, X_n} \Vert_{L^2(\Omega)} \\
&\leq c_f h_{X_n}^{\mu} + C_{H^{\beta'}(\Omega) \hookrightarrow L^2(\Omega)} \Vert \tilde{f} - s_{f, X_n} \Vert_{H^{\beta'}(\Omega)} \stackrel{n \rightarrow \infty}{\longrightarrow} 0.
\end{align*}
Thus we have $f = \tilde{f}$ in $L^2(\Omega)$ an $f$ can be seen as a particular representative of $\tilde{f}$, 
i.e.\ we have proven that $f \in H^{\beta'}(\Omega)$ for any $\beta' < \beta$. 
\end{proof}

\noindent In the following we collect some remarks on the theorem and the proof:
\begin{itemize}
%
%
%
\item First we remark that the inequalities in Eq.~\eqref{eq:application_hoelder} essentially bounds a strong norm (the $\Vert \cdot \Vert_{H^{\beta'}(\Omega)}$ norm) via weaker norms 
(the $\Vert \cdot \Vert_{L^2(\Omega)}$ norm and the $\Vert \cdot \Vert_{L^2(X_m \cup X_{m+1})}$ norm, which can be again bounded via the $\Vert \cdot \Vert_{L^2(\Omega)}$ norm).
Thus this can be seen as a Bernstein inequality on general Lipschitz regions $\Omega \subset \R^d$.
However it is not a Bernstein inequality in its full generality, 
as it holds only with respect to the derived sequence $(X_m)_{m \geq m_0} \subset \Omega$ of  sets of points.
%
%
%
%
%
%
%
%

%
%
%
%
%
%
%
%
%
%
%
%
%
%
%
%
%
%
\item Second we remark that the statement of \Cref{th:inverse_statement_generalized} can be refined by additionally taking into account logarithmic decay rates, 
i.e.\ within \Cref{th:inverse_statement_generalized} we could assume a decay as
\begin{align*}
\Vert f - s_{f, X} \Vert_{L^2(\Omega)} \leq c_f h_X^{\beta} \cdot |\log(h_X)|^{-\nu}
\end{align*}
for all $X \subset \Omega$ with $h_X \leq h_0$ and $\rho_X \leq 44$ (quasi-uniformly distributed) and some $\nu \in \R$. \\
%
Due to $h_{X_m} \stackrel{\bigcdot}{\leq} a^m$ in the proof of \Cref{th:inverse_statement_generalized} (where $\stackrel{\bigcdot}{\leq}$ denotes inequality up to a constant), 
the decay $h_{X_m}^{\mu}$ turned into quantities like $a^{\mu m}$.
And thus the parts $|\log(h_{X_m})|^{-\nu}$ will turn into $m^{-\nu} |\log(a)|^{-\nu} \stackrel{\bigcdot}{=} m^{-\nu}$.
Therefore, when showing the Cauchy sequence property in Eq.~\eqref{eq:cauchy_sequence_calculation}, 
we would not obtain the pure geometric series $\sum_{\ell=m}^{\infty} a^{\ell\delta}$ anymore.
Instead, in particular for $\beta = \beta'$, we would obtain a harmonic sum like $\sum_{\ell=m}^{\infty} \ell^{-\tilde{\delta} \nu}$ (for some constant $\tilde{\delta} > 0$) which still converges for $\nu > 0$ large enough.
We waived to include this more detailed but more complicated analysis here in order to keep a clear outline of the exposition.
%
%
%

%
%
\end{itemize}

%




%
%
%
%

Finally we point out that in view of \Cref{th:error_estimate} the inverse statement in \Cref{th:inverse_statement_generalized} is sharp for $\Omega \subset \R^d$ with Lipschitz boundary (up to an arbitrary small $\varepsilon > 0$):

For $f \in H^\beta(\Omega)$ with $d/2 < \beta \leq \tau$, 
the direct statement \Cref{th:error_estimate} provides an $\Vert \cdot \Vert_{L^2(\Omega)}$ error decay for asymptotically uniformly distributed points $X \subset \Omega$ according to $h_{X}^\beta$.
On the other hand given an approximation rate as $h_X^\beta$, the inverse statement \Cref{th:inverse_statement_generalized} already provides $f \in H^{\beta'}(\Omega)$ for any $\beta' < \beta$.
%
This establishes a one-to-one correspondence between smoothness and approximation rate for kernel based interpolation using finitely smooth kernels.
%
%





%
%
%
%
%
%
%
%
%
%
%







%
%
%
%
%
%
%
%



%
%
%
%
%
%
%
%
%
%
%
%
%
%
%
%
%
%
%
%
%
%
%
%
%
%
%
%
%
%
%
%
%
%
%
%
%
%
%
%
%
%
%
%
%
%
%
%
%
%
%
%
%
%
%
%
%
%
%
%
%
%
%
%
%
%
%
%
%
%
%
%
%
%
%
%
%
%
%
%
%


%
\section{Discussion} \label{sec:discussion_remarks}

%
%
%
%
%
%
%
%
%
%
%
%
%
%
%
%
%
%
%
%
%

We start by commenting on the relation between the main result \Cref{th:inverse_statement_generalized} and the related literature discussed in \Cref{subsec:related_work}:
The standard approach for showing inverse statements is with help of Bernstein inequalities.
As these are not available in the literature for the general case of compact Lipschitz domains because of challenges due to the boundary,
we employed the utility statements of \Cref{sec:utility_statements} instead.
These statements allowed to turn the $L^2(\Omega)$ convergence rates into point based estimates, 
which enabled the construction of a sequence of interpolants that turned out to be a Cauchy sequence.
As elaborated in the previous section, the resulting estimate can also be understood as a Bernstein inequality.
While it is evident that the utility statements of \Cref{sec:utility_statements} are generalizable to other $L^p(\Omega)$ norms, $1 \leq p < \infty$,
the interpolation inequality within the key step Eq.~\eqref{eq:application_hoelder} does not transfer in a straightforward way.
Thus it remains a future task to transfer the $L^2(\Omega)$ error analysis of the inverse statement from \Cref{th:inverse_statement_generalized} into $L^\infty(\Omega)$ or other $L^p(\Omega)$ error statements that are sharp as well. 

For the case of $\Omega$ being a compact Lipschitz region, 
the only inverse statement so far was given in \cite{schaback2002inverse} and was formulated in terms of $L^\infty(\Omega)$ convergence rates.
In contrast to \Cref{th:inverse_statement_generalized}, 
it covered only the edge case of smoothness $\tau$ (instead of intermediate smoothnesses $\beta \in (0, \tau]$) and did not yield a sharp inverse statement, thus leaving a ``gap of $d/2$''.
Such a gap is no longer existent in the present approach which considers $L^2(\Omega)$ convergence rates.
In fact, our inverse statement \Cref{th:inverse_statement_generalized} is exactly \cite[Theorem 6.2]{narcowich2007direct},
however for the case of a compact Lipschitz region $\Omega$ instead of the sphere $\mathbb{S}^{d-1}$. 

\Cref{th:inverse_statement_generalized} makes use of well distributed points due to considering the fill distance $h_{X}$.
This is a frequent assumption in kernel based interpolation, 
and also useful in applications since greedy algorithms as the $P$-greedy algorithm are known to provide well distributed points \cite[Theorem 19 \& 20]{wenzel2021novel}.
It remains open whether and how inverse statements can be generalized to adaptively distributed points, 
such as greedily selected points \cite{wenzel2023analysis, santin2024optimality}.

The key ingredients for the proof of \Cref{th:inverse_statement_generalized} have been the technical statement of \Cref{prop:disc_L2_from_cont_L2_new} as well as standard RKHS analysis tools including stability estimates as in \Cref{th:stability_statement}.
These tools do not seem to rely too much on the considered domain $\Omega \subset \R^d$, such that an extension to manifolds seems possible: 
In the case of manifolds, the asymptotic of fill and separation distance for well distributed points still coincide 
-- in view of Eq.~\eqref{eq:geometric_bounds} they will however depend on the effective dimension of the manifold instead of the dimension of the ambient space.
%
Also weakening the assumption of a Lipschitz boundary or cone condition might be possible, 
as \cite{wenzel2024stability} showed that the $L^\infty(\Omega)$ error estimates do not depend significantly on the shape of the region $\Omega \subset \R^d$ 
-- though it remains open whether this holds as well for the $L^2(\Omega)$ error.

%
%
%
%
%
%
%
%
%
%
%
%
%
%

%
We continue with a remark on the sharpness of the inverse statements, which allowed to conclude the one-to-one correspondence between smoothness and approximation rates as elaborated after \Cref{th:inverse_statement_generalized}:
We considered continuous kernels satisfying a norm-equivalence $\ns \asymp H^\tau(\Omega$ for some $\tau > d/2$, 
which allowed to make use of stability statements which were originally derived for RBF kernels, see \Cref{th:stability_statement}.
Such RBF kernels are usually characterized with help of the decay of their Fourier transform, i.e.\ Eq.~\eqref{eq:fourier_decay}:
The upper bound on the Fourier decay usually yields the error estimates (via the fill distance), i.e.\ the direct statements,
while the lower bound on the Fourier decay yields stability estimates (via the separation distance) which are required for the inverse statements \cite{schaback2002inverse, schaback1995error}.
In case of a gap between the lower and upper bound, i.e.
\begin{align*}
c_\Phi (1+\Vert \omega \Vert_2^2)^{-\tau_1} \leq \hat{\Phi}(\omega) \leq C_\Phi (1 + \Vert \omega \Vert_2^2)^{-\tau_2}
\end{align*}
for $d/2 < \tau_1 < \tau_2$,
this could result in a gap between direct and inverse statements.
Such different exponents, i.e.\ $\tau_1 < \tau_2$, can happen if the Fourier transform $\hat{\Phi}$ has a staircase behavior.
We remark that such a kernel is probably rather exotic and of little practical use.

%
%

In this paper we considered kernels of finite smoothness, i.e.\ $\ns \asymp H^\tau(\Omega)$.
However there are also kernels of infinite smoothness, such as the Gaussian kernel $k(x, z) = \exp(- \Vert x - z \Vert_2^2)$.
A one-to-one correspondence as discussed after \Cref{th:inverse_statement_generalized}) might not be derivable here:
For example for the Gaussian kernel, the known error estimates give rates as $\exp(c \log(h_{X}) / h_{X})$ \cite[Theorem 11.22]{wendland2005scattered}.
However the corresponding stability estimates look like $\exp(-40.71d^2/q_X^2) q_X^{-d}$ \cite[Corollary 12.4]{wendland2005scattered}, 
i.e.\ there is a large gap even for well distributed points, i.e.\ $q_X \asymp h_{X}$.
A detailed investigation on this possible \textit{analytic gap} between direct and inverse statements seems to be an interesting future research direction.
Another possible direction for future work is to improve the case $\beta > \tau$ within \Cref{th:inverse_statement_generalized}, 
as one would naturally expect to conclude increased smoothness here as well.
Within the presented analysis this was not possible, as \Cref{prop:hoelder_interpol} is no longer applicable for $\beta > \tau$.



Note that \Cref{th:inverse_statement_generalized} and thus the one-to-one correspondence also hold for large values of the dimension $d$,
as only the asymptotic behavior of the error for $h_{X} \rightarrow 0$ is of importance.
Though the dimensionality will affect the constants involved.
However when dealing with only finitely many scattered data points in high dimensions, 
the fill distance is usually larger than one.
We point to \cite{lin2024kernel} for more details on kernel interpolation of scattered data in high dimensions.



%
%
%
%
%
%





Note that \Cref{th:inverse_statement_generalized} also works for large values of the dimension $d$,
as long as Eq.~\eqref{eq:inverse_statement_assumption} is satisfied.
The dimensionality only affects the constants involved, however not the asymptotic behavior of the error for $h_{X, \Omega} \rightarrow 0$,
which is only linked to the Sobolev smoothness .
We point to \cite{lin2024kernel} for details on kernel interpolation in high-dimensional domains.
%
%
%
%
%
%
%
%


%
%

%
%
%
%

Finally, we conclude with a numerical example related to the one-to-one correspondence,
i.e.\ \Cref{th:error_estimate} and \Cref{th:inverse_statement_generalized}:

\begin{example} %
We consider the reconstruction of continuous functions on $\Omega := [0, 1]$ using the Matérn kernel $k(x, z) = \exp(-\Vert x - z \Vert_2)$,
 which is an RBF kernel with $\Phi(r) = \exp(-r)$ and corresponding Fourier transform $\hat{\Phi}(\omega) = \sqrt{\frac{2}{\pi}} \cdot \frac{1}{\omega^2 + 1}$.
Thus, the RKHS $\ns$ is norm equivalent to the Sobolev space $H^1(\Omega)$, i.e.\ $\tau = 1$.
We consider the functions
\begin{align*}
f_\sigma: \Omega \rightarrow \R, x \mapsto x^\sigma \cdot (1-x)^\sigma
\end{align*}
for $\sigma > 0$ and it can be shown that $f_\sigma \in H^{\sigma' + 1/2}(\Omega)$ ~ $\forall \sigma' \in (0, \sigma)$ for any $\sigma \in (0, 2) \setminus \{ 1 \}$.
%
For the numerical experiment, we make use of the values %
$\sigma \in \{ 0.21, 0.36, 0.51, 0.81, 1.11, 1.51 \}$ 
and consider the $\Vert f_\sigma - s_{f_\sigma, X} \Vert_{L^2(\Omega)}$ error using up to $2^{13} + 1 = 8193$ equidistant interpolation points $X \subset \Omega$,
where each both boundary points are included.
The $\Vert f_\sigma - s_{f_\sigma, X} \Vert_{L^2(\Omega)}$ error was evaluated by discretizing $\Omega$ using $2^{20} + 1$ equally spaced points.

The results are visualized in \Cref{fig:matern_L2_error}:
For every value of $\sigma$, the decay of the error $\Vert f_\sigma - s_{f_\sigma, X} \Vert_{L^2(\Omega)}$ in the number of interpolation points $n$ is depicted.
The legend additionally lists the numerically computed decay rates,
and the dashed black lines indicate the theoretically expected decay.
For $\sigma \in \{ 0.21, 0.36 \}$ we have $f_\sigma \in H^{\sigma' + 1/2} \supset H^\tau(\Omega)$ ~ $\forall \sigma' \in (0, \sigma)$, 
i.e.\ exactly the situation described in \Cref{th:error_estimate} and \Cref{th:inverse_statement_generalized}: 
We obtain a decay as $n^{-0.69}$ and $n^{-0.85}$, which almost exactly match the expected decay of $n^{-0.71}$ respectively $n^{-0.86}$.
For $\sigma \in \{ 0.51, 0.81, 1.11, 1.51 \}$ we have $\sigma + 1/2 > 1$ and thus $f_\sigma \in H^1(\Omega)$, 
i.e.\ additional smoothness and apparently an increased convergence rate. %
The dashed black lines indicate likewise a decay as $n^{-1.01}$, $n^{-1.31}$, $n^{-1.61}$ and $n^{-2.01}$, i.e.\ scaling according to the smoothness of $f_\sigma$.
This seems to be an instance of superconvergence, see e.g.\ \cite{schaback2018superconvergence}.

%


%
%


\begin{figure}[hbt!]
\centering
\setlength\fwidth{.65\textwidth}
%
%
%
%
\begin{tikzpicture}

\definecolor{crimson2143940}{RGB}{214,39,40}
\definecolor{darkgray176}{RGB}{176,176,176}
\definecolor{darkorange25512714}{RGB}{255,127,14}
\definecolor{forestgreen4416044}{RGB}{44,160,44}
\definecolor{lightgray204}{RGB}{204,204,204}
\definecolor{mediumpurple148103189}{RGB}{148,103,189}
\definecolor{sienna1408675}{RGB}{140,86,75}
\definecolor{steelblue31119180}{RGB}{31,119,180}

\begin{axis}[
legend cell align={left},
legend style={
  fill opacity=0.8,
  draw opacity=1,
  text opacity=1,
  at={(0.03,0.03)},
  anchor=south west,
  draw=lightgray204
},
log basis x={10},
log basis y={10},
tick align=outside,
tick pos=left,
x grid style={darkgray176},
xmin=3.45339580257623, xmax=11862.2371549303,
xmode=log,
xtick style={color=black},
xtick={0.1,1,10,100,1000,10000,100000,1000000},
xticklabels={
  \(\displaystyle {10^{-1}}\),
  \(\displaystyle {10^{0}}\),
  \(\displaystyle {10^{1}}\),
  \(\displaystyle {10^{2}}\),
  \(\displaystyle {10^{3}}\),
  \(\displaystyle {10^{4}}\),
  \(\displaystyle {10^{5}}\),
  \(\displaystyle {10^{6}}\)
},
y grid style={darkgray176},
ymin=5.78622928893346e-09, ymax=0.280003478947575,
ymode=log,
ytick style={color=black},
ytick={1e-10,1e-09,1e-08,1e-07,1e-06,1e-05,0.0001,0.001,0.01,0.1,1,10},
yticklabels={
  \(\displaystyle {10^{-10}}\),
  \(\displaystyle {10^{-9}}\),
  \(\displaystyle {10^{-8}}\),
  \(\displaystyle {10^{-7}}\),
  \(\displaystyle {10^{-6}}\),
  \(\displaystyle {10^{-5}}\),
  \(\displaystyle {10^{-4}}\),
  \(\displaystyle {10^{-3}}\),
  \(\displaystyle {10^{-2}}\),
  \(\displaystyle {10^{-1}}\),
  \(\displaystyle {10^{0}}\),
  \(\displaystyle {10^{1}}\)
}
]
\addplot [semithick, steelblue31119180, mark=x, mark size=5, mark options={solid}, only marks]
table {%
5 0.0466967065630209
9 0.0478759498328178
17 0.0366928179821262
33 0.024860053254987
65 0.0159639205368953
129 0.00999936475371457
257 0.00619026530859231
513 0.00381240517195436
1025 0.00234479220433814
2049 0.00144477068896539
4097 0.000895486726452085
8193 0.000562119428182609
};
\addlegendentry{$\sigma$ = 0.21, $n^{-0.69}$}
\addplot [semithick, darkorange25512714, mark=x, mark size=5, mark options={solid}, only marks]
table {%
5 0.0697358374924028
9 0.0569537027628708
17 0.0371838280793163
33 0.0221739431899068
65 0.0126942524694815
129 0.00712789011388194
257 0.00396513965147717
513 0.00219609361367931
1025 0.00121422393322513
2049 0.00067141472289871
4097 0.000372011851195291
8193 0.000207155566541475
};
\addlegendentry{$\sigma$ = 0.36, $n^{-0.85}$}
\addplot [semithick, forestgreen4416044, mark=x, mark size=5, mark options={solid}, only marks]
table {%
5 0.075541297598699
9 0.0532783496987298
17 0.0308363709378836
33 0.016457413027936
65 0.0084639073698255
129 0.00427644264895892
257 0.00214216475070379
513 0.00106863856079846
1025 0.00053215515675901
2049 0.000264915239379014
4097 0.000132009547594661
8193 6.5968050457666e-05
};
\addlegendentry{$\sigma$ = 0.51, $n^{-1.00}$}
\addplot [semithick, crimson2143940, mark=x, mark size=5, mark options={solid}, only marks]
table {%
5 0.0648008188362111
9 0.0366075789899073
17 0.0171349829740573
33 0.00741653916940206
65 0.00309616456120359
129 0.00127025922857185
257 0.000516734408817335
513 0.000209337643578268
1025 8.46457933686391e-05
2049 3.42056082368218e-05
4097 1.38275281985738e-05
8193 5.59816640782513e-06
};
\addlegendentry{$\sigma$ = 0.81, $n^{-1.31}$}
\addplot [semithick, mediumpurple148103189, mark=x, mark size=5, mark options={solid}, only marks]
table {%
5 0.047421540343049
9 0.0221105283577681
17 0.0084724650073319
33 0.00299030003759234
65 0.00101594650587007
129 0.00033887945146583
257 0.00011202360055803
513 3.68690778296606e-05
1025 1.2109266953291e-05
2049 3.97406564258091e-06
4097 1.30437114496383e-06
8193 4.28573685975758e-07
};
\addlegendentry{$\sigma$ = 1.11, $n^{-1.61}$}
\addplot [semithick, sienna1408675, mark=x, mark size=5, mark options={solid}, only marks]
table {%
5 0.0283674790211751
9 0.0104127842643092
17 0.00308245144902635
33 0.00083354641630121
65 0.000216154221300006
129 5.49354201340526e-05
257 1.38254619897806e-05
513 3.46284208466582e-06
1025 8.6537885824424e-07
2049 2.16063513561059e-07
4097 5.3941825847535e-08
8193 1.34765354532366e-08
};
\addlegendentry{$\sigma$ = 1.51, $n^{-2.00}$}
\addplot [semithick, black, dashed, forget plot]
table {%
5 0.111634733840691
9 0.073545458981321
17 0.046821929873124
33 0.0292363496419669
65 0.0180675623091796
129 0.0111057872618015
257 0.0068079369782053
513 0.00416758738826269
1025 0.00254949462044316
2049 0.00155909742364156
4097 0.000953272818856703
8193 0.000582805372242441
};
\addplot [semithick, black, dashed, forget plot]
table {%
5 0.125272516182551
9 0.0755651321550609
17 0.0437304485456777
33 0.0247199743646695
65 0.0137995370750179
129 0.00765354876153395
257 0.00423084892923338
513 0.0023349042213256
1025 0.00128750071857208
2049 0.000709649531044706
4097 0.000391065294134118
8193 0.000215481040096465
};
\addplot [semithick, black, dashed, forget plot]
table {%
5 0.118084133203615
9 0.0652178257152781
17 0.0343081926259354
33 0.017557075273129
65 0.00885337296697497
129 0.00443052960605404
257 0.00220860878768589
513 0.00109883557571443
1025 0.000546160322562205
2049 0.00027132753693993
4097 0.000134759883478144
8193 6.69227634405898e-05
};
\addplot [semithick, black, dashed, forget plot]
table {%
5 0.0971492117820942
9 0.0449813307451723
17 0.0195524822357925
33 0.00820043188796158
65 0.00337422073423558
129 0.00137473017837649
257 0.000557285310245906
513 0.000225338935663181
1025 9.10000693577902e-05
2049 3.67256948581593e-05
4097 1.48169758903464e-05
8193 5.97695128479253e-06
};
\addplot [semithick, black, dashed, forget plot]
table {%
5 0.0636908754939917
9 0.0247222779882753
17 0.00887965642656872
33 0.00305218941941043
65 0.00102477476292188
129 0.000339914344334921
257 0.000112053802093367
513 3.68239157183055e-05
1025 1.20824081527805e-05
2049 3.96128758364237e-06
4097 1.29822128026726e-06
8193 4.25378728130128e-07
};
\addplot [semithick, black, dashed, forget plot]
table {%
5 0.0373933088478114
9 0.011473506375836
17 0.00319537088182732
33 0.000842384924720835
65 0.00021565908509298
129 5.43798853456246e-05
257 1.36068634779611e-05
513 3.39146782627912e-06
1025 8.43662286884708e-07
2049 2.09664193991324e-07
4097 5.20795253861513e-08
8193 1.2933118773864e-08
};
\end{axis}

\end{tikzpicture}
	
%
\caption{Visualization of the $L^2(\Omega)$ error decay in the ($y$-axis) for interpolation with the Matérn kernel $k(x, z) = \exp(-\Vert x - z \Vert)$ on $\Omega = [0,1]$ in the number of equidistant interpolation points ($x$-axis) for different $\sigma$ parameters of the functions $x^\sigma \cdot (1-x)^\sigma$.
This class of functions scales linearly in the Sobolev spaces $H^\tau(\Omega)$: 
It holds $x^\sigma \cdot (1-x)^\sigma \in H^{\sigma + 1/2}(\Omega)$ ~ $\forall \sigma' \in (0, \sigma)$.
The dashed black lines indicate the convergence rates $n^{-0.71}, n^{-0.86}, n^{-1.01}$, $n^{-1.31}$, $n^{-1.61}$ and $n^{-2.01}$,
while the legend lists numerically computed convergence rates.}
\label{fig:matern_L2_error}
\end{figure}

\end{example}

\section{Conclusion} \label{sec:conclusion}

The current paper proved sharp inverse statements for kernel interpolation using finitely smooth kernels for which the RKHS $\ns$ is norm-equivalent to a Sobolev space $H^\tau(\Omega)$, $\tau > d/2$.
This result includes in particular popular RBF kernels like the class of Matérn or Wendland kernels.
As a techincal tool to do so, discrete $L^2(X)$ error estimates based of continuous $L^2(\Omega)$ error estimates were derived.

%
%
As a result of the sharpness of the inverse statement, 
a one-to-one correspondence between smoothness (of the function to approximate) and the corresponding approximation rate was concluded. %
A numerical experiment illustrated the result.
The derived techniques and results are expected to yield further insights, e.g.\ for kernel quadrature, PDEs approximation with kernels or in statistical learning theory. \\
For future research, we aim at several directions:
First, investigations of other kernels such as analytic kernels as the Gaussian and other $L^p(\Omega)$, $1 \leq p \leq \infty$ norms seems meaningful.
Second, the results should be extended to the superconvergence case, i.e.\ the case where the considered function is smoother than the RKHS. \\


\textbf{Acknowledgements:}
The work was mainly conducted while the author was affiliated with the University of Stuttgart.
The author thanks Prof.\ Bernard Haasdonk and Prof.\ Gabriele Santin for support, discussions and feedback
and acknowledges funding of the project by the Deutsche Forschungsgemeinschaft (DFG, German Research Foundation) under Germany's Excellence Strategy - EXC 2075 - 390740016 and funding by the BMBF under contract 05M20VSA as well as support from the Studienstiftung des deutschen Volkes (German national Academic Foundation).
Part of the work was inspired and motivated by a research stay of the author with Prof.\ Lorenzo Rosasco at MaLGa, DIBRIS, Universit`a di Genova.
 
 
 

\IfFileExists{/home/wenzel/references.bib}{
\bibliography{/home/wenzel/references}			%
}{
\bibliography{/home/math/wenzel/references}			%
}
\bibliographystyle{abbrv} 
 
 
\appendix 

\section{Further proofs} \label{sec:proofs}

\subsection{Proof of \Cref{th:utility_theorem_new}}

The proof of \Cref{th:utility_theorem_new} is constructive, but a bit technical.
For enhanced readability, the most important steps of the proof are visualized in \Cref{fig:vis_proof}.

\begin{figure}[t]
\setlength\fwidth{.7\textwidth}
%
\begin{tikzpicture}

\definecolor{darkgray176}{RGB}{176,176,176}
\definecolor{darkviolet1910191}{RGB}{191,0,191}
\definecolor{gray}{RGB}{128,128,128}

\begin{axis}[
width=0.75\fwidth,
height=0.75\fwidth,
at={(0\fwidth,0\fwidth)},
tick align=outside,
tick pos=left,
x grid style={darkgray176},
xmin=-0.02, xmax=1.02,
xtick style={color=black},
y grid style={darkgray176},
ymin=-0.02, ymax=1.02,
ytick style={color=black},
legend style={at={(0.03,0.03)}, anchor=south west, legend cell align=left, align=left, draw=white!15!black}
]
\addplot [semithick, red, mark=*, mark size=2.5, mark options={solid}, only marks]
table {%
0.16 0.6
0.7 0.3
0.85 0.85
};
\addlegendentry{$Y_0$}
\addplot [semithick, darkviolet1910191, mark=*, mark size=2, mark options={solid}, only marks]
table {%
0.996246337890625 0.989500245558773
0.1885986328125 0.112787684804146
0.216583251953125 0.903994987214009
0.926422119140625 0.250588494301343
0.49468994140625 0.498602855255804
0.01495361328125 0.520245897474978
0.58428955078125 0.0354112686074277
0.61822509765625 0.876746430930244
0.9354248046875 0.624650713813951
0.2484130859375 0.374536401971244
0.27093505859375 0.640247929685515
0.4156494140625 0.205151653711324
0.679443359375 0.648529187623838
0.06268310546875 0.736066656505614
0.8076171875 0.436671239140375
0.4246826171875 0.971701468272113
0.0450439453125 0.304831580551745
0.4803466796875 0.704669003708784
0.7752685546875 0.11436264797033
0.6463623046875 0.287202154143169
0.83905029296875 0.827262104354011
};
\addlegendentry{$Y_1'$}
\addplot [line width=0.24pt, black]
table {%
0.5 1
0.476209042088129 0.999433669591504
0.452471978347909 0.997735961286542
0.428842580863357 0.994910720940466
0.405374377819795 0.990964348631353
0.382120532245286 0.985905784161771
0.359133721579285 0.979746486807249
0.336466018341289 0.972500409357334
0.314168772169836 0.964183966508036
0.292292493499057 0.954815997677259
0.270886739136295 0.944417724327462
0.25 0.933012701892219
0.229679591272201 0.92062676641559
0.209971545214401 0.907287976025168
0.190920506889697 0.893026547371394
0.172569633027358 0.877874787177129
0.154960494258944 0.861867019052535
0.138132980947465 0.845039505741056
0.122125212822871 0.827430366972643
0.106973452628606 0.809079493110303
0.0927120239748321 0.790028454785599
0.0793732335844095 0.770320408727799
0.0669872981077808 0.75
0.0555822756725383 0.729113260863705
0.0451840023227408 0.707707506500943
0.0358160334919637 0.685831227830164
0.0274995906426658 0.663533981658711
0.0202535131927513 0.640866278420715
0.0140942158382292 0.617879467754714
0.0090356513686467 0.594625622180205
0.00508927905953366 0.571157419136643
0.00226403871345771 0.547528021652091
0.000566330408496019 0.523790957911871
0 0.5
0.000566330408496019 0.476209042088129
0.00226403871345765 0.452471978347909
0.0050892790595336 0.428842580863358
0.00903565136864665 0.405374377819795
0.0140942158382291 0.382120532245286
0.0202535131927513 0.359133721579285
0.0274995906426658 0.336466018341289
0.0358160334919637 0.314168772169836
0.0451840023227407 0.292292493499057
0.0555822756725383 0.270886739136295
0.0669872981077806 0.25
0.0793732335844094 0.229679591272201
0.0927120239748321 0.209971545214401
0.106973452628606 0.190920506889698
0.122125212822871 0.172569633027358
0.138132980947465 0.154960494258944
0.154960494258944 0.138132980947465
0.172569633027357 0.122125212822871
0.190920506889697 0.106973452628606
0.209971545214401 0.0927120239748322
0.229679591272201 0.0793732335844094
0.25 0.0669872981077808
0.270886739136295 0.0555822756725383
0.292292493499056 0.045184002322741
0.314168772169836 0.0358160334919638
0.336466018341289 0.0274995906426658
0.359133721579285 0.0202535131927513
0.382120532245287 0.0140942158382291
0.405374377819795 0.0090356513686467
0.428842580863357 0.00508927905953366
0.452471978347908 0.00226403871345771
0.476209042088129 0.000566330408496019
0.5 0
0.523790957911871 0.000566330408496019
0.547528021652091 0.00226403871345771
0.571157419136642 0.0050892790595336
0.594625622180205 0.00903565136864659
0.617879467754713 0.0140942158382291
0.640866278420715 0.0202535131927513
0.663533981658711 0.0274995906426657
0.685831227830164 0.0358160334919637
0.707707506500943 0.0451840023227407
0.729113260863705 0.0555822756725382
0.75 0.0669872981077805
0.770320408727799 0.0793732335844093
0.790028454785599 0.0927120239748321
0.809079493110303 0.106973452628606
0.827430366972643 0.122125212822871
0.845039505741056 0.138132980947465
0.861867019052535 0.154960494258944
0.877874787177129 0.172569633027357
0.893026547371394 0.190920506889697
0.907287976025168 0.209971545214401
0.92062676641559 0.229679591272201
0.933012701892219 0.25
0.944417724327462 0.270886739136295
0.954815997677259 0.292292493499057
0.964183966508036 0.314168772169836
0.972500409357334 0.336466018341289
0.979746486807249 0.359133721579285
0.985905784161771 0.382120532245286
0.990964348631353 0.405374377819795
0.994910720940466 0.428842580863357
0.997735961286542 0.452471978347908
0.999433669591504 0.476209042088129
1 0.5
1 1
0.5 1
};
\addplot [semithick, gray, mark=*, mark size=0.25, mark options={solid}, only marks]
table {%
0.996246337890625 0.989500245558773
0.49468994140625 0.498602855255804
0.494384765625 0.00152415790275873
0.0008544921875 0.500736676319667
0.494171142578125 0.999339531575471
0.992919921875 0.495808565767414
0.24700927734375 0.248539348676523
0.741943359375 0.253467459228776
0.24615478515625 0.751308235533201
0.74237060546875 0.742010872326373
0.744384765625 0.996037189452827
0.24786376953125 0.500127013158563
0.4957275390625 0.750190519737845
0.49530029296875 0.247015190773764
0.99761962890625 0.744144693390235
0.7486572265625 0.497332723670172
0.86871337890625 0.870344967738658
0.37103271484375 0.873139257227049
0.37017822265625 0.376060559874003
0.37164306640625 0.628461108570848
0.61822509765625 0.876746430930244
0.874267578125 0.622008840115836
0.36932373046875 0.126098663821572
0.12188720703125 0.377584717776762
0.1268310546875 0.623278971701468
0.868408203125 0.373266270385612
0.613433837890625 0.372216294941489
0.62213134765625 0.613371945333537
0.618438720703125 0.128892953309963
0.87054443359375 0.995529136818574
0.4981689453125 0.625819234872733
0.744659423828125 0.865975715084082
0.8692626953125 0.745160798658741
0.620147705078125 0.744652746024488
0.249053955078125 0.62531118223848
0.12420654296875 0.749072803942488
0.49298095703125 0.874155362495554
0.24871826171875 0.875273078290911
0.37237548828125 0.50652847635015
0.6180419921875 0.998679063150942
0.2484130859375 0.374536401971244
0.368682861328125 0.250063506579282
0.8709716796875 0.494792460498908
0.123077392578125 0.499110907890057
0.865264892578125 0.248082101305695
0.7369384765625 0.376619417771681
0.36846923828125 0.75283239343596
0.4923095703125 0.378295991464716
0.99945068359375 0.864451557181324
0.623779296875 0.493065081542448
0.4940185546875 0.126708326982675
0.74737548828125 0.617639587461261
0.124786376953125 0.24738776270555
0.99420166015625 0.62099273484733
0.24957275390625 0.124574505918813
0.621002197265625 0.252722315365205
0.741485595703125 0.134379921759894
0.680206298828125 0.679418787786415
0.5579833984375 0.812122135853274
0.06201171875 0.563328760859625
0.0640869140625 0.312147538484987
0.31304931640625 0.44190418127318
0.80718994140625 0.68378804044099
0.8052978515625 0.31351928059747
0.80755615234375 0.935782147030432
0.935821533203125 0.925417873291673
0.431304931640625 0.186726278175752
0.810455322265625 0.557249064336399
0.682037353515625 0.804467476163864
0.6829833984375 0.0713813951125337
0.30657958984375 0.690240308896002
0.30908203125 0.311842706904435
0.93316650390625 0.809988314789412
0.185821533203125 0.814306762180562
0.436676025390625 0.685430743958407
0.431884765625 0.440481633897272
0.93499755859375 0.682771935172484
0.552978515625 0.310166133211401
0.434600830078125 0.315297464817355
0.931640625 0.314128943758573
0.55682373046875 0.1878270588833
0.18426513671875 0.688716150993243
0.55615234375 0.688157293095565
0.551788330078125 0.0642686582329929
0.43487548828125 0.0620840319057054
0.179443359375 0.315195854290504
0.06109619140625 0.440989686531525
0.187408447265625 0.188250436078511
0.559326171875 0.438195397043134
0.8076171875 0.436671239140375
0.30712890625 0.937357110196616
0.432373046875 0.941624752324341
0.68670654296875 0.934257989127673
0.307769775390625 0.81163101830683
0.433258056640625 0.570814069670951
0.31146240234375 0.0647259056038205
0.93267822265625 0.561245745059188
0.679229736328125 0.31362089112432
0.18304443359375 0.439973581263019
0.55841064453125 0.553624955545394
0.429718017578125 0.812782604277803
0.928375244140625 0.432268116310183
0.55194091796875 0.934054768073972
0.80609130859375 0.808972209520906
0.308380126953125 0.566546427543227
0.682281494140625 0.55572490643364
0.1824951171875 0.557435350302291
0.308624267578125 0.192755169435554
0.062042236328125 0.688106487832139
0.67694091796875 0.193314027333232
0.804840087890625 0.193889820318718
0.684844970703125 0.441650154956053
0.43255615234375 0.750596961845247
0.432037353515625 0.248911920608308
0.6810302734375 0.87247878880252
0.99859619140625 0.557587766092567
0.8677978515625 0.807346441091297
0.2447509765625 0.43834781283341
0.67523193359375 0.37824518620129
0.80902099609375 0.999187115785195
0.558868408203125 0.615404155870548
0.555694580078125 0.997053294721333
0.801300048828125 0.374892038815221
0.6187744140625 0.678402682517909
0.37030029296875 0.691459635218209
0.371490478515625 0.934782976849735
0.87493896484375 0.932124168063811
0.247833251953125 0.311402394621416
0.926422119140625 0.250588494301343
0.31121826171875 0.50490270792054
0.805450439453125 0.872377178275669
0.866363525390625 0.310945147250588
0.065673828125 0.624295076969974
0.87152099609375 0.556571660824061
0.681304931640625 0.742281833731308
0.93023681640625 0.747802672356856
0.68341064453125 0.998069399989839
0.18743896484375 0.376568612508256
0.80828857421875 0.74678656708835
0.370147705078125 0.189097190468933
0.618682861328125 0.805619062134837
0.185089111328125 0.752121119748006
0.931427001953125 0.993158224525394
0.430206298828125 0.123863232230859
0.802734375 0.133058984910837
0.49798583984375 0.560382055580958
0.060546875 0.373113854595336
0.616119384765625 0.937983708445528
0.74530029296875 0.80257074632932
0.4935302734375 0.810750393740791
0.30865478515625 0.380937865162831
0.68621826171875 0.616013819031652
0.625152587890625 0.432572947890735
0.24847412109375 0.187573032566174
0.309844970703125 0.626835340141239
0.3746337890625 0.0643194634964182
0.74029541015625 0.316465985876137
0.9967041015625 0.681146166742874
0.558624267578125 0.748310724991109
0.62127685546875 0.0671137529848092
0.43463134765625 0.502260834222425
0.99200439453125 0.435655133871869
0.747955322265625 0.927619434706769
0.744903564453125 0.436230926857356
0.5584716796875 0.124422090128537
0.242462158203125 0.561906213483717
0.55535888671875 0.248386932886247
0.1856689453125 0.255448864502362
0.875030517578125 0.685210587816898
0.6170654296875 0.190265711527714
0.93084716796875 0.497840776304425
0.81072998046875 0.624041050652848
0.430816650390625 0.875645650222696
0.43133544921875 0.375044454605497
0.24578857421875 0.685058172026622
0.36944580078125 0.566275466138292
0.37481689453125 0.316313570085861
0.7440185546875 0.682263882538231
0.49237060546875 0.186455316770817
0.8056640625 0.253772290809328
0.18218994140625 0.128232484885434
0.189727783203125 0.627749834882894
0.36993408203125 0.813798709546309
0.86517333984375 0.185896458873139
0.612213134765625 0.312469305153347
0.49346923828125 0.308387948991516
0.55328369140625 0.379261291469796
0.186920166015625 0.873511829158834
0.31036376953125 0.129756642788193
0.682525634765625 0.131399346305611
0.3729248046875 0.439465528628766
0.244781494140625 0.814984165692899
0.994781494140625 0.92609527680401
0.499267578125 0.436823654930651
0.1256103515625 0.185947264136565
0.8094482421875 0.498450439465529
0.55743408203125 0.875527104608037
0.304840087890625 0.749445375874274
0.619384765625 0.551592745008383
0.31005859375 0.878829446730681
0.744415283203125 0.192907585225829
0.49041748046875 0.932683025961489
0.120513916015625 0.560754627512744
0.74591064453125 0.559111923995326
0.49462890625 0.690443529949703
0.86676025390625 0.431387491744145
0.12310791015625 0.690951582583956
0.99835205078125 0.806330335822791
0.9324951171875 0.372250165117106
0.9354248046875 0.624650713813951
0.740478515625 0.0632525529644871
0.05853271484375 0.502768886856678
0.30450439453125 0.250469948686684
0.68170166015625 0.25062236447696
0.122589111328125 0.307676675303561
0.1253662109375 0.434689833866789
0.560333251953125 0.496587579806601
0.6907958984375 0.499364934207184
0.43719482421875 0.627445003302342
0.128692626953125 0.805229555115243
0.0048828125 0.438957475994513
0.93206787109375 0.869328862470152
0.184234619140625 0.501617300885705
0.490264892578125 0.0628969161205101
0.646148681640625 0.841047265830073
0.96466064453125 0.775847177767617
0.59130859375 0.841792409693644
0.405517578125 0.535589087029416
0.650909423828125 0.964741147182848
0.4046630859375 0.284001422547376
0.337890625 0.721536351165981
0.340240478515625 0.342190384257142
0.58428955078125 0.0354112686074277
0.96563720703125 0.340547680739725
0.5228271484375 0.965757252451354
0.7095947265625 0.83752476756592
0.279541015625 0.413046791647615
0.467315673828125 0.343341970228116
0.6495361328125 0.709495503734187
0.276885986328125 0.531727887009094
0.52069091796875 0.34146217548138
0.766021728515625 0.0975630408643669
0.84014892578125 0.965452420870802
0.527587890625 0.782045419905502
0.4022216796875 0.408372707412488
0.5244140625 0.278463648834019
0.52813720703125 0.842605293908449
0.09368896484375 0.656404003454758
0.216583251953125 0.903994987214009
0.96417236328125 0.527104608037393
0.647857666015625 0.647174380599163
0.52532958984375 0.159376111365137
0.90228271484375 0.465731849819641
0.835113525390625 0.903537739843181
0.84423828125 0.656149977137631
0.533233642578125 0.649545292892344
0.840484619140625 0.464580263848668
0.965576171875 0.956713915561652
0.5867919921875 0.40608647055835
0.965118408203125 0.83762637809277
0.094635009765625 0.78023336551
0.151763916015625 0.840590018459246
0.40240478515625 0.660773256109333
0.96795654296875 0.712035766905451
0.214630126953125 0.344324205321005
0.6463623046875 0.287202154143169
0.461700439453125 0.0945994004978916
0.9029541015625 0.77991159884164
0.653564453125 0.908245694253925
0.40325927734375 0.717675151145659
0.837921142578125 0.591932124168064
0.15484619140625 0.589137834679673
0.1524658203125 0.468830970888584
0.2158203125 0.407864654778235
0.654541015625 0.524157902758726
0.22119140625 0.15546410608139
0.3375244140625 0.8471269623533
0.21728515625 0.470050297210791
0.0902099609375 0.473098613016308
0.89971923828125 0.345424986028553
0.900177001953125 0.400565631932802
0.46832275390625 0.593710308387949
0.4638671875 0.844078646547782
0.46240234375 0.972107910379515
0.40093994140625 0.59736828735457
0.70989990234375 0.219732764314383
0.966461181640625 0.585903232908263
0.95556640625 0.464106081390032
0.462890625 0.277091906721536
0.71630859375 0.101051668952904
0.71099853515625 0.162830869278057
0.403350830078125 0.907890057409948
0.521575927734375 0.0954292198005047
0.8394775390625 0.342783112330437
0.09112548828125 0.531219834374841
0.65313720703125 0.101864553167708
0.58941650390625 0.0939389320733628
0.153472900390625 0.655252417483785
0.027496337890625 0.532710122101983
0.34149169921875 0.0952090636589951
0.092681884765625 0.407119510914664
0.5281982421875 0.52314179749022
0.65106201171875 0.583092008332063
0.466217041015625 0.468390658605565
0.590606689453125 0.222239157310031
0.96722412109375 0.891276736269877
0.212677001953125 0.28945452082169
0.77008056640625 0.402123660011177
0.34027099609375 0.962150078748158
0.0926513671875 0.277599959355789
0.338897705078125 0.781689783061525
0.779083251953125 0.965723382275737
0.033355712890625 0.59855374350116
0.904693603515625 0.713932496739996
0.21368408203125 0.0977493268302596
0.654205322265625 0.773984318108689
0.404205322265625 0.218428762553134
0.3365478515625 0.276482243560433
0.5863037109375 0.158969669257735
0.708282470703125 0.342884722857288
0.34124755859375 0.596352182086064
0.400054931640625 0.779318870768345
0.0374755859375 0.407254991617132
0.27783203125 0.904435299497028
0.0927734375 0.347508001828989
0.775146484375 0.717116293247981
0.1551513671875 0.722044403800234
0.3995361328125 0.153939948178631
0.28045654296875 0.65853782451862
0.091064453125 0.72306050906874
0.2130126953125 0.782197835695778
0.841339111328125 0.715084082710969
0.718994140625 0.649748513946045
0.58709716796875 0.966976578773561
0.462188720703125 0.412843570593913
0.2821044921875 0.476045318294975
0.77362060546875 0.221663364324544
0.646728515625 0.162017985063253
0.277435302734375 0.59337160663178
0.39752197265625 0.0303815475283239
0.52203369140625 0.214652237971854
0.5311279296875 0.585327439922776
0.34039306640625 0.0358685159782553
0.593017578125 0.646700198140527
0.464752197265625 0.536672932649156
0.462982177734375 0.214770783586513
0.712982177734375 0.770326339142068
0.7147216796875 0.963928262968043
0.8968505859375 0.222273027485648
0.032989501953125 0.470524479669427
0.83392333984375 0.778489051465732
0.1539306640625 0.346898338667886
0.280517578125 0.0911446425849718
0.342529296875 0.661789361377839
0.7750244140625 0.587867703094041
0.2808837890625 0.163084895595184
0.09478759765625 0.59218615048519
0.652587890625 0.0413046791647615
0.15985107421875 0.224152822232383
0.521728515625 0.902758725803993
0.27484130859375 0.278108011990042
0.27435302734375 0.843875425494081
0.27587890625 0.344764517604024
0.836700439453125 0.279784585683077
0.39617919921875 0.473809886704263
0.460113525390625 0.718352554657996
0.3494873046875 0.472387339328354
0.523223876953125 0.0356822300123626
0.8341064453125 0.843926230757506
0.77630615234375 0.34318955443784
0.457550048828125 0.782299446222628
0.77642822265625 0.84108113600569
0.46221923828125 0.900980541584108
0.20794677734375 0.58873139257227
0.95794677734375 0.403546207387085
0.71380615234375 0.71355992480821
0.46380615234375 0.158004369252655
0.900543212890625 0.528594895764534
0.2196044921875 0.846415688665346
0.8985595703125 0.843316567596403
0.7161865234375 0.894426662602245
0.58270263671875 0.906823146878016
0.904052734375 0.898186252095717
0.339111328125 0.222831885383326
0.589111328125 0.778387440938881
0.835235595703125 0.159071279784586
0.59246826171875 0.467865670883504
0.277374267578125 0.785347762028146
0.71990966796875 0.467662449829802
0.21966552734375 0.654575013971447
0.8980712890625 0.589696692577351
0.40069580078125 0.846110857084794
0.907257080078125 0.955901031346847
0.77044677734375 0.650459787633999
0.5804443359375 0.279682975156226
0.775390625 0.462277091906722
0.964263916015625 0.655404833274061
0.835723876953125 0.406052600382733
0.46234130859375 0.0311944317431286
0.530670166015625 0.466104421751427
0.837158203125 0.219631153787532
0.20892333984375 0.222933495910176
0.21099853515625 0.718386424833613
0.33795166015625 0.53457298176091
0.71356201171875 0.286795712035767
0.336669921875 0.409388812680994
0.582763671875 0.339429964944368
0.89556884765625 0.275923385662755
0.15338134765625 0.526952192247117
0.77435302734375 0.288319869938526
0.401123046875 0.349032159731748
0.77978515625 0.531778692272519
0.649322509765625 0.339904147403004
0.704620361328125 0.405070365289844
0.782012939453125 0.905298988975258
0.71844482421875 0.590407966265305
0.58935546875 0.578417924096936
0.77276611328125 0.774018188284306
0.5308837890625 0.718640451150739
0.33935546875 0.91175125743027
0.15252685546875 0.276990296194686
0.036651611328125 0.34059848600315
0.8424072265625 0.522024081694864
0.5850830078125 0.714982472184118
0.15667724609375 0.410709749530051
0.219482421875 0.527358634354519
0.2686767578125 0.719097698521567
0.33380126953125 0.162678453487781
0.469970703125 0.659960371894528
0.394561767578125 0.966417720875883
0.27740478515625 0.216328811664888
0.718536376953125 0.5313383799895
0.15386962890625 0.775694761977341
0.1007080078125 0.217040085352843
0.652618408203125 0.4672560077224
0.03179931640625 0.664126403495402
0.90576171875 0.657978966620942
0.524444580078125 0.404460702128741
0.6456298828125 0.21566834324036
0.7686767578125 0.163542142966011
0.406494140625 0.0941929583904892
0.1556396484375 0.155057663973988
0.5897216796875 0.519483818523599
0.6439208984375 0.400599502108418
0.3695068359375 0.283798201493675
0.776611328125 0.998780673677793
0.96392822265625 0.989229284153838
0.34112548828125 0.308997612152619
0.746856689453125 0.962979898050771
0.999420166015625 0.589561211874884
0.551910400390625 0.346610442175143
0.936279296875 0.71528730376467
0.503570556640625 0.659977306982337
0.867706298828125 0.901641010008637
0.93212890625 0.529949702789209
0.99951171875 0.526291723822588
0.555419921875 0.964944368236549
0.12042236328125 0.465376212975664
0.40570068359375 0.629883655946756
0.929931640625 0.592135345221765
0.61767578125 0.713763145861911
0.36767578125 0.158207590306356
0.525390625 0.684499314128944
0.843414306640625 0.624312012057783
0.80303955078125 0.405781638977798
0.64971923828125 0.995630747345425
0.679443359375 0.648529187623838
0.77813720703125 0.620383071686227
0.367431640625 0.40695016003658
0.24261474609375 0.279022506731697
0.6513671875 0.807041609510745
0.4013671875 0.25148605395519
0.494537353515625 0.841504513200901
0.39886474609375 0.748158309200833
0.37286376953125 0.722349235380785
0.963653564453125 0.806601297227726
0.92926025390625 0.283239343595997
0.80035400390625 0.10216938474826
0.30975341796875 0.345577401818828
0.961334228515625 0.494453758742739
0.525054931640625 0.00154109299056716
0.6112060546875 0.281715185693238
0.34063720703125 0.879642330945486
0.2451171875 0.342021033379058
0.594757080078125 0.437382512828329
0.313507080078125 0.474419549865366
0.5537109375 0.0333790580704161
0.094970703125 0.437738149672306
0.467681884765625 0.629341733136886
0.058868408203125 0.282070822537215
0.80865478515625 0.714271198496164
0.49615478515625 0.529086013310979
0.901153564453125 0.312774136733899
0.772247314453125 0.131704177886162
0.030059814453125 0.371759047570662
0.775177001953125 0.252417483784653
0.461822509765625 0.376941184440041
0.28082275390625 0.445562160239801
0.12298583984375 0.338159833358736
0.0609130859375 0.654371792917746
0.900665283203125 0.99126149469085
0.9981689453125 0.959152568206066
0.677001953125 0.842402072854748
0.688507080078125 0.585530660976477
0.3074951171875 0.779657572524513
0.864044189453125 0.839455367576081
0.77606201171875 0.805314230554285
0.996978759765625 0.776118139172552
0.800811767578125 0.163948585073414
0.838775634765625 0.748683296922895
0.743804931640625 0.347220105336246
0.713531494140625 0.999559687716981
0.27960205078125 0.127317990143779
0.985595703125 0.466544734034446
0.465118408203125 0.504293044759437
0.15924072265625 0.625717624345882
0.58453369140625 0.675557587766093
0.3128662109375 0.656912056089011
0.867279052734375 0.217057020440651
0.871337890625 0.461057765584515
0.3377685546875 0.373621907229589
0.8079833984375 0.589899913631052
0.588409423828125 0.125234974343342
0.836029052734375 0.809649613033243
0.4920654296875 0.968043489305492
0.250885009765625 0.404376026689698
0.524169921875 0.372351775643957
0.279815673828125 0.623177361174618
0.87164306640625 0.961794441904181
0.03509521484375 0.634456129655032
0.300811767578125 0.719504140628969
0.21844482421875 0.257074632931972
0.367919921875 0.903215973174821
0.705169677734375 0.375264610747007
0.73956298828125 0.222577859066199
0.62078857421875 0.969008789310572
0.6510009765625 0.870446578265508
0.272216796875 0.563023929279073
0.618927001953125 0.342952463208522
0.24493408203125 0.591576487324087
0.962188720703125 0.746176903927247
0.680938720703125 0.968399126149469
0.34014892578125 0.632119087537469
0.8074951171875 0.224102016968958
0.06402587890625 0.594624803129604
0.4329833984375 0.404714728445867
0.49908447265625 0.590052329421328
0.74365234375 0.836305441243713
0.277496337890625 0.310487899879761
0.275543212890625 0.936002303171942
0.33807373046875 0.81745668851293
0.5562744140625 0.21749733272367
0.341064453125 0.500838286846517
0.182037353515625 0.471134142830531
0.683685302734375 0.902013581940422
0.903594970703125 0.812020525326424
0.67413330078125 0.414012091652695
0.7498779296875 0.898084641568866
0.5869140625 0.871056241426612
0.248291015625 0.092059137326627
0.5264892578125 0.747904282883707
0.5267333984375 0.614591271655743
0.5926513671875 0.610933292689123
0.45928955078125 0.813189046385205
0.70977783203125 0.679825229893817
0.45977783203125 0.124269674338261
0.19171142578125 0.157496316618402
0.24798583984375 0.78260427780318
0.812652587890625 0.654795170112957
0.77349853515625 0.681349387796576
0.2479248046875 0.217243306406544
0.778228759765625 0.430439126826873
0.28094482421875 0.874358583549255
0.56072998046875 0.84626327287507
0.81256103515625 0.527714271198496
0.528350830078125 0.13011227963217
0.49169921875 0.71970736168267
0.05908203125 0.534064929126657
0.030426025390625 0.561973953834951
0.36700439453125 0.843062541279276
0.405670166015625 0.0586970143440194
0.715972900390625 0.560602211722468
0.716461181640625 0.808125455130485
0.839141845703125 0.685820250978001
0.654815673828125 0.437992175989432
0.620635986328125 0.0996291215769954
0.649566650390625 0.246016020593067
0.4613037109375 0.936747447035513
0.154541015625 0.190824569425392
0.220489501953125 0.618672627817575
0.89935302734375 0.436468018086674
0.803375244140625 0.284119968162035
0.554351806640625 0.155176209588647
0.58526611328125 0.996240410506528
0.404052734375 0.564852918762384
0.61968994140625 0.906010262663212
0.99798583984375 0.893715388914291
0.31182861328125 0.100492811055225
0.80615234375 0.465935070873342
0.497802734375 0.466087486663618
0.340484619140625 0.131246930515335
0.36871337890625 0.537011634405324
0.7191162109375 0.932632220698064
0.212982177734375 0.436993005808735
0.18505859375 0.656607224508459
0.562255859375 0.467459228776101
0.4332275390625 0.157597927145252
0.837554931640625 0.371911463360937
0.6832275390625 0.713153482700808
0.7755126953125 0.876848041457095
0.183746337890625 0.28579654185507
0.21929931640625 0.373317075649037
0.2535400390625 0.470355128791343
0.2774658203125 0.691053193110806
0.529296875 0.55281207133059
0.875274658203125 0.651069450795102
0.12005615234375 0.528374739623025
0.67962646484375 0.347762028146116
0.398956298828125 0.815221256922217
0.00396728515625 0.471472844586699
0.093017578125 0.313366864807194
0.5640869140625 0.64548087181832
0.437286376953125 0.469609984927772
0.40032958984375 0.937153889142915
0.33551025390625 0.690646751003404
0.64227294921875 0.370929228268049
0.48638916015625 0.781486562007824
0.965362548828125 0.621805619062135
0.68304443359375 0.773306914596352
0.43304443359375 0.217751359040797
0.841705322265625 0.996206540330911
0.613037109375 0.039475689681451
0.249176025390625 0.90765296618063
0.9002685546875 0.558807092414774
0.21551513671875 0.194888990499416
0.42608642578125 0.0340395264949449
0.3089599609375 0.843468983386679
0.3702392578125 0.0935832952293858
0.15185546875 0.689529035208047
0.09808349609375 0.624701519077376
0.58111572265625 0.371945333536554
0.25341796875 0.154549611339735
0.67547607421875 0.162424427170655
0.217041015625 0.68465172991922
0.55633544921875 0.782451862012904
0.212890625 0.499314128943759
0.434234619140625 0.279395078663483
0.811676025390625 0.969381361242358
0.710357666015625 0.252112652204102
0.435211181640625 0.845162492167522
0.839111328125 0.556165218716659
0.65264892578125 0.74323019864858
0.40264892578125 0.187674643093024
0.30828857421875 0.413453233755017
0.93719482421875 0.960778336635675
0.905059814453125 0.927314603126217
0.53076171875 0.43575674439872
0.213897705078125 0.559467560839303
0.367462158203125 0.784128435705939
0.33441162109375 0.749987298684144
0.522125244140625 0.308811326186726
0.781982421875 0.490321597317482
0.0294189453125 0.502362444749276
0.46051025390625 0.246202306558959
0.907958984375 0.616521871665905
0.86248779296875 0.777676167250927
0.9666748046875 0.920947010110247
0.375 0.222222222222222
0.625 0.777777777777778
0.65142822265625 0.618858913783468
0.18060302734375 0.786262256769801
0.65185546875 0.133973479652492
0.009185791015625 0.410777489881285
0.55047607421875 0.903165167911396
0.751495361328125 0.528527155413301
0.8724365234375 0.524056292231875
0.401031494140625 0.444004132161425
0.09454345703125 0.684854950972921
0.779541015625 0.746380124980948
0.59197998046875 0.253670680282477
0.8055419921875 0.77645684092872
0.14862060546875 0.554996697657877
0.59307861328125 0.0634557740181883
0.156890869140625 0.808887534081864
0.25189208984375 0.657318498196413
0.369140625 0.0301783264746228
0.64776611328125 0.551795966062084
0.89764404296875 0.74495757760504
0.58642578125 0.549154092363969
0.871612548828125 0.341970228115633
0.67401123046875 0.281105522532134
0.30853271484375 0.280546664634456
0.275146484375 0.383782959914647
0.187896728515625 0.529661806296466
0.276641845703125 0.74754864603973
0.776641845703125 0.191993090484174
0.71881103515625 0.503022913173805
0.2181396484375 0.74765025656658
0.647613525390625 0.681315517620959
0.397613525390625 0.125759962065403
0.8370361328125 0.931717725956409
0.4620361328125 0.746532540771224
0.154205322265625 0.440650984775356
0.3092041015625 0.594726413656455
0.9359130859375 0.781943809378651
0.937164306640625 0.34447662111128
0.6500244140625 0.069349184575522
0.5250244140625 0.810089925316263
0.18194580078125 0.34405324391607
0.4246826171875 0.971701468272113
0.13092041015625 0.213585327439923
0.304443359375 0.537418076512727
0.151885986328125 0.309505664786872
0.777587890625 0.55982319768328
0.7447509765625 0.771681146166743
0.5533447265625 0.0967840268251791
0.62249755859375 0.460549712950262
0.43994140625 0.657521719250114
0.58740234375 0.194330132601738
0.33648681640625 0.562617487171671
0.685302734375 0.527815881725347
0.6629638671875 0.495706955240563
0.5902099609375 0.806431946349642
0.310791015625 0.906873952141442
0.61590576171875 0.402022049484327
0.77337646484375 0.372453386170807
0.121612548828125 0.658842656099172
0.90283203125 0.497027892089621
0.37054443359375 0.662195803485241
0.408935546875 0.689224203627496
0.5181884765625 0.0638622161255906
0.09259033203125 0.560483666107809
0.40289306640625 0.875374688817761
0.681060791015625 0.102135514572643
0.4337158203125 0.542905044962658
0.2166748046875 0.809835898999136
0.6871337890625 0.471726870903826
0.52630615234375 0.993395315754712
0.712921142578125 0.0734136056495453
0.18634033203125 0.412335517959661
0.12738037109375 0.593608697861098
0.68133544921875 0.0417111212721638
0.179779052734375 0.84668665007028
0.182891845703125 0.722857288015038
0.120391845703125 0.278412843570594
0.5924072265625 0.744246303917086
0.83734130859375 0.315145049027079
0.4276123046875 0.785144540974445
0.555938720703125 0.523954681705025
0.8671875 0.275720164609053
0.4921875 0.0905349794238683
0.218994140625 0.316415180612711
0.1495361328125 0.376162170400853
0.96527099609375 0.554742671340751
0.51806640625 0.933241883859168
0.46441650390625 0.87171670985114
0.40521240234375 0.311639485850734
0.8662109375 0.403749428440786
0.9970703125 0.716506630086877
0.587646484375 0.939338515470203
0.83624267578125 0.252146522379718
0.80499267578125 0.844739114972311
0.49365234375 0.280749885688157
0.180572509765625 0.599163406662264
0.3458251953125 0.444749276024996
0.71087646484375 0.742823756541178
0.46087646484375 0.187268200985622
0.963897705078125 0.374282375654118
0.71368408203125 0.431082660163593
0.805419921875 0.347660417619265
0.4051513671875 0.499822181578011
0.957122802734375 0.431048789987976
0.773529052734375 0.315822452539416
0.52044677734375 0.872682009856221
0.643585205078125 0.936696641772087
0.522491455078125 0.250825585530661
0.3424072265625 0.18869074836153
0.371826171875 0.343545191281817
0.4373779296875 0.342529086013311
0.9306640625 0.401920438957476
0.707427978515625 0.189927009771546
0.990631103515625 0.648088875340819
0.929595947265625 0.8370844552829
0.4947509765625 0.216125590611187
0.712677001953125 0.622787854155024
0.74151611328125 0.409134786363867
0.680328369140625 0.22041016782672
0.49603271484375 0.0336330843875425
0.842071533203125 0.87603515724229
0.404693603515625 0.380599163406662
0.310455322265625 0.223915731003065
0.619171142578125 0.221561753797693
0.77862548828125 0.938627241782249
0.864654541015625 0.587732222391573
0.83843994140625 0.189960879947163
0.707672119140625 0.867872444918627
0.45953369140625 0.564446476654981
0.468414306640625 0.439126826872597
0.433380126953125 0.714694575691375
0.130035400390625 0.406509847753561
0.808380126953125 0.89987976087656
0.435333251953125 0.0891801723991939
0.89532470703125 0.373469491439313
0.434722900390625 0.597639248759505
0.437042236328125 0.910328710054362
0.066680908203125 0.413758065335569
0.063018798828125 0.345475791291978
0.46124267578125 0.0669613371945334
0.2469482421875 0.535487476502566
0.71209716796875 0.127470405934055
0.6195068359375 0.83935375704923
0.45941162109375 0.305542854239699
0.2401123046875 0.711070466900371
0.095947265625 0.251181222374638
0.92919921875 0.460448102423411
0.56298828125 0.582075903063557
0.463043212890625 0.689088722925028
0.274200439453125 0.24274754864604
0.908721923828125 0.687649240461312
0.93060302734375 0.897373367880912
0.5302734375 0.187014174668496
0.03143310546875 0.439770360209318
0.208740234375 0.128638926992836
0.61956787109375 0.64710664024793
0.2735595703125 0.0655387898186252
0.905517578125 0.868922420362749
0.49322509765625 0.337651780724483
0.2481689453125 0.848041457094955
0.12109375 0.77914951989026
0.08740234375 0.749885688157293
0.58575439453125 0.312198343748412
0.0650634765625 0.718183203779912
0.96136474609375 0.686429914139105
0.11761474609375 0.723466951176142
0.3707275390625 0.602042371589697
0.489898681640625 0.153804467476164
0.739898681640625 0.709360023031719
0.962646484375 0.865264441396129
0.15106201171875 0.24975867499873
0.624786376953125 0.580721096038883
0.557159423828125 0.717827566935935
0.307159423828125 0.162272011380379
0.277008056640625 0.817727649917865
0.6256103515625 0.519280597469898
0.4678955078125 0.998933089468069
0.8350830078125 0.492760249961896
0.086090087890625 0.379075005503904
0.74359130859375 0.290453691002388
0.273681640625 0.505715592135345
0.61541748046875 0.154905248183712
0.493438720703125 0.906670731087741
0.184112548828125 0.214398211654727
0.49761962890625 0.410811360056902
0.5540771484375 0.274399227759996
0.828277587890625 0.128046198919541
0.21392822265625 0.878118173042727
0.6419677734375 0.312808006909516
0.9310302734375 0.650256566580298
0.8736572265625 0.719554945892395
0.83819580078125 0.438703449677387
0.348907470703125 0.0671645582482345
0.0865478515625 0.498704465782655
0.3402099609375 0.250876390794086
0.99468994140625 0.831936188589138
0.061553955078125 0.477163034090332
0.36444091796875 0.971091805111009
0.52545166015625 0.497535944723873
0.705596923828125 0.317278870090941
0.156036376953125 0.494301342952463
0.54998779296875 0.407305796880557
0.89410400390625 0.250317532896408
0.7364501953125 0.160798658741046
0.7427978515625 0.585124218869075
0.75067138671875 0.465122186658538
0.158355712890625 0.746701891649308
0.74371337890625 0.648122745516435
0.284454345703125 0.187877864146726
0.598876953125 0.496723060509069
0.88299560546875 0.764873240867754
0.64727783203125 0.18599806939999
0.74102783203125 0.0872326373012244
0.37310791015625 0.468729360361733
0.7037353515625 0.91434232586496
0.334228515625 0.93979576284103
0.08203125 0.230452674897119
0.7430419921875 0.111619163745364
0.599273681640625 0.692899117681925
0.79095458984375 0.661433724533862
0.385894775390625 0.26842114176362
0.512481689453125 0.765381293502007
0.6015625 0.893004115226337
0.2314453125 0.111568358481939
0.5712890625 0.0809327846364883
0.885284423828125 0.885180104658842
0.29095458984375 0.328100391200528
0.573638916015625 0.8282443394469
0.6337890625 0.821673525377229
0.699615478515625 0.980050466561669
0.7713623046875 0.509424376365392
0.946533203125 0.762841030330742
0.726470947265625 0.269183220714999
0.75994873046875 0.700858608951887
0.297760009765625 0.459245711189012
0.485260009765625 0.607393859337161
0.50726318359375 0.789767819946146
0.896240234375 0.634811766499009
0.48004150390625 0.575420413554844
0.98095703125 0.793324188385917
0.388092041015625 0.106250740910092
0.187255859375 0.578570339887212
0.94805908203125 0.603922166336432
0.665191650390625 0.937374045284425
0.3843994140625 0.797744246303917
0.76226806640625 0.270436417212823
0.88726806640625 0.418584565360971
0.72662353515625 0.854188893969415
0.763427734375 0.980490778844688
0.485626220703125 0.643296245491033
0.7930908203125 0.826855662246609
0.8865966796875 0.295432606818066
0.44671630859375 0.525021592236956
0.286407470703125 0.363460854544531
0.599029541015625 0.0198309878236719
0.415435791015625 0.830530576301038
0.5616455078125 0.667428745618046
0.656982421875 0.26809937509526
0.385284423828125 0.551846771325509
0.4150390625 0.463648834019204
0.6954345703125 0.275415333028502
0.2933349609375 0.0533455265965554
0.6680908203125 0.604633440024387
0.19964599609375 0.838693288624701
0.51312255859375 0.701976324747244
0.91937255859375 0.331605954376873
0.4788818359375 0.827008078036884
0.9805908203125 0.975003810394757
0.5081787109375 0.982014936747447
0.607208251953125 0.595353011905367
0.44464111328125 0.983894731494183
0.3839111328125 0.581923487273281
0.912384033203125 0.760808819793731
0.7877197265625 0.269623532998019
0.9127197265625 0.417771681146167
0.078857421875 0.576741350403902
0.98486328125 0.57384545038866
0.45458984375 0.609967992684042
0.1737060546875 0.392826296804349
0.606689453125 0.858862978204542
0.57354736328125 0.732865924909821
0.10870361328125 0.421480465376213
0.82476806640625 0.863029009805416
0.509918212890625 0.734356212636962
0.018341064453125 0.580263848668055
0.417236328125 0.517299192196312
0.104736328125 0.295076969974089
0.604736328125 0.628410303307423
0.4189453125 0.331961591220851
0.75787353515625 0.330183407000965
0.19976806640625 0.30747345424986
0.69976806640625 0.640806787583194
0.2618408203125 0.328913275415333
0.63226318359375 0.0490270792054057
0.229644775390625 0.391877931887077
0.947509765625 0.576284103033074
0.9456787109375 0.512879134278311
0.419708251953125 0.484241900794256
0.3551025390625 0.733729614388051
0.637939453125 0.118122237463801
0.137939453125 0.673677793019357
0.383819580078125 0.392115023116395
0.0224609375 0.391403749428441
0.574951171875 0.240664532845603
0.2352294921875 0.418432149570696
0.51885986328125 0.634100492811055
0.169708251953125 0.706464123016478
0.0799560546875 0.639739877051262
0.4735107421875 0.0498907686836356
0.94744873046875 0.849006757100035
0.200592041015625 0.328472963132314
0.420074462890625 0.645192975325577
0.916656494140625 0.481041169198462
0.634521484375 0.88995579942082
0.726806640625 0.913122999542753
0.19915771484375 0.453386170807296
0.69915771484375 0.786719504140629
0.633453369140625 0.0228793036291893
0.506256103515625 0.952615624312012
0.355499267578125 0.324442412233907
0.447540283203125 0.36163186506122
0.8218994140625 0.39033683889651
0.352294921875 0.113092516384697
0.6361083984375 0.268912259310065
0.976806640625 0.394604481024234
0.539306640625 0.86374028349337
0.226806640625 0.579789666209419
0.981842041015625 0.604193127741367
0.57275390625 0.143118427069044
0.4432373046875 0.423004623278972
0.824737548828125 0.572422903012752
0.762603759765625 0.825500855221934
0.11090087890625 0.397093938932073
0.289154052734375 0.575081711798676
0.9820556640625 0.853071178174059
0.505767822265625 0.362461684363833
0.388214111328125 0.365289844027841
0.139678955078125 0.822842046436011
0.390533447265625 0.701281986147098
0.386810302734375 0.295246320852174
0.663970947265625 0.565479517011296
0.20172119140625 0.391606970482142
0.536285400390625 0.887991329235042
0.6903076171875 0.177462785144541
0.5653076171875 0.918203525885282
0.8214111328125 0.174516079865874
0.44482421875 0.961133973479652
0.79315185546875 0.63501498755271
0.3819580078125 0.142966011278768
0.569366455078125 0.0532947213331301
0.6319580078125 0.698521566834324
0.97735595703125 0.417365239038764
0.69598388671875 0.29776964893563
0.231475830078125 0.636285119138343
0.574310302734375 0.110061135666988
0.1063232421875 0.486104760453183
0.08306884765625 0.609256718996088
0.41741943359375 0.612813087435858
0.10284423828125 0.514149265863944
0.289794921875 0.55204999237921
0.664794921875 0.663161103490322
0.295196533203125 0.641467256007722
0.51611328125 0.450388660265203
0.23583984375 0.140832190214906
0.820465087890625 0.918169655709665
0.445465087890625 0.73298447052448
0.199981689453125 0.481430676218056
0.351806640625 0.802011888431642
0.537506103515625 0.228335788921066
0.855224609375 0.607376924249352
0.327423095703125 0.615861403241376
0.698272705078125 0.146573184981964
0.2305908203125 0.863892699283646
0.2930908203125 0.493522328913275
0.481475830078125 0.0190011685210588
0.387115478515625 0.424494911006114
0.2496337890625 0.730986130163085
0.952178955078125 0.637656861250826
0.44622802734375 0.201900116852106
0.69622802734375 0.757455672407661
0.16839599609375 0.542396992328405
0.67169189453125 0.822486409592034
0.28692626953125 0.261443885586547
0.84942626953125 0.298480922623584
0.16876220703125 0.361123812426967
0.85430908203125 0.357008586089519
0.57196044921875 0.420311944317431
0.79150390625 0.513488797439415
0.22906494140625 0.2269979169842
0.7197265625 0.36076817558299
0.41131591796875 0.954630899761215
0.476348876953125 0.319632847296313
0.351348876953125 0.764077291740758
0.478271484375 0.235634811766499
0.56927490234375 0.170350048265
0.353271484375 0.680079256210943
0.728271484375 0.791190367322055
0.661651611328125 0.229487374892039
0.041046142578125 0.579586445155718
0.858551025390625 0.673085064946062
0.858795166015625 0.638943927924266
0.200347900390625 0.358956121187488
0.884552001953125 0.795627360327863
0.66473388671875 0.791596809429457
0.29388427734375 0.0756998425036834
0.41473388671875 0.236041253873901
0.4176025390625 0.585581466239902
0.855621337890625 0.94011752950939
0.644927978515625 0.486223306067842
0.44366455078125 0.023065589595082
0.91058349609375 0.291368185744043
0.94329833984375 0.480363765686125
0.7659912109375 0.85855814662399
0.478302001953125 0.363528594895765
0.043975830078125 0.488136970990195
0.446044921875 0.107605547934766
0.26171875 0.79835390946502
0.54296875 0.761316872427984
0.3232421875 0.7946959304984
0.47747802734375 0.399430981049637
0.97747802734375 0.73276431438297
0.567535400390625 0.29539873664245
0.947357177734375 0.943165845314908
0.660552978515625 0.3216142525699
0.474029541015625 0.79760876560145
0.388580322265625 0.613490490948196
0.980224609375 0.7555250723975
0.480224609375 0.422191739064167
0.3914794921875 0.327897170146827
0.41668701171875 0.422598181171569
0.60418701171875 0.53370929228268
0.72918701171875 0.978153736727125
0.598907470703125 0.919016410100086
0.33270263671875 0.0549712950261647
0.38153076171875 0.488441802570746
0.79022216796875 0.448458060255042
0.47772216796875 0.263272875069857
0.35272216796875 0.707717319514302
0.54022216796875 0.670680282477265
0.20440673828125 0.859828278209622
0.41534423828125 0.736371488086166
0.577911376953125 0.481955663940118
0.376617431640625 0.95490186116615
0.10980224609375 0.361327033480669
0.69635009765625 0.234771122288269
0.164794921875 0.329827770156988
0.700592041015625 0.661806296465647
0.13482666015625 0.70741248793375
0.63385009765625 0.629832850683331
0.38336181640625 0.644922013920642
0.88336181640625 0.978255347253975
0.91363525390625 0.382004775694762
0.54266357421875 0.508103439516334
0.2908935546875 0.516283086927806
0.22821044921875 0.728953919626073
0.757476806640625 0.755999254856137
0.48077392578125 0.483970939389321
0.98077392578125 0.817304272722654
0.47589111328125 0.292536706802825
0.386474609375 0.520957171162932
0.35736083984375 0.642686582329929
0.85736083984375 0.976019915663263
0.54052734375 0.326017375400091
0.53839111328125 0.699944114210232
0.60675048828125 0.568256871411878
0.6654052734375 0.0823553320123965
0.3529052734375 0.860133109790174
0.5404052734375 0.823096072753137
0.7904052734375 0.600873850530915
0.16693115234375 0.511913834273231
0.60211181640625 0.953563989229284
0.821563720703125 0.729715998577453
0.786407470703125 0.696794187877864
0.723907470703125 0.178275669359346
0.294586181640625 0.795779776118139
0.57574462890625 0.763349082964995
0.32574462890625 0.20779352740944
0.105987548828125 0.757608088197937
0.2847900390625 0.766651425087639
0.724395751953125 0.392335179257904
0.91961669921875 0.971752273535538
0.70086669921875 0.601381903165168
0.38311767578125 0.857084793984657
0.62774658203125 0.859066199258243
0.706451416015625 0.482565327101221
0.72650146484375 0.446527460244881
0.736846923828125 0.514809734288472
0.6331787109375 0.142508763907941
0.51104736328125 0.140273332317228
0.38604736328125 0.918051110095006
0.823638916015625 0.606022117224678
0.475616455078125 0.13971447441955
0.698638916015625 0.0875035987061593
0.725616455078125 0.695270029975105
0.7039794921875 0.455469186607733
0.726348876953125 0.875188402851869
0.1351318359375 0.292028654168572
0.888580322265625 0.946823824281529
0.260040283203125 0.608545445308134
0.9119873046875 0.361276228217243
0.4744873046875 0.916831783772799
0.848297119140625 0.793798370844553
0.1834716796875 0.235533201239648
0.7254638671875 0.199410658944267
0.56951904296875 0.950718894477468
0.788970947265625 0.787701739233518
0.13507080078125 0.574505918813189
0.878997802734375 0.233517925790445
0.3248291015625 0.144795000762079
0.23187255859375 0.670426256160138
0.232666015625 0.487120865721689
0.853515625 0.857338820301783
0.671661376953125 0.630103812088266
0.572662353515625 0.261257599620654
0.66119384765625 0.300614743687446
0.862945556640625 0.701129570356822
0.820220947265625 0.326796389439279
0.66448974609375 0.118528679571204
0.16448974609375 0.674084235126759
0.639923095703125 0.726972514352487
0.827423095703125 0.949194736574709
0.389923095703125 0.171416958796931
0.168243408203125 0.825280699080424
0.540313720703125 0.951938220799675
0.321563720703125 0.396382665244119
0.292022705078125 0.924350962759742
0.917022705078125 0.516943555352334
0.85137939453125 0.386272417822486
0.53887939453125 0.015902047452116
0.482330322265625 0.51472505884943
0.8511962890625 0.540313976527968
0.473388671875 0.952598689224204
0.598388671875 0.174820911446426
0.3511962890625 0.206980643194635
0.6011962890625 0.762536198750191
0.75701904296875 0.728496672255246
0.82501220703125 0.422852207488696
0.88153076171875 0.821775135904079
0.20355224609375 0.608240613727582
0.66400146484375 0.885484936239394
0.138671875 0.762688614540466
0.385498046875 0.886755067825027
0.879638671875 0.199512269471117
0.35626220703125 0.58334603464919
0.361572265625 0.489864349946654
0.6116943359375 0.477213839353757
0.8231201171875 0.761824925062236
0.31988525390625 0.328506833307931
0.230255126953125 0.179715151823062
0.31793212890625 0.734542498602855
0.51483154296875 0.423969923284052
0.48284912109375 0.108011990042168
0.825958251953125 0.669427085979441
0.294952392578125 0.671950414062897
0.22711181640625 0.768378804044099
0.20086669921875 0.268048569831835
0.3262939453125 0.0826093583295229
0.327056884765625 0.579959017087504
0.161163330078125 0.570441497739166
0.64190673828125 0.600569018950363
0.25885009765625 0.51872173957222
0.509307861328125 0.114261037443479
0.6322021484375 0.663973987705126
0.40966796875 0.79652491998171
0.569976806640625 0.978221477078359
0.819976806640625 0.360937526461075
0.171875 0.641975308641975
0.3614501953125 0.0496875476299345
0.38818359375 0.236854138088706
0.939605712890625 0.265220410167827
0.76220703125 0.447645176040238
0.63818359375 0.792409693644262
0.259033203125 0.577655845145557
0.8546142578125 0.821978356957781
0.7947998046875 0.575267997764568
0.821044921875 0.292790733119951
0.97479248046875 0.508814713204288
0.60369873046875 0.330488238581517
0.1658935546875 0.294060864705584
0.543121337890625 0.56974715913902
0.260986328125 0.764212772443225
0.34930419921875 0.42442717065488
0.94891357421875 0.730325661738556
0.88092041015625 0.324696438551034
0.1378173828125 0.639536655997561
0.699737548828125 0.0539043844942336
0.5050048828125 0.166285627190977
0.91888427734375 0.853477620281461
0.73724365234375 0.944012599705329
0.85357666015625 0.571610018797948
0.2628173828125 0.861758878219783
0.32470703125 0.706904435299497
0.44970703125 0.262459990855053
0.237060546875 0.516384697454656
0.41973876953125 0.672966519331403
0.264892578125 0.486206370980033
0.1431884765625 0.174973327236702
0.6431884765625 0.508306660570035
0.260284423828125 0.107402326881065
0.44537353515625 0.891225931006452
0.818389892578125 0.14931666920693
0.756744384765625 0.946214161120425
0.41845703125 0.855052583447645
0.69970703125 0.818015546410608
0.63525390625 0.98262459990855
0.4493408203125 0.049077884468831
0.13800048828125 0.481837118325458
0.8243408203125 0.234263069654016
0.9241943359375 0.699436061575979
0.574737548828125 0.794645125234974
0.293487548828125 0.831682162272011
0.074737548828125 0.461311791901641
0.324737548828125 0.239089569679419
0.50616455078125 0.825534725397551
0.3958740234375 0.0754966214499822
0.257476806640625 0.422665921522803
0.108428955078125 0.575928466189097
0.198272705078125 0.70212874053752
0.43878173828125 0.769293298785754
0.260162353515625 0.890887229250284
0.732452392578125 0.486765228877712
0.540069580078125 0.20692983793121
0.66455078125 0.697302240512117
0.41455078125 0.141746684956562
0.78863525390625 0.233856627546614
0.5394287109375 0.290656912056089
0.85205078125 0.919524462734339
0.47705078125 0.734339277549154
0.54150390625 0.735711019661637
0.31976318359375 0.826804856983183
0.918212890625 0.543362292333486
0.261962890625 0.358177107148301
0.82373046875 0.45496113397348
0.10882568359375 0.70395773002083
0.45184326171875 0.455519991871158
0.85888671875 0.509830818472794
0.60748291015625 0.662144998221816
0.91851806640625 0.912411725854799
0.04644775390625 0.517807244830565
0.89019775390625 0.480770207793527
0.29644775390625 0.295585022608342
0.732666015625 0.820454199055022
0.14239501953125 0.511100950058426
0.448394775390625 0.861013734356213
0.754058837890625 0.391420684516249
0.350982666015625 0.832359565784349
0.04278564453125 0.615353350607123
0.97039794921875 0.445003302342123
0.76165771484375 0.548036376568612
0.27093505859375 0.640247929685515
0.292510986328125 0.235431590712798
0.20263671875 0.645633287608596
0.57318115234375 0.697099019458416
0.82318115234375 0.474876797236194
0.5963134765625 0.360158512421887
0.44659423828125 0.143778895493573
0.69659423828125 0.699334451049129
0.92254638671875 0.668140019306
0.76312255859375 0.479754102525022
0.3917236328125 0.680688919372047
0.172027587890625 0.165083235956578
0.75469970703125 0.367982522989382
0.28839111328125 0.144388558654677
0.15045166015625 0.608647055834985
0.599700927734375 0.391725516096801
0.7525634765625 0.23670172229843
0.2603759765625 0.26550830666057
0.697998046875 0.11446425849718
0.94677734375 0.548239597622314
0.11260986328125 0.202001727378956
0.42498779296875 0.296194685769446
0.95050048828125 0.296651933140273
0.8485107421875 0.235075953868821
0.88177490234375 0.392572270487222
0.68951416015625 0.393283544175177
0.507843017578125 0.232535690697556
0.645294189453125 0.419702281156328
0.8203125 0.209876543209877
0.41748046875 0.269775948788294
0.3203125 0.765432098765432
0.727813720703125 0.672102829853173
0.200469970703125 0.799674846314078
0.950469970703125 0.910785957425189
0.231719970703125 0.207082253721486
0.54412841796875 0.538993039678911
0.66912841796875 0.983437484123355
0.5093994140625 0.0199664685261393
0.41705322265625 0.359599654524209
0.353607177734375 0.362918931734661
0.350860595703125 0.550017781842199
0.8839111328125 0.915256820606615
0.510528564453125 0.485613642906738
0.323028564453125 0.522650679943775
0.725372314453125 0.304543684059002
0.38330078125 0.833104709647919
0.698150634765625 0.954444613795323
0.98297119140625 0.700248945790784
0.91265869140625 0.733170756490372
0.82659912109375 0.984555199918711
0.38909912109375 0.206777422140934
0.63909912109375 0.762332977696489
0.326690673828125 0.425646496977087
0.357940673828125 0.919473657470914
0.629058837890625 0.169198462294027
0.23486328125 0.462734339277549
0.657012939453125 0.386780470456739
0.7689208984375 0.919118020626937
0.07464599609375 0.394248844180257
0.23333740234375 0.830158004369253
0.98333740234375 0.941269115480364
0.79339599609375 0.209063658995072
0.109527587890625 0.609527680401023
0.573883056640625 0.627055496282748
0.073883056640625 0.293722162949415
0.57757568359375 0.888533252044912
0.116485595703125 0.640552761266067
0.8504638671875 0.421632881166489
0.5379638671875 0.0512625107961185
0.5106201171875 0.391454554691866
0.237091064453125 0.893021050314146
0.60919189453125 0.081745668851293
0.81976318359375 0.271249301427628
0.0511474609375 0.670629477213839
0.7308349609375 0.757049230300259
0.5447998046875 0.797490219986791
0.66107177734375 0.724533861708073
0.47357177734375 0.76157089874511
0.04901123046875 0.392216633643245
0.5745849609375 0.0163084895595184
0.6697998046875 0.0567494792460499
0.170654296875 0.14327084285932
0.4808349609375 0.201493674744704
0.6307373046875 0.385967586241935
0.01568603515625 0.614134024284916
0.41107177734375 0.168978306152517
0.56793212890625 0.401209165269522
0.789459228515625 0.14877474639706
0.512969970703125 0.19473657470914
0.59991455078125 0.726769293298786
0.823486328125 0.702484377381497
0.9482421875 0.387288523090992
0.5107421875 0.856424325560128
0.7607421875 0.634202103337906
0.54119873046875 0.626784534877813
0.53729248046875 0.977950515673424
0.04681396484375 0.360107707158462
0.105499267578125 0.54666463445613
0.6739501953125 0.457094955037342
0.14044189453125 0.324544022760758
0.9150390625 0.796982167352538
0.296875 0.864197530864197
0.197174072265625 0.671103659672475
0.41900634765625 0.897322562617487
0.8846435546875 0.701468272112991
0.417999267578125 0.768886856678352
0.633087158203125 0.236803332825281
0.34381103515625 0.391911802062694
0.729644775390625 0.72521126522041
0.66497802734375 0.3623939440126
0.479644775390625 0.169655709664855
0.10675048828125 0.234923538078545
0.483184814453125 0.89027756608918
0.7283935546875 0.331097901742621
0.35723876953125 0.228522074886958
0.60723876953125 0.784077630442514
0.19744873046875 0.737895645988924
0.4803466796875 0.704669003708784
0.32415771484375 0.675608393029518
0.44915771484375 0.231163948585073
0.91790771484375 0.26820098562211
0.825836181640625 0.635285948957645
0.54522705078125 0.485342681501804
0.23272705078125 0.300157496316618
0.42022705078125 0.0779352740943962
0.20147705078125 0.892750088909211
0.605743408203125 0.417873291673017
0.69500732421875 0.425798912767363
0.85076904296875 0.893105725753188
0.41326904296875 0.11532794797541
0.353515625 0.52400548696845
0.29296875 0.205761316872428
0.295166015625 0.116750495351318
0.0450439453125 0.304831580551745
0.36163330078125 0.302900980541584
0.82574462890625 0.541126860742773
0.75250244140625 0.665498145607885
0.321197509765625 0.360480279090247
0.723541259765625 0.142373283205473
0.598541259765625 0.98187945604498
0.97918701171875 0.360869786109841
0.0733642578125 0.546258192348727
0.6595458984375 0.203068637910888
0.16644287109375 0.264390590865214
0.112091064453125 0.448576605869701
0.956298828125 0.356805365035818
0.8804931640625 0.604328608443835
0.856353759765625 0.726735423123169
0.928131103515625 0.944385171637115
0.69952392578125 0.885891378346797
0.13848876953125 0.730579688055683
0.6651611328125 0.544886450236244
0.887725830078125 0.451099933953158
0.4161376953125 0.395366559975613
0.0184326171875 0.453182949753594
0.6298828125 0.327846364883402
0.04754638671875 0.551541939744957
0.32666015625 0.297210791037952
0.07666015625 0.519433013260174
0.73291015625 0.605852766346594
0.85009765625 0.480566986739826
0.35009765625 0.147233653406493
0.38238525390625 0.772951277752375
0.5098876953125 0.60524310318549
0.82354736328125 0.510643702687598
0.6973876953125 0.358329522938576
0.132965087890625 0.362614100154109
0.33026123046875 0.466290707717319
0.072784423828125 0.329624549103287
0.734649658203125 0.631865061220342
0.23687744140625 0.358685159782553
0.34857177734375 0.613422750596962
0.33038330078125 0.895493573134177
0.04974365234375 0.462531118223848
0.23724365234375 0.610679266371996
0.573150634765625 0.510000169350878
0.6063232421875 0.819438093786516
0.3563232421875 0.263882538230961
0.257720947265625 0.289759352402242
0.637115478515625 0.0911615776727802
0.20526123046875 0.244068485495097
0.70526123046875 0.577401818828431
0.32586669921875 0.490270792054057
0.538604736328125 0.264238175074938
0.539703369140625 0.171027451777337
0.291748046875 0.892242036274958
0.1663818359375 0.456637707666514
0.78399658203125 0.389930396789107
0.16778564453125 0.763501498755271
0.91326904296875 0.448661281308744
0.854400634765625 0.452387000626598
0.06903076171875 0.266219580348524
0.56903076171875 0.599552913681857
0.910430908203125 0.574251892496063
0.5421142578125 0.45160798658741
0.7218017578125 0.237616217040085
0.8507080078125 0.328151196463954
0.8150634765625 0.829294314891023
0.50811767578125 0.079307016206879
0.795684814453125 0.329235042083693
0.881866455078125 0.571813239851649
0.9490966796875 0.888025199410659
0.45068359375 0.829446730681298
0.31866455078125 0.862571762434588
0.78851318359375 0.86384189402022
0.88189697265625 0.857542041355484
0.53680419921875 0.387390133617843
0.078887939453125 0.695422445765381
0.453887939453125 0.917644667987603
0.9808349609375 0.534827008078037
0.6668701171875 0.144540974444952
0.882598876953125 0.35666988433335
0.509765625 0.322359396433471
0.81597900390625 0.793527409439618
0.63970947265625 0.454249860285526
0.95220947265625 0.824620230655896
0.54595947265625 0.602398008433674
0.45220947265625 0.491286897322563
0.825592041015625 0.884028518687869
0.537139892578125 0.0752425951328558
0.662353515625 0.85337600975461
0.0494384765625 0.32312147538485
0.07574462890625 0.430015749631662
0.7603759765625 0.598841639993903
0.669830322265625 0.477688021812393
0.888458251953125 0.663940117529509
0.26129150390625 0.229741401209165
0.418365478515625 0.918322071499941
0.543365478515625 0.140544293722163
0.79034423828125 0.921556673271351
0.53485107421875 0.113041711121272
0.8466796875 0.173296753543667
0.79443359375 0.422039323273891
0.203155517578125 0.179037748310725
0.7496337890625 0.175430574607529
0.566436767578125 0.331301122796322
0.78729248046875 0.36066656505614
0.34991455078125 0.17121373774323
0.60406494140625 0.115886805873088
0.29547119140625 0.613829192704364
0.450103759765625 0.640315670036749
0.950103759765625 0.973649003370082
0.67572021484375 0.918406746938983
0.952880859375 0.32754153330285
0.762847900390625 0.420684516249217
0.75384521484375 0.152974648173551
0.0797119140625 0.67428745618046
0.60784912109375 0.515419397449576
0.197906494140625 0.764991786482413
0.4130859375 0.0210333790580704
0.90966796875 0.240969364426155
0.792999267578125 0.954072041863537
0.7880859375 0.304983996342021
0.57867431640625 0.854798557130519
0.760986328125 0.20865721688767
0.135986328125 0.541990550221003
0.941650390625 0.44459686023472
0.7962646484375 0.546004166031601
0.7960205078125 0.986587410455723
0.10748291015625 0.328811664888482
0.20062255859375 0.516791139562059
0.28961181640625 0.398008433673729
0.78961181640625 0.731341767007062
0.477294921875 0.557536960829142
0.23004150390625 0.797642635777067
0.98004150390625 0.908753746888178
0.139892578125 0.263984148757811
0.4588623046875 0.324239191180206
0.0118408203125 0.551135497637555
0.266845703125 0.676421277244322
0.141845703125 0.4541990550221
0.329864501953125 0.116615014648851
0.6024169921875 0.208555606360819
0.757598876953125 0.307287168283968
0.230621337890625 0.0882656776575386
0.353546142578125 0.949956815526088
0.521514892578125 0.569069755626683
0.78350830078125 0.175328964080679
0.28350830078125 0.730884519636234
0.851348876953125 0.208521736185202
0.630401611328125 0.953767210282985
0.6380615234375 0.532286744906772
0.7901611328125 0.767108672458467
0.02264404296875 0.485698318345781
0.50958251953125 0.268302596148961
0.355560302734375 0.889667902928077
0.574615478515625 0.535606022117225
0.482818603515625 0.676895459702959
0.577545166015625 0.564869853850192
0.70123291015625 0.513996850073668
0.1126708984375 0.676319666717472
0.2669677734375 0.201696895798405
0.27886962890625 0.109028095310674
0.38836669921875 0.0458263476096124
0.173614501953125 0.486985385019221
0.486114501953125 0.857355755389592
0.9769287109375 0.677488187776254
0.416534423828125 0.705481887923589
0.915557861328125 0.632779555961998
0.2884521484375 0.706497993192095
0.72467041015625 0.415231417974902
0.98260498046875 0.870954630899761
0.3251953125 0.183813443072702
0.572509765625 0.354061880810852
0.6732177734375 0.510338871107047
0.3626708984375 0.45409744449525
0.6317138671875 0.298176091043032
0.944244384765625 0.699300580873512
0.576416015625 0.211400701112635
0.69580078125 0.20347508001829
0.223663330078125 0.274145201442869
0.912750244140625 0.827329844705245
0.513092041015625 0.513658148317499
0.758819580078125 0.57730020830158
0.542755126953125 0.364900337008247
0.72943115234375 0.548950871310268
0.44818115234375 0.585987908347305
0.459625244140625 0.395231079273146
0.678375244140625 0.950786634828702
0.26434326171875 0.439059086521364
0.5994873046875 0.139054005995021
0.22747802734375 0.917949499568155
0.14404296875 0.42615454961134
0.263580322265625 0.169046046503751
0.701690673828125 0.536757608088198
0.05224609375 0.705989940557842
0.317047119140625 0.262934173313689
0.01495361328125 0.520245897474978
0.69512939453125 0.855408220291622
0.81573486328125 0.115581974292537
0.1666259765625 0.240816948635879
0.68609619140625 0.329878575420414
0.57745361328125 0.45851750241325
0.0220947265625 0.356043286084438
0.26544189453125 0.54676624498298
0.511993408203125 0.549560534471371
0.7752685546875 0.11436264797033
0.267974853515625 0.144659520059611
0.954498291015625 0.670036749140544
0.887359619140625 0.546884790597639
0.44024658203125 0.797337804196515
0.65301513671875 0.824518620129045
0.22900390625 0.550525834476452
0.15118408203125 0.394045623126556
0.1656494140625 0.427373875933547
0.4156494140625 0.205151653711324
0.6656494140625 0.76070720926688
0.13555908203125 0.789107351521617
0.61004638671875 0.445917797083778
0.381011962890625 0.985385019221324
0.20391845703125 0.536706802824772
0.501861572265625 0.889210655557249
0.2020263671875 0.820809835898999
0.791534423828125 0.890667073108774
0.3223876953125 0.642280140222527
0.8597412109375 0.759792714525225
0.1075439453125 0.798658741045572
0.045074462890625 0.422970753103355
0.885650634765625 0.7322223915731
0.755126953125 0.0769699740893156
0.50738525390625 0.291469796270894
0.668365478515625 0.584988738166607
0.692291259765625 0.833731307896831
0.95184326171875 0.788853325204491
0.54833984375 0.424782807498857
0.2613525390625 0.832495046486816
0.32147216796875 0.554082202916222
0.16961669921875 0.860641162424427
0.25830078125 0.388660265203475
0.87652587890625 0.261291469796271
0.755279541015625 0.785263086589104
0.146575927734375 0.206540330911616
0.969482421875 0.63846974546563
0.51861572265625 0.668241629832851
0.297637939453125 0.176903927246863
0.449859619140625 0.674456807058545
0.316497802734375 0.04833274060526
0.943023681640625 0.42129417941032
0.5052490234375 0.921861504851903
0.35540771484375 0.0830158004369253
0.531524658203125 0.922674389066707
0.31719970703125 0.923538078544937
0.1717529296875 0.20261139054006
0.699310302734375 0.554505580111433
0.22607421875 0.448102423411065
0.261260986328125 0.926789615404156
0.1395263671875 0.228217243306407
0.7330322265625 0.892394452065234
0.48095703125 0.459990855052583
0.42041015625 0.0502972107910379
0.99969482421875 0.664482040339379
0.63232421875 0.924096936442615
0.36065673828125 0.390692475740487
0.99981689453125 0.699029619468577
0.095611572265625 0.45711189012515
0.89129638671875 0.51450490270792
0.921905517578125 0.882741452014429
0.489288330078125 0.545750139714474
0.9783935546875 0.997764568409287
0.045440673828125 0.647868719199309
0.925811767578125 0.608393029517858
0.891876220703125 0.68033328252807
0.17718505859375 0.615556571660824
0.7652587890625 0.891479957323579
0.1885986328125 0.112787684804146
0.60296630859375 0.241070974953005
0.51092529296875 0.999898389473149
0.26531982421875 0.0842351267591322
0.452667236328125 0.548188792358888
0.235748291015625 0.262629341733137
0.173736572265625 0.802790902470829
0.104766845703125 0.377178275669359
0.664947509765625 0.175295093905062
0.6336669921875 0.356703754508967
0.80267333984375 0.482192755169436
0.58587646484375 0.594675608393029
0.97698974609375 0.488899049941574
0.201446533203125 0.4192450337855
0.867950439453125 0.477315449880608
0.70172119140625 0.724940303815475
0.59912109375 0.293095564700503
0.45172119140625 0.16938474825992
0.58001708984375 0.389828786262257
0.10552978515625 0.734847330183407
0.0777587890625 0.484072549916171
0.41259765625 0.986739826245999
0.545806884765625 0.839218276346763
0.795806884765625 0.616996054124541
0.7340087890625 0.570492303002591
0.222686767578125 0.596047350505512
0.21539306640625 0.702535182644922
0.453521728515625 0.702501312469305
0.27130126953125 0.458974749784078
0.9951171875 0.453132144490169
0.083251953125 0.360920591373266
0.76422119140625 0.354569933445105
0.29876708984375 0.426865823299294
0.83905029296875 0.827262104354011
0.7459716796875 0.272570238276685
0.5494384765625 0.656454808718183
0.9476318359375 0.995732357872276
0.232177734375 0.326169791190367
0.9134521484375 0.595386882080984
0.9285888671875 0.767311893512168
0.63079833984375 0.20875882741452
0.67889404296875 0.5721180714322
0.914581298828125 0.999796778946299
0.1993408203125 0.56759640298735
0.508087158203125 0.681247777269725
0.09136962890625 0.331250317532896
0.729034423828125 0.0758522582939593
0.30682373046875 0.521160392216634
0.7276611328125 0.116902911141594
0.56890869140625 0.313417670070619
0.667633056640625 0.336246168436383
0.385040283203125 0.756693593456282
0.201690673828125 0.203424274754865
0.420166015625 0.561194939795763
0.173004150390625 0.735727954749445
0.8704833984375 0.293603617334756
0.20831298828125 0.629375603312503
0.58148193359375 0.658080577147793
0.85235595703125 0.269217090890616
0.91680908203125 0.307625870040136
0.47735595703125 0.0840319057054311
0.3897705078125 0.505105928974242
0.6163330078125 0.0236244474927602
0.765380859375 0.179393385154702
0.265380859375 0.734948940710258
0.505645751953125 0.0466561669122254
0.296539306640625 0.0934478145269183
0.480987548828125 0.979830310420159
0.020538330078125 0.422293349591018
0.51446533203125 0.531677081745669
0.15863037109375 0.791139562058629
0.446441650390625 0.381818489728869
0.9215087890625 0.348270080780369
0.629913330078125 0.409947670578672
0.55877685546875 0.363410049281105
0.1826171875 0.770004572473708
0.566497802734375 0.900184592457112
0.26300048828125 0.704059340547681
0.4554443359375 0.761164456637708
0.113800048828125 0.26378092770411
0.343414306640625 0.290978678724449
0.206451416015625 0.149231993767888
0.7584228515625 0.115988416399939
0.97900390625 0.661636945587563
0.89776611328125 0.329573743839862
0.24407958984375 0.171721790377483
0.84600830078125 0.767921556673271
0.53240966796875 0.41827973378042
0.73529052734375 0.461159376111365
0.92950439453125 0.731951430168165
0.0458984375 0.446273433927755
0.796966552734375 0.118291588341886
0.29290771484375 0.379312096733222
0.79290771484375 0.712645430066555
0.968017578125 0.473860691967688
0.672393798828125 0.431895544378398
0.28204345703125 0.85769445714576
0.269073486328125 0.184219885180105
0.630279541015625 0.563040864366882
0.0775146484375 0.731189351216786
0.078643798828125 0.246710359193212
0.489044189453125 0.950566478687192
0.2342529296875 0.696438551033887
0.255279541015625 0.451929753255771
0.517120361328125 0.651983945536757
0.397491455078125 0.732307067012143
0.5943603515625 0.272367017222984
0.632904052734375 0.19236566241596
0.8388671875 0.732967535436671
0.6162109375 0.922267946959305
0.711700439453125 0.946451252349743
0.45025634765625 0.00843367372859828
0.1844482421875 0.831783772798862
0.85272216796875 0.152161763958746
0.57171630859375 0.932428999644363
0.057525634765625 0.637572185811784
0.60223388671875 0.0508560686887161
0.09344482421875 0.701519077376416
0.2255859375 0.243255601280293
0.100830078125 0.642584971803079
0.1951904296875 0.136767769140883
0.32440185546875 0.597977950515673
0.57452392578125 0.441446933902352
0.44757080078125 0.944876289183559
0.599365234375 0.708885840573083
0.092071533203125 0.764924046131179
0.90478515625 0.976223136716964
0.48443603515625 0.441294518112076
0.63507080078125 0.907839252146522
0.558380126953125 0.233213094209893
0.966552734375 0.601889955799421
0.469879150390625 0.612271164625989
0.44879150390625 0.0815932530610171
0.13580322265625 0.3842910125489
0.440704345703125 0.928618604887466
0.165191650390625 0.604040711951091
0.65887451171875 0.923284052227811
0.7100830078125 0.270538027739674
0.614044189453125 0.172788700909414
0.401275634765625 0.267201815441413
0.80535888671875 0.915053599552914
0.43035888671875 0.729868414367728
0.228057861328125 0.126606716455825
0.3448486328125 0.92760249961896
0.32598876953125 0.952801910277905
0.8236083984375 0.19483818523599
0.6207275390625 0.268709038256363
0.44512939453125 0.299852664736067
0.351959228515625 0.408034005656319
0.731597900390625 0.0996968619282291
0.843505859375 0.640298734948941
0.32525634765625 0.880861657267693
0.6771240234375 0.235990448610476
0.50091552734375 0.642229334959102
0.291168212890625 0.343409710579349
0.791168212890625 0.676743043912683
0.118194580078125 0.224762485393487
0.168731689453125 0.18513437992176
0.449462890625 0.506325255296449
0.636962890625 0.469288218259412
0.553619384765625 0.0490948195566394
0.8890380859375 0.716100187979475
0.14306640625 0.748056698673983
0.6156005859375 0.296347101559722
0.54608154296875 0.917797083777879
0.3857421875 0.350251486053955
0.05706787109375 0.610069603210893
0.157745361328125 0.170502464055276
0.28192138671875 0.428085149621501
0.14080810546875 0.159934969262816
0.54705810546875 0.641416450744297
0.573028564453125 0.189317346610442
0.53131103515625 0.354874765025657
0.489166259765625 0.0477230774441565
0.859893798828125 0.65411776660062
0.5142822265625 0.590611187319006
0.1947021484375 0.780572067266169
0.669097900390625 0.395993158224525
0.563140869140625 0.203949262476926
0.43792724609375 0.484783823604125
0.109161376953125 0.346153194804315
0.399566650390625 0.5793493539264
0.899566650390625 0.912682687259733
0.451934814453125 0.440346153194804
0.918060302734375 0.778556791816965
0.666107177734375 0.869091771240834
0.724700927734375 0.836169960541245
0.310089111328125 0.0854544530813392
0.12432861328125 0.322715033277447
0.38983154296875 0.312858812172941
0.127105712890625 0.450405595353012
0.301422119140625 0.361699605412454
0.88238525390625 0.21739572219682
0.603271484375 0.34674592287761
0.822021484375 0.815881725346746
0.76104736328125 0.806939998983895
0.293701171875 0.27770156988264
0.33416748046875 0.323629528019103
0.417755126953125 0.253789225897136
0.294921875 0.589849108367627
0.3770751953125 0.740435909160189
0.50946044921875 0.716608240613728
0.461456298828125 0.518924960625921
0.42938232421875 0.23238327490728
0.386260986328125 0.186048874663415
0.636260986328125 0.741604430218971
0.078399658203125 0.446679876035157
0.654571533203125 0.950109231316364
0.748321533203125 0.851343799217599
0.998443603515625 0.911463360937526
0.67938232421875 0.787938830462836
0.851959228515625 0.741367338989653
0.387939453125 0.451455570797135
0.8037109375 0.700045724737083
0.03253173828125 0.547071076563532
0.59185791015625 0.826754051719758
0.264923095703125 0.912157699537672
0.38555908203125 0.566885129299395
0.88555908203125 0.900218462632729
0.745635986328125 0.54407356602144
0.59503173828125 0.452420870802215
0.6475830078125 0.56683432403597
0.19952392578125 0.552558045013463
0.7725830078125 0.789056546258192
0.366943359375 0.142356348117665
0.616943359375 0.697911903673221
0.3883056640625 0.666057003505563
0.088287353515625 0.421751426781148
0.776397705078125 0.665091703500483
0.06268310546875 0.736066656505614
0.380950927734375 0.083083540788159
0.9962158203125 0.48117664990093
0.891754150390625 0.649308201663026
0.84820556640625 0.698419956307473
0.4315185546875 0.521770055377737
0.753448486328125 0.9126149469085
0.721435546875 0.0595945739978662
0.62603759765625 0.283544175176548
0.40582275390625 0.890006604684245
0.7545166015625 0.511050144795001
0.354522705078125 0.628054666463446
0.854522705078125 0.961387999796779
0.862335205078125 0.171264543006655
0.402130126953125 0.550085522193433
0.902130126953125 0.883418855526766
0.985504150390625 0.55054276956426
0.510650634765625 0.806296465647174
0.21112060546875 0.112381242696743
0.47613525390625 0.937560331250317
0.60113525390625 0.15978255347254
0.12493896484375 0.8210130569527
0.137725830078125 0.339988822842046
0.9453125 0.654320987654321
0.45166015625 0.346593507087334
0.88580322265625 0.495402123660011
0.864166259765625 0.232908262629342
0.873931884765625 0.666988433335027
0.8536376953125 0.284255448864502
0.85125732421875 0.252959406594523
0.47625732421875 0.067774221409338
0.50750732421875 0.573947060915511
0.2171630859375 0.209927348473302
0.69439697265625 0.684702535182645
0.927734375 0.577503429355281
0.44439697265625 0.129146979627089
0.810089111328125 0.418787786414673
0.8883056640625 0.999390336838896
0.343902587890625 0.486070890277566
0.765777587890625 0.523107927314603
0.961090087890625 0.934630561059459
0.7701416015625 0.0830666057003506
0.737213134765625 0.36185202120273
0.4317626953125 0.26367931717726
0.149810791015625 0.361394773831902
0.781890869140625 0.475554200748531
0.219390869140625 0.512591237785568
0.183837890625 0.300563938424021
0.9798583984375 0.836813493877966
0.6600341796875 0.418889396941523
0.30340576171875 0.142762790225067
0.896728515625 0.828684651729919
0.31964111328125 0.539450287049738
0.693603515625 0.260783417162018
0.774078369140625 0.146336093752646
0.4180908203125 0.93796677335772
0.2476806640625 0.643194634964182
0.791290283203125 0.250520753950109
0.916290283203125 0.398668902098257
0.351654052734375 0.278785415502379
0.18798828125 0.396890717878372
0.578155517578125 0.0679266371996139
0.86212158203125 0.883757557282934
0.42462158203125 0.105979779505157
0.1533203125 0.704160951074531
0.635986328125 0.875323883554336
0.49609375 0.705075445816187
0.729736328125 0.776558451455571
0.479736328125 0.221002895900015
0.276031494140625 0.295855984013277
0.626129150390625 0.0814069670951244
0.480865478515625 0.622025775203644
0.51702880859375 0.471523649850124
0.536041259765625 0.0929905671560907
0.512359619140625 0.620958864671713
0.522216796875 0.22969059594574
0.494903564453125 0.769564260190689
0.48919677734375 0.662805466646345
0.81378173828125 0.954478483970939
0.0477294921875 0.566580297718844
0.228546142578125 0.43143829700757
0.417022705078125 0.183610222019001
0.667022705078125 0.739165777574557
0.418060302734375 0.445223458483632
0.36773681640625 0.858913783467967
0.684967041015625 0.727649917864824
0.434967041015625 0.172094362309269
0.982421875 0.775034293552812
0.524322509765625 0.32344324205321
0.758941650390625 0.994987214008705
0.253692626953125 0.138562888448577
0.325714111328125 0.809734288472286
0.18603515625 0.173754000914495
0.04205322265625 0.532439160697048
0.29205322265625 0.310216938474826
0.35504150390625 0.1309759691104
0.35479736328125 0.341919422852207
0.144683837890625 0.625988585750817
0.652496337890625 0.0265372825958103
0.7333984375 0.927754915409236
0.716552734375 0.824112178021643
0.323272705078125 0.035462073870853
0.691497802734375 0.159443851716371
0.87628173828125 0.510034039526495
0.894073486328125 0.961997662957882
0.87371826171875 0.541939744957577
0.77166748046875 0.953259157648732
0.5841064453125 0.17725956409084
0.848388671875 0.841487578113092
0.4591064453125 0.955037341868617
0.84075927734375 0.606564040034548
0.81842041015625 0.341157343900828
0.67657470703125 0.299395417365239
0.09027099609375 0.295483412081492
0.40277099609375 0.517705634303714
0.59027099609375 0.628816745414825
0.804107666015625 0.178038578130028
};
\addplot [semithick, black]
table {%
0.26 0.6
0.259798667647188 0.606342391965656
0.25919548128308 0.612659245357375
0.258192869726271 0.618925124436041
0.256794870139636 0.625114798718108
0.255007111774095 0.631203344569849
0.252836793301607 0.637166245566033
0.250292653828662 0.642979491208917
0.247384937706979 0.648619673610047
0.244125353283118 0.65406408174556
0.240527025753106 0.659290792905464
0.236604444311898 0.664278760968654
0.232373403810507 0.669007901148211
0.227850941155713 0.673459170865753
0.223055266708452 0.677614646429176
0.21800569095712 0.681457595205034
0.21272254676105 0.684972542994951
0.207227107477268 0.688145336344758
0.201541501300189 0.690963199535452
0.195688622159187 0.693414786026511
0.189692037532827 0.695490224144407
0.183575893550943 0.697181156832354
0.177364817766693 0.698480775301221
0.171083819990101 0.699383846446125
0.164758191582374 0.699886733918301
0.158413403616519 0.699987412767388
0.152075004314321 0.699685477595194
0.145768516172671 0.698982144188093
0.139519333193481 0.697880244621478
0.133352618630997 0.696384215855994
0.127293203668258 0.694500081871467
0.121365487430687 0.692235429410458
0.115593338739423 0.689599377429134
0.11 0.686602540378444
0.104607993613389 0.683256985463477
0.0994390312862333 0.679576184053083
0.0945139266054715 0.675574957435426
0.0898525112293679 0.671269417137886
0.0854735550324245 0.666676900051629
0.0813946905257212 0.66181589862206
0.0776323418570167 0.656705986386277
0.0742016586765023 0.651367739157341
0.0711164551345077 0.645822652172741
0.068389154256793 0.640093053540661
0.0660307379214092 0.634202014332567
0.0640507026385503 0.628173255684143
0.0624570213114593 0.622031053278654
0.0612561111323606 0.615800139597335
0.0604528077426915 0.609505604330418
0.0600503457616815 0.603172793349807
0.0600503457616815 0.596827206650193
0.0604528077426915 0.590494395669582
0.0612561111323606 0.584199860402665
0.0624570213114593 0.577968946721346
0.0640507026385503 0.571826744315857
0.0660307379214091 0.565797985667433
0.068389154256793 0.559906946459339
0.0711164551345077 0.554177347827259
0.0742016586765023 0.548632260842659
0.0776323418570167 0.543294013613723
0.0813946905257213 0.538184101377939
0.0854735550324245 0.533323099948371
0.0898525112293679 0.528730582862114
0.0945139266054715 0.524425042564574
0.0994390312862333 0.520423815946917
0.104607993613389 0.516743014536523
0.11 0.513397459621556
0.115593338739423 0.510400622570866
0.121365487430687 0.507764570589542
0.127293203668258 0.505499918128533
0.133352618630997 0.503615784144006
0.139519333193481 0.502119755378522
0.145768516172671 0.501017855811907
0.152075004314321 0.500314522404806
0.158413403616519 0.500012587232612
0.164758191582374 0.500113266081699
0.171083819990101 0.500616153553875
0.177364817766693 0.501519224698779
0.183575893550943 0.502818843167646
0.189692037532827 0.504509775855593
0.195688622159187 0.506585213973489
0.201541501300189 0.509036800464548
0.207227107477268 0.511854663655242
0.21272254676105 0.515027457005049
0.21800569095712 0.518542404794966
0.223055266708452 0.522385353570824
0.227850941155713 0.526540829134247
0.232373403810507 0.530992098851789
0.236604444311898 0.535721239031346
0.240527025753106 0.540709207094536
0.244125353283118 0.54593591825444
0.247384937706979 0.551380326389953
0.250292653828662 0.557020508791083
0.252836793301607 0.562833754433967
0.255007111774095 0.568796655430151
0.256794870139636 0.574885201281892
0.258192869726271 0.581074875563959
0.25919548128308 0.587340754642625
0.259798667647188 0.593657608034344
0.26 0.6
};
\addplot [semithick, black]
table {%
0.8 0.3
0.799798667647188 0.306342391965656
0.799195481283079 0.312659245357375
0.798192869726271 0.318925124436041
0.796794870139636 0.325114798718108
0.795007111774094 0.331203344569849
0.792836793301607 0.337166245566033
0.790292653828662 0.342979491208917
0.787384937706978 0.348619673610047
0.784125353283118 0.35406408174556
0.780527025753106 0.359290792905464
0.776604444311898 0.364278760968654
0.772373403810507 0.369007901148211
0.767850941155713 0.373459170865753
0.763055266708452 0.377614646429176
0.75800569095712 0.381457595205034
0.75272254676105 0.384972542994951
0.747227107477268 0.388145336344758
0.741541501300189 0.390963199535452
0.735688622159187 0.393414786026511
0.729692037532827 0.395490224144407
0.723575893550943 0.397181156832354
0.717364817766693 0.398480775301221
0.711083819990101 0.399383846446125
0.704758191582374 0.399886733918301
0.698413403616519 0.399987412767387
0.692075004314321 0.399685477595194
0.685768516172671 0.398982144188093
0.679519333193481 0.397880244621478
0.673352618630996 0.396384215855994
0.667293203668258 0.394500081871467
0.661365487430687 0.392235429410458
0.655593338739422 0.389599377429134
0.65 0.386602540378444
0.644607993613389 0.383256985463477
0.639439031286233 0.379576184053083
0.634513926605471 0.375574957435426
0.629852511229368 0.371269417137886
0.625473555032425 0.366676900051629
0.621394690525721 0.36181589862206
0.617632341857017 0.356705986386277
0.614201658676502 0.351367739157341
0.611116455134508 0.345822652172741
0.608389154256793 0.340093053540661
0.606030737921409 0.334202014332567
0.60405070263855 0.328173255684143
0.602457021311459 0.322031053278654
0.601256111132361 0.315800139597335
0.600452807742691 0.309505604330418
0.600050345761681 0.303172793349807
0.600050345761681 0.296827206650193
0.600452807742691 0.290494395669582
0.601256111132361 0.284199860402665
0.602457021311459 0.277968946721346
0.60405070263855 0.271826744315857
0.606030737921409 0.265797985667433
0.608389154256793 0.259906946459339
0.611116455134508 0.254177347827259
0.614201658676502 0.248632260842659
0.617632341857017 0.243294013613723
0.621394690525721 0.238184101377939
0.625473555032425 0.233323099948371
0.629852511229368 0.228730582862114
0.634513926605471 0.224425042564574
0.639439031286233 0.220423815946917
0.644607993613389 0.216743014536523
0.65 0.213397459621556
0.655593338739423 0.210400622570866
0.661365487430687 0.207764570589542
0.667293203668258 0.205499918128533
0.673352618630996 0.203615784144006
0.679519333193481 0.202119755378522
0.685768516172671 0.201017855811907
0.692075004314321 0.200314522404806
0.698413403616519 0.200012587232612
0.704758191582374 0.200113266081699
0.711083819990101 0.200616153553875
0.717364817766693 0.201519224698779
0.723575893550943 0.202818843167646
0.729692037532827 0.204509775855593
0.735688622159187 0.206585213973489
0.741541501300189 0.209036800464548
0.747227107477268 0.211854663655242
0.75272254676105 0.215027457005049
0.75800569095712 0.218542404794966
0.763055266708452 0.222385353570824
0.767850941155713 0.226540829134247
0.772373403810507 0.230992098851789
0.776604444311898 0.235721239031346
0.780527025753106 0.240709207094536
0.784125353283118 0.24593591825444
0.787384937706978 0.251380326389953
0.790292653828662 0.257020508791083
0.792836793301607 0.262833754433967
0.795007111774094 0.268796655430151
0.796794870139636 0.274885201281892
0.798192869726271 0.281074875563959
0.799195481283079 0.287340754642625
0.799798667647188 0.293657608034344
0.8 0.3
};
\addplot [semithick, black]
table {%
0.95 0.85
0.949798667647188 0.856342391965656
0.94919548128308 0.862659245357375
0.948192869726271 0.868925124436041
0.946794870139636 0.875114798718108
0.945007111774095 0.881203344569849
0.942836793301607 0.887166245566033
0.940292653828662 0.892979491208917
0.937384937706979 0.898619673610047
0.934125353283118 0.90406408174556
0.930527025753106 0.909290792905464
0.926604444311898 0.914278760968654
0.922373403810507 0.919007901148211
0.917850941155713 0.923459170865753
0.913055266708452 0.927614646429176
0.90800569095712 0.931457595205034
0.90272254676105 0.934972542994951
0.897227107477268 0.938145336344758
0.891541501300189 0.940963199535452
0.885688622159187 0.943414786026511
0.879692037532827 0.945490224144407
0.873575893550943 0.947181156832354
0.867364817766693 0.948480775301221
0.861083819990101 0.949383846446125
0.854758191582374 0.949886733918301
0.848413403616519 0.949987412767388
0.842075004314321 0.949685477595194
0.835768516172671 0.948982144188093
0.829519333193481 0.947880244621478
0.823352618630997 0.946384215855994
0.817293203668258 0.944500081871467
0.811365487430687 0.942235429410458
0.805593338739423 0.939599377429134
0.8 0.936602540378444
0.794607993613389 0.933256985463477
0.789439031286233 0.929576184053083
0.784513926605471 0.925574957435426
0.779852511229368 0.921269417137886
0.775473555032425 0.916676900051629
0.771394690525721 0.91181589862206
0.767632341857017 0.906705986386277
0.764201658676502 0.901367739157341
0.761116455134508 0.895822652172741
0.758389154256793 0.890093053540661
0.756030737921409 0.884202014332567
0.75405070263855 0.878173255684143
0.752457021311459 0.872031053278654
0.751256111132361 0.865800139597335
0.750452807742692 0.859505604330418
0.750050345761682 0.853172793349807
0.750050345761682 0.846827206650193
0.750452807742692 0.840494395669582
0.751256111132361 0.834199860402665
0.752457021311459 0.827968946721346
0.75405070263855 0.821826744315857
0.756030737921409 0.815797985667433
0.758389154256793 0.809906946459339
0.761116455134508 0.804177347827259
0.764201658676502 0.798632260842659
0.767632341857017 0.793294013613723
0.771394690525721 0.788184101377939
0.775473555032425 0.783323099948371
0.779852511229368 0.778730582862114
0.784513926605471 0.774425042564574
0.789439031286233 0.770423815946917
0.794607993613389 0.766743014536523
0.8 0.763397459621556
0.805593338739423 0.760400622570866
0.811365487430687 0.757764570589542
0.817293203668258 0.755499918128533
0.823352618630997 0.753615784144006
0.829519333193481 0.752119755378522
0.835768516172671 0.751017855811907
0.842075004314321 0.750314522404806
0.848413403616519 0.750012587232612
0.854758191582374 0.750113266081699
0.861083819990101 0.750616153553875
0.867364817766693 0.751519224698779
0.873575893550943 0.752818843167646
0.879692037532827 0.754509775855593
0.885688622159187 0.756585213973489
0.891541501300189 0.759036800464548
0.897227107477268 0.761854663655242
0.90272254676105 0.765027457005049
0.90800569095712 0.768542404794966
0.913055266708452 0.772385353570824
0.917850941155713 0.776540829134247
0.922373403810507 0.780992098851789
0.926604444311898 0.785721239031346
0.930527025753106 0.790709207094536
0.934125353283118 0.79593591825444
0.937384937706979 0.801380326389953
0.940292653828662 0.807020508791083
0.942836793301607 0.812833754433967
0.945007111774095 0.818796655430151
0.946794870139636 0.824885201281892
0.948192869726271 0.831074875563959
0.94919548128308 0.837340754642625
0.949798667647188 0.843657608034344
0.95 0.85
};
\end{axis}

\end{tikzpicture}

%
\begin{tikzpicture}

\definecolor{darkgray176}{RGB}{176,176,176}
\definecolor{darkviolet1910191}{RGB}{191,0,191}
\definecolor{gray}{RGB}{128,128,128}

\begin{axis}[
width=0.75\fwidth,
height=0.75\fwidth,
at={(0\fwidth,0\fwidth)},
tick align=outside,
tick pos=left,
x grid style={darkgray176},
xmin=-0.02, xmax=1.02,
xtick style={color=black},
y grid style={darkgray176},
ymin=-0.02, ymax=1.02,
ytick style={color=black},
legend style={at={(0.03,0.03)}, anchor=south west, legend cell align=left, align=left, draw=white!15!black}
]
\addplot [semithick, darkviolet1910191, mark=*, mark size=2, mark options={solid}, only marks]
table {%
0.996246337890625 0.989500245558773
0.1885986328125 0.112787684804146
0.216583251953125 0.903994987214009
0.926422119140625 0.250588494301343
0.49468994140625 0.498602855255804
0.01495361328125 0.520245897474978
0.58428955078125 0.0354112686074277
0.61822509765625 0.876746430930244
0.9354248046875 0.624650713813951
0.2484130859375 0.374536401971244
0.27093505859375 0.640247929685515
0.4156494140625 0.205151653711324
0.679443359375 0.648529187623838
0.06268310546875 0.736066656505614
0.8076171875 0.436671239140375
0.4246826171875 0.971701468272113
0.0450439453125 0.304831580551745
0.4803466796875 0.704669003708784
0.7752685546875 0.11436264797033
%
%
};
\addlegendentry{$Y_1''$}
\addplot [semithick, red, mark=*, mark size=2.5, mark options={solid}, only marks]
table {%
0.16 0.6
0.7 0.3
0.85 0.85
};
\addplot [line width=0.24pt, black]
table {%
0.5 1
0.476209042088129 0.999433669591504
0.452471978347909 0.997735961286542
0.428842580863357 0.994910720940466
0.405374377819795 0.990964348631353
0.382120532245286 0.985905784161771
0.359133721579285 0.979746486807249
0.336466018341289 0.972500409357334
0.314168772169836 0.964183966508036
0.292292493499057 0.954815997677259
0.270886739136295 0.944417724327462
0.25 0.933012701892219
0.229679591272201 0.92062676641559
0.209971545214401 0.907287976025168
0.190920506889697 0.893026547371394
0.172569633027358 0.877874787177129
0.154960494258944 0.861867019052535
0.138132980947465 0.845039505741056
0.122125212822871 0.827430366972643
0.106973452628606 0.809079493110303
0.0927120239748321 0.790028454785599
0.0793732335844095 0.770320408727799
0.0669872981077808 0.75
0.0555822756725383 0.729113260863705
0.0451840023227408 0.707707506500943
0.0358160334919637 0.685831227830164
0.0274995906426658 0.663533981658711
0.0202535131927513 0.640866278420715
0.0140942158382292 0.617879467754714
0.0090356513686467 0.594625622180205
0.00508927905953366 0.571157419136643
0.00226403871345771 0.547528021652091
0.000566330408496019 0.523790957911871
0 0.5
0.000566330408496019 0.476209042088129
0.00226403871345765 0.452471978347909
0.0050892790595336 0.428842580863358
0.00903565136864665 0.405374377819795
0.0140942158382291 0.382120532245286
0.0202535131927513 0.359133721579285
0.0274995906426658 0.336466018341289
0.0358160334919637 0.314168772169836
0.0451840023227407 0.292292493499057
0.0555822756725383 0.270886739136295
0.0669872981077806 0.25
0.0793732335844094 0.229679591272201
0.0927120239748321 0.209971545214401
0.106973452628606 0.190920506889698
0.122125212822871 0.172569633027358
0.138132980947465 0.154960494258944
0.154960494258944 0.138132980947465
0.172569633027357 0.122125212822871
0.190920506889697 0.106973452628606
0.209971545214401 0.0927120239748322
0.229679591272201 0.0793732335844094
0.25 0.0669872981077808
0.270886739136295 0.0555822756725383
0.292292493499056 0.045184002322741
0.314168772169836 0.0358160334919638
0.336466018341289 0.0274995906426658
0.359133721579285 0.0202535131927513
0.382120532245287 0.0140942158382291
0.405374377819795 0.0090356513686467
0.428842580863357 0.00508927905953366
0.452471978347908 0.00226403871345771
0.476209042088129 0.000566330408496019
0.5 0
0.523790957911871 0.000566330408496019
0.547528021652091 0.00226403871345771
0.571157419136642 0.0050892790595336
0.594625622180205 0.00903565136864659
0.617879467754713 0.0140942158382291
0.640866278420715 0.0202535131927513
0.663533981658711 0.0274995906426657
0.685831227830164 0.0358160334919637
0.707707506500943 0.0451840023227407
0.729113260863705 0.0555822756725382
0.75 0.0669872981077805
0.770320408727799 0.0793732335844093
0.790028454785599 0.0927120239748321
0.809079493110303 0.106973452628606
0.827430366972643 0.122125212822871
0.845039505741056 0.138132980947465
0.861867019052535 0.154960494258944
0.877874787177129 0.172569633027357
0.893026547371394 0.190920506889697
0.907287976025168 0.209971545214401
0.92062676641559 0.229679591272201
0.933012701892219 0.25
0.944417724327462 0.270886739136295
0.954815997677259 0.292292493499057
0.964183966508036 0.314168772169836
0.972500409357334 0.336466018341289
0.979746486807249 0.359133721579285
0.985905784161771 0.382120532245286
0.990964348631353 0.405374377819795
0.994910720940466 0.428842580863357
0.997735961286542 0.452471978347908
0.999433669591504 0.476209042088129
1 0.5
1 1
0.5 1
};
%
%
%
%
%
%
%
%
%
%
%
%
%
%
%
%
%
%
%
%
%
%
%
%
%
%
%
%
%
%
%
%
%
%
%
%
%
%
%
%
%
%
%
%
%
%
%
%
%
%
%
%
%
%
%
%
%
%
%
%
%
%
%
%
%
%
%
%
%
%
%
%
%
%
%
%
%
%
%
%
%
%
%
%
%
%
%
%
%
%
%
%
%
%
%
%
%
%
%
%
%
%
%
%
%
%
%
%
%
%
%
%
%
%
%
%
%
%
%
%
%
%
%
%
%
%
%
%
%
%
%
%
%
%
%
%
%
%
%
%
%
%
%
%
%
%
%
%
%
%
%
%
%
%
%
%
%
%
%
%
%
%
%
%
%
%
%
%
%
%
%
%
%
%
%
%
%
%
%
%
%
%
%
%
%
%
%
%
%
%
%
%
%
%
%
%
%
%
%
%
%
%
%
%
%
%
%
%
%
%
%
%
%
%
%
%
%
%
%
%
%
%
%
%
%
%
%
%
%
%
%
%
%
%
%
%
%
%
%
%
%
%
%
%
%
%
%
%
%
%
%
%
%
%
%
%
%
%
%
%
%
%
%
%
%
%
%
%
%
%
%
%
%
%
%
%
%
%
%
%
%
%
%
%
%
%
%
%
%
%
%
%
%
%
%
%
%
%
%
%
%
%
%
%
%
%
%
%
%
\addplot [semithick, darkviolet1910191, mark=*, mark size=2, mark options={solid}, only marks]
table {%
0.996246337890625 0.989500245558773
0.1885986328125 0.112787684804146
0.216583251953125 0.903994987214009
0.926422119140625 0.250588494301343
0.49468994140625 0.498602855255804
0.01495361328125 0.520245897474978
0.58428955078125 0.0354112686074277
0.61822509765625 0.876746430930244
0.9354248046875 0.624650713813951
0.2484130859375 0.374536401971244
0.27093505859375 0.640247929685515
0.4156494140625 0.205151653711324
0.679443359375 0.648529187623838
0.06268310546875 0.736066656505614
0.8076171875 0.436671239140375
0.4246826171875 0.971701468272113
0.0450439453125 0.304831580551745
0.4803466796875 0.704669003708784
0.7752685546875 0.11436264797033
};
\addplot [semithick, gray, mark=*, mark size=0.375, mark options={solid}, only marks]
table {%
0.996246337890625 0.989500245558773
0.931427001953125 0.993158224525394
0.994781494140625 0.92609527680401
0.965576171875 0.956713915561652
0.96392822265625 0.989229284153838
0.9981689453125 0.959152568206066
0.93719482421875 0.960778336635675
0.9805908203125 0.975003810394757
0.98333740234375 0.941269115480364
0.950103759765625 0.973649003370082
0.9783935546875 0.997764568409287
0.9476318359375 0.995732357872276
0.961090087890625 0.934630561059459
};
\addplot [semithick, gray, mark=*, mark size=0.375, mark options={solid}, only marks]
table {%
0.24957275390625 0.124574505918813
0.18218994140625 0.128232484885434
0.22119140625 0.15546410608139
0.21368408203125 0.0977493268302596
0.1556396484375 0.155057663973988
0.248291015625 0.092059137326627
0.19171142578125 0.157496316618402
0.208740234375 0.128638926992836
0.2314453125 0.111568358481939
0.23583984375 0.140832190214906
0.172027587890625 0.165083235956578
0.170654296875 0.14327084285932
0.230621337890625 0.0882656776575386
0.1885986328125 0.112787684804146
0.206451416015625 0.149231993767888
0.1951904296875 0.136767769140883
0.228057861328125 0.126606716455825
0.157745361328125 0.170502464055276
0.21112060546875 0.112381242696743
0.18603515625 0.173754000914495
};
\addplot [semithick, gray, mark=*, mark size=0.375, mark options={solid}, only marks]
table {%
0.24871826171875 0.875273078290911
0.186920166015625 0.873511829158834
0.216583251953125 0.903994987214009
0.27783203125 0.904435299497028
0.2196044921875 0.846415688665346
0.249176025390625 0.90765296618063
0.2481689453125 0.848041457094955
0.21392822265625 0.878118173042727
0.2305908203125 0.863892699283646
0.20440673828125 0.859828278209622
0.2628173828125 0.861758878219783
0.260162353515625 0.890887229250284
0.237091064453125 0.893021050314146
0.20147705078125 0.892750088909211
0.22747802734375 0.917949499568155
0.16961669921875 0.860641162424427
0.261260986328125 0.926789615404156
0.264923095703125 0.912157699537672
};
\addplot [semithick, gray, mark=*, mark size=0.375, mark options={solid}, only marks]
table {%
0.865264892578125 0.248082101305695
0.931640625 0.314128943758573
0.926422119140625 0.250588494301343
0.8968505859375 0.222273027485648
0.89556884765625 0.275923385662755
0.92926025390625 0.283239343595997
0.8671875 0.275720164609053
0.89410400390625 0.250317532896408
0.8865966796875 0.295432606818066
0.91058349609375 0.291368185744043
0.878997802734375 0.233517925790445
0.939605712890625 0.265220410167827
0.95050048828125 0.296651933140273
0.91790771484375 0.26820098562211
0.90966796875 0.240969364426155
0.87652587890625 0.261291469796271
0.91680908203125 0.307625870040136
0.88238525390625 0.21739572219682
0.864166259765625 0.232908262629342
};
\addplot [semithick, gray, mark=*, mark size=0.375, mark options={solid}, only marks]
table {%
0.49468994140625 0.498602855255804
0.49798583984375 0.560382055580958
0.43463134765625 0.502260834222425
0.499267578125 0.436823654930651
0.560333251953125 0.496587579806601
0.5281982421875 0.52314179749022
0.466217041015625 0.468390658605565
0.464752197265625 0.536672932649156
0.530670166015625 0.466104421751427
0.49615478515625 0.529086013310979
0.465118408203125 0.504293044759437
0.497802734375 0.466087486663618
0.529296875 0.55281207133059
0.437286376953125 0.469609984927772
0.555938720703125 0.523954681705025
0.468414306640625 0.439126826872597
0.52545166015625 0.497535944723873
0.44671630859375 0.525021592236956
0.51611328125 0.450388660265203
0.54266357421875 0.508103439516334
0.48077392578125 0.483970939389321
0.482330322265625 0.51472505884943
0.45184326171875 0.455519991871158
0.54412841796875 0.538993039678911
0.510528564453125 0.485613642906738
0.54522705078125 0.485342681501804
0.45220947265625 0.491286897322563
0.477294921875 0.557536960829142
0.513092041015625 0.513658148317499
0.511993408203125 0.549560534471371
0.48095703125 0.459990855052583
0.489288330078125 0.545750139714474
0.452667236328125 0.548188792358888
0.51446533203125 0.531677081745669
0.48443603515625 0.441294518112076
0.449462890625 0.506325255296449
0.43792724609375 0.484783823604125
0.461456298828125 0.518924960625921
0.51702880859375 0.471523649850124
};
\addplot [semithick, gray, mark=*, mark size=0.375, mark options={solid}, only marks]
table {%
0.0008544921875 0.500736676319667
0.06201171875 0.563328760859625
0.05853271484375 0.502768886856678
0.027496337890625 0.532710122101983
0.032989501953125 0.470524479669427
0.05908203125 0.534064929126657
0.030426025390625 0.561973953834951
0.00396728515625 0.471472844586699
0.0294189453125 0.502362444749276
0.061553955078125 0.477163034090332
0.018341064453125 0.580263848668055
0.041046142578125 0.579586445155718
0.043975830078125 0.488136970990195
0.04644775390625 0.517807244830565
0.0733642578125 0.546258192348727
0.04754638671875 0.551541939744957
0.07666015625 0.519433013260174
0.0118408203125 0.551135497637555
0.02264404296875 0.485698318345781
0.01495361328125 0.520245897474978
0.03253173828125 0.547071076563532
0.0477294921875 0.566580297718844
0.04205322265625 0.532439160697048
};
\addplot [semithick, gray, mark=*, mark size=0.375, mark options={solid}, only marks]
table {%
0.551788330078125 0.0642686582329929
0.62127685546875 0.0671137529848092
0.58428955078125 0.0354112686074277
0.58941650390625 0.0939389320733628
0.523223876953125 0.0356822300123626
0.5537109375 0.0333790580704161
0.613037109375 0.039475689681451
0.59307861328125 0.0634557740181883
0.5712890625 0.0809327846364883
0.599029541015625 0.0198309878236719
0.63226318359375 0.0490270792054057
0.633453369140625 0.0228793036291893
0.569366455078125 0.0532947213331301
0.53887939453125 0.015902047452116
0.5379638671875 0.0512625107961185
0.60919189453125 0.081745668851293
0.5745849609375 0.0163084895595184
0.537139892578125 0.0752425951328558
0.6163330078125 0.0236244474927602
0.60223388671875 0.0508560686887161
0.553619384765625 0.0490948195566394
0.578155517578125 0.0679266371996139
0.626129150390625 0.0814069670951244
};
\addplot [semithick, gray, mark=*, mark size=0.375, mark options={solid}, only marks]
table {%
0.61822509765625 0.876746430930244
0.6810302734375 0.87247878880252
0.616119384765625 0.937983708445528
0.55743408203125 0.875527104608037
0.646148681640625 0.841047265830073
0.59130859375 0.841792409693644
0.653564453125 0.908245694253925
0.58270263671875 0.906823146878016
0.6510009765625 0.870446578265508
0.5869140625 0.871056241426612
0.56072998046875 0.84626327287507
0.61968994140625 0.906010262663212
0.643585205078125 0.936696641772087
0.6195068359375 0.83935375704923
0.6015625 0.893004115226337
0.573638916015625 0.8282443394469
0.6337890625 0.821673525377229
0.606689453125 0.858862978204542
0.634521484375 0.88995579942082
0.598907470703125 0.919016410100086
0.62774658203125 0.859066199258243
0.66400146484375 0.885484936239394
0.57757568359375 0.888533252044912
0.6063232421875 0.819438093786516
0.662353515625 0.85337600975461
0.57867431640625 0.854798557130519
0.65301513671875 0.824518620129045
0.63232421875 0.924096936442615
0.566497802734375 0.900184592457112
0.6162109375 0.922267946959305
0.63507080078125 0.907839252146522
0.65887451171875 0.923284052227811
0.666107177734375 0.869091771240834
0.59185791015625 0.826754051719758
0.635986328125 0.875323883554336
};
\addplot [semithick, gray, mark=*, mark size=0.375, mark options={solid}, only marks]
table {%
0.874267578125 0.622008840115836
0.99420166015625 0.62099273484733
0.93499755859375 0.682771935172484
0.93267822265625 0.561245745059188
0.9354248046875 0.624650713813951
0.966461181640625 0.585903232908263
0.8980712890625 0.589696692577351
0.964263916015625 0.655404833274061
0.90576171875 0.657978966620942
0.929931640625 0.592135345221765
0.875274658203125 0.651069450795102
0.965362548828125 0.621805619062135
0.907958984375 0.616521871665905
0.990631103515625 0.648088875340819
0.9310302734375 0.650256566580298
0.896240234375 0.634811766499009
0.94805908203125 0.603922166336432
0.947509765625 0.576284103033074
0.981842041015625 0.604193127741367
0.952178955078125 0.637656861250826
0.92254638671875 0.668140019306
0.8804931640625 0.604328608443835
0.910430908203125 0.574251892496063
0.888458251953125 0.663940117529509
0.915557861328125 0.632779555961998
0.954498291015625 0.670036749140544
0.969482421875 0.63846974546563
0.925811767578125 0.608393029517858
0.9134521484375 0.595386882080984
0.97900390625 0.661636945587563
0.966552734375 0.601889955799421
0.891754150390625 0.649308201663026
0.9453125 0.654320987654321
0.927734375 0.577503429355281
};
\addplot [semithick, gray, mark=*, mark size=0.375, mark options={solid}, only marks]
table {%
0.2484130859375 0.374536401971244
0.2447509765625 0.43834781283341
0.247833251953125 0.311402394621416
0.18743896484375 0.376568612508256
0.30865478515625 0.380937865162831
0.279541015625 0.413046791647615
0.214630126953125 0.344324205321005
0.2158203125 0.407864654778235
0.27587890625 0.344764517604024
0.2451171875 0.342021033379058
0.250885009765625 0.404376026689698
0.21929931640625 0.373317075649037
0.275146484375 0.383782959914647
0.218994140625 0.316415180612711
0.29095458984375 0.328100391200528
0.286407470703125 0.363460854544531
0.2618408203125 0.328913275415333
0.229644775390625 0.391877931887077
0.2352294921875 0.418432149570696
0.200592041015625 0.328472963132314
0.20172119140625 0.391606970482142
0.200347900390625 0.358956121187488
0.257476806640625 0.422665921522803
0.261962890625 0.358177107148301
0.23687744140625 0.358685159782553
0.28961181640625 0.398008433673729
0.26434326171875 0.439059086521364
0.25830078125 0.388660265203475
0.201446533203125 0.4192450337855
0.232177734375 0.326169791190367
0.29290771484375 0.379312096733222
0.291168212890625 0.343409710579349
0.28192138671875 0.428085149621501
0.301422119140625 0.361699605412454
0.18798828125 0.396890717878372
0.228546142578125 0.43143829700757
};
\addplot [semithick, gray, mark=*, mark size=0.375, mark options={solid}, only marks]
table {%
0.249053955078125 0.62531118223848
0.30657958984375 0.690240308896002
0.309844970703125 0.626835340141239
0.24578857421875 0.685058172026622
0.28045654296875 0.65853782451862
0.277435302734375 0.59337160663178
0.21966552734375 0.654575013971447
0.3128662109375 0.656912056089011
0.279815673828125 0.623177361174618
0.24493408203125 0.591576487324087
0.220489501953125 0.618672627817575
0.2774658203125 0.691053193110806
0.25189208984375 0.657318498196413
0.3092041015625 0.594726413656455
0.231475830078125 0.636285119138343
0.295196533203125 0.641467256007722
0.327423095703125 0.615861403241376
0.260040283203125 0.608545445308134
0.23187255859375 0.670426256160138
0.294952392578125 0.671950414062897
0.259033203125 0.577655845145557
0.27093505859375 0.640247929685515
0.32415771484375 0.675608393029518
0.23724365234375 0.610679266371996
0.29547119140625 0.613829192704364
0.266845703125 0.676421277244322
0.3223876953125 0.642280140222527
0.222686767578125 0.596047350505512
0.20831298828125 0.629375603312503
0.26300048828125 0.704059340547681
0.294921875 0.589849108367627
0.2476806640625 0.643194634964182
};
\addplot [semithick, gray, mark=*, mark size=0.375, mark options={solid}, only marks]
table {%
0.368682861328125 0.250063506579282
0.431304931640625 0.186726278175752
0.432037353515625 0.248911920608308
0.370147705078125 0.189097190468933
0.404205322265625 0.218428762553134
0.3995361328125 0.153939948178631
0.462982177734375 0.214770783586513
0.4013671875 0.25148605395519
0.4332275390625 0.157597927145252
0.43304443359375 0.217751359040797
0.40264892578125 0.187674643093024
0.46051025390625 0.246202306558959
0.375 0.222222222222222
0.46087646484375 0.187268200985622
0.44622802734375 0.201900116852106
0.41473388671875 0.236041253873901
0.389923095703125 0.171416958796931
0.3511962890625 0.206980643194635
0.38818359375 0.236854138088706
0.44970703125 0.262459990855053
0.41455078125 0.141746684956562
0.41748046875 0.269775948788294
0.38909912109375 0.206777422140934
0.4808349609375 0.201493674744704
0.41107177734375 0.168978306152517
0.35723876953125 0.228522074886958
0.44915771484375 0.231163948585073
0.4156494140625 0.205151653711324
0.45172119140625 0.16938474825992
0.401275634765625 0.267201815441413
0.417755126953125 0.253789225897136
0.42938232421875 0.23238327490728
0.386260986328125 0.186048874663415
0.4317626953125 0.26367931717726
0.479736328125 0.221002895900015
0.417022705078125 0.183610222019001
0.434967041015625 0.172094362309269
};
\addplot [semithick, gray, mark=*, mark size=0.375, mark options={solid}, only marks]
table {%
0.680206298828125 0.679418787786415
0.68621826171875 0.616013819031652
0.647857666015625 0.647174380599163
0.718994140625 0.649748513946045
0.679443359375 0.648529187623838
0.688507080078125 0.585530660976477
0.70977783203125 0.679825229893817
0.6832275390625 0.713153482700808
0.65142822265625 0.618858913783468
0.647613525390625 0.681315517620959
0.712677001953125 0.622787854155024
0.61956787109375 0.64710664024793
0.74371337890625 0.648122745516435
0.6680908203125 0.604633440024387
0.69976806640625 0.640806787583194
0.664794921875 0.663161103490322
0.700592041015625 0.661806296465647
0.63385009765625 0.629832850683331
0.70086669921875 0.601381903165168
0.725616455078125 0.695270029975105
0.671661376953125 0.630103812088266
0.64190673828125 0.600569018950363
0.6322021484375 0.663973987705126
0.66455078125 0.697302240512117
0.69659423828125 0.699334451049129
0.727813720703125 0.672102829853173
0.734649658203125 0.631865061220342
0.668365478515625 0.584988738166607
0.69439697265625 0.684702535182645
};
\addplot [semithick, gray, mark=*, mark size=0.375, mark options={solid}, only marks]
table {%
0.12420654296875 0.749072803942488
0.062042236328125 0.688106487832139
0.094635009765625 0.78023336551
0.091064453125 0.72306050906874
0.09454345703125 0.684854950972921
0.08740234375 0.749885688157293
0.0650634765625 0.718183203779912
0.11761474609375 0.723466951176142
0.105987548828125 0.757608088197937
0.10882568359375 0.70395773002083
0.0511474609375 0.670629477213839
0.078887939453125 0.695422445765381
0.0797119140625 0.67428745618046
0.05224609375 0.705989940557842
0.10552978515625 0.734847330183407
0.0775146484375 0.731189351216786
0.09344482421875 0.701519077376416
0.092071533203125 0.764924046131179
0.06268310546875 0.736066656505614
};
\addplot [semithick, gray, mark=*, mark size=0.375, mark options={solid}, only marks]
table {%
0.8076171875 0.436671239140375
0.801300048828125 0.374892038815221
0.744903564453125 0.436230926857356
0.8094482421875 0.498450439465529
0.86676025390625 0.431387491744145
0.840484619140625 0.464580263848668
0.77008056640625 0.402123660011177
0.775390625 0.462277091906722
0.835723876953125 0.406052600382733
0.80303955078125 0.405781638977798
0.778228759765625 0.430439126826873
0.80615234375 0.465935070873342
0.781982421875 0.490321597317482
0.8350830078125 0.492760249961896
0.83819580078125 0.438703449677387
0.75067138671875 0.465122186658538
0.8218994140625 0.39033683889651
0.79022216796875 0.448458060255042
0.82501220703125 0.422852207488696
0.76220703125 0.447645176040238
0.82373046875 0.45496113397348
0.82318115234375 0.474876797236194
0.76312255859375 0.479754102525022
0.8504638671875 0.421632881166489
0.85009765625 0.480566986739826
0.78399658203125 0.389930396789107
0.854400634765625 0.452387000626598
0.79443359375 0.422039323273891
0.762847900390625 0.420684516249217
0.80267333984375 0.482192755169436
0.810089111328125 0.418787786414673
0.781890869140625 0.475554200748531
};
\addplot [semithick, gray, mark=*, mark size=0.375, mark options={solid}, only marks]
table {%
0.432373046875 0.941624752324341
0.371490478515625 0.934782976849735
0.46240234375 0.972107910379515
0.394561767578125 0.966417720875883
0.4613037109375 0.936747447035513
0.40032958984375 0.937153889142915
0.4246826171875 0.971701468272113
0.437042236328125 0.910328710054362
0.4678955078125 0.998933089468069
0.36444091796875 0.971091805111009
0.44464111328125 0.983894731494183
0.44482421875 0.961133973479652
0.41131591796875 0.954630899761215
0.376617431640625 0.95490186116615
0.38604736328125 0.918051110095006
0.473388671875 0.952598689224204
0.453887939453125 0.917644667987603
0.418365478515625 0.918322071499941
0.381011962890625 0.985385019221324
0.41259765625 0.986739826245999
0.480987548828125 0.979830310420159
0.44757080078125 0.944876289183559
0.440704345703125 0.928618604887466
0.47613525390625 0.937560331250317
0.4180908203125 0.93796677335772
0.4591064453125 0.955037341868617
};
\addplot [semithick, gray, mark=*, mark size=0.375, mark options={solid}, only marks]
table {%
0.0640869140625 0.312147538484987
0.0926513671875 0.277599959355789
0.0927734375 0.347508001828989
0.036651611328125 0.34059848600315
0.058868408203125 0.282070822537215
0.093017578125 0.313366864807194
0.063018798828125 0.345475791291978
0.104736328125 0.295076969974089
0.073883056640625 0.293722162949415
0.04681396484375 0.360107707158462
0.0450439453125 0.304831580551745
0.072784423828125 0.329624549103287
0.06903076171875 0.266219580348524
0.0494384765625 0.32312147538485
0.0220947265625 0.356043286084438
0.09136962890625 0.331250317532896
0.09027099609375 0.295483412081492
};
\addplot [semithick, gray, mark=*, mark size=0.375, mark options={solid}, only marks]
table {%
0.4957275390625 0.750190519737845
0.436676025390625 0.685430743958407
0.43255615234375 0.750596961845247
0.49462890625 0.690443529949703
0.460113525390625 0.718352554657996
0.5308837890625 0.718640451150739
0.469970703125 0.659960371894528
0.503570556640625 0.659977306982337
0.525390625 0.684499314128944
0.5264892578125 0.747904282883707
0.49169921875 0.71970736168267
0.4620361328125 0.746532540771224
0.43994140625 0.657521719250114
0.433380126953125 0.714694575691375
0.463043212890625 0.689088722925028
0.485626220703125 0.643296245491033
0.51312255859375 0.701976324747244
0.509918212890625 0.734356212636962
0.445465087890625 0.73298447052448
0.53839111328125 0.699944114210232
0.47705078125 0.734339277549154
0.47357177734375 0.76157089874511
0.4803466796875 0.704669003708784
0.482818603515625 0.676895459702959
0.416534423828125 0.705481887923589
0.51861572265625 0.668241629832851
0.449859619140625 0.674456807058545
0.453521728515625 0.702501312469305
0.508087158203125 0.681247777269725
0.4554443359375 0.761164456637708
0.517120361328125 0.651983945536757
0.43035888671875 0.729868414367728
0.50091552734375 0.642229334959102
0.50946044921875 0.716608240613728
0.49609375 0.705075445816187
0.494903564453125 0.769564260190689
0.48919677734375 0.662805466646345
};
\addplot [semithick, gray, mark=*, mark size=0.375, mark options={solid}, only marks]
table {%
0.741485595703125 0.134379921759894
0.802734375 0.133058984910837
0.740478515625 0.0632525529644871
0.766021728515625 0.0975630408643669
0.71630859375 0.101051668952904
0.7686767578125 0.163542142966011
0.80035400390625 0.10216938474826
0.772247314453125 0.131704177886162
0.800811767578125 0.163948585073414
0.71209716796875 0.127470405934055
0.828277587890625 0.128046198919541
0.7364501953125 0.160798658741046
0.74102783203125 0.0872326373012244
0.7430419921875 0.111619163745364
0.818389892578125 0.14931666920693
0.789459228515625 0.14877474639706
0.723541259765625 0.142373283205473
0.7496337890625 0.175430574607529
0.75384521484375 0.152974648173551
0.78350830078125 0.175328964080679
0.81573486328125 0.115581974292537
0.7752685546875 0.11436264797033
0.755126953125 0.0769699740893156
0.729034423828125 0.0758522582939593
0.7276611328125 0.116902911141594
0.765380859375 0.179393385154702
0.7584228515625 0.115988416399939
0.796966552734375 0.118291588341886
0.731597900390625 0.0996968619282291
0.7701416015625 0.0830666057003506
0.774078369140625 0.146336093752646
};
\end{axis}

\end{tikzpicture}
		%

%
\begin{tikzpicture}

\definecolor{darkgray176}{RGB}{176,176,176}
\definecolor{darkviolet1910191}{RGB}{191,0,191}
\definecolor{gray}{RGB}{128,128,128}

\begin{axis}[
width=0.75\fwidth,
height=0.75\fwidth,
at={(0\fwidth,0\fwidth)},
tick align=outside,
tick pos=left,
x grid style={darkgray176},
xmin=-0.02, xmax=1.02,
xtick style={color=black},
y grid style={darkgray176},
ymin=-0.02, ymax=1.02,
ytick style={color=black}
]
\addplot [line width=0.24pt, black]
table {%
0.5 1
0.476209042088129 0.999433669591504
0.452471978347909 0.997735961286542
0.428842580863357 0.994910720940466
0.405374377819795 0.990964348631353
0.382120532245286 0.985905784161771
0.359133721579285 0.979746486807249
0.336466018341289 0.972500409357334
0.314168772169836 0.964183966508036
0.292292493499057 0.954815997677259
0.270886739136295 0.944417724327462
0.25 0.933012701892219
0.229679591272201 0.92062676641559
0.209971545214401 0.907287976025168
0.190920506889697 0.893026547371394
0.172569633027358 0.877874787177129
0.154960494258944 0.861867019052535
0.138132980947465 0.845039505741056
0.122125212822871 0.827430366972643
0.106973452628606 0.809079493110303
0.0927120239748321 0.790028454785599
0.0793732335844095 0.770320408727799
0.0669872981077808 0.75
0.0555822756725383 0.729113260863705
0.0451840023227408 0.707707506500943
0.0358160334919637 0.685831227830164
0.0274995906426658 0.663533981658711
0.0202535131927513 0.640866278420715
0.0140942158382292 0.617879467754714
0.0090356513686467 0.594625622180205
0.00508927905953366 0.571157419136643
0.00226403871345771 0.547528021652091
0.000566330408496019 0.523790957911871
0 0.5
0.000566330408496019 0.476209042088129
0.00226403871345765 0.452471978347909
0.0050892790595336 0.428842580863358
0.00903565136864665 0.405374377819795
0.0140942158382291 0.382120532245286
0.0202535131927513 0.359133721579285
0.0274995906426658 0.336466018341289
0.0358160334919637 0.314168772169836
0.0451840023227407 0.292292493499057
0.0555822756725383 0.270886739136295
0.0669872981077806 0.25
0.0793732335844094 0.229679591272201
0.0927120239748321 0.209971545214401
0.106973452628606 0.190920506889698
0.122125212822871 0.172569633027358
0.138132980947465 0.154960494258944
0.154960494258944 0.138132980947465
0.172569633027357 0.122125212822871
0.190920506889697 0.106973452628606
0.209971545214401 0.0927120239748322
0.229679591272201 0.0793732335844094
0.25 0.0669872981077808
0.270886739136295 0.0555822756725383
0.292292493499056 0.045184002322741
0.314168772169836 0.0358160334919638
0.336466018341289 0.0274995906426658
0.359133721579285 0.0202535131927513
0.382120532245287 0.0140942158382291
0.405374377819795 0.0090356513686467
0.428842580863357 0.00508927905953366
0.452471978347908 0.00226403871345771
0.476209042088129 0.000566330408496019
0.5 0
0.523790957911871 0.000566330408496019
0.547528021652091 0.00226403871345771
0.571157419136642 0.0050892790595336
0.594625622180205 0.00903565136864659
0.617879467754713 0.0140942158382291
0.640866278420715 0.0202535131927513
0.663533981658711 0.0274995906426657
0.685831227830164 0.0358160334919637
0.707707506500943 0.0451840023227407
0.729113260863705 0.0555822756725382
0.75 0.0669872981077805
0.770320408727799 0.0793732335844093
0.790028454785599 0.0927120239748321
0.809079493110303 0.106973452628606
0.827430366972643 0.122125212822871
0.845039505741056 0.138132980947465
0.861867019052535 0.154960494258944
0.877874787177129 0.172569633027357
0.893026547371394 0.190920506889697
0.907287976025168 0.209971545214401
0.92062676641559 0.229679591272201
0.933012701892219 0.25
0.944417724327462 0.270886739136295
0.954815997677259 0.292292493499057
0.964183966508036 0.314168772169836
0.972500409357334 0.336466018341289
0.979746486807249 0.359133721579285
0.985905784161771 0.382120532245286
0.990964348631353 0.405374377819795
0.994910720940466 0.428842580863357
0.997735961286542 0.452471978347908
0.999433669591504 0.476209042088129
1 0.5
1 1
0.5 1
};
\addplot [semithick, red, mark=*, mark size=2.5, mark options={solid}, only marks]
table {%
0.16 0.6
0.7 0.3
0.85 0.85
};
%
%
%
%
%
%
%
%
%
%
%
%
%
%
%
%
%
%
%
%
%
%
%
%
%
%
%
%
%
%
%
%
%
%
%
%
%
%
%
%
%
%
%
%
%
%
%
%
%
%
%
%
%
%
%
%
%
%
%
%
%
%
%
%
%
%
%
%
%
%
%
%
%
%
%
%
%
%
%
%
%
%
%
%
%
%
%
%
%
%
%
%
%
%
%
%
%
%
%
%
%
%
%
%
%
%
%
%
%
%
%
%
%
%
%
%
%
%
%
%
%
%
%
%
%
%
%
%
%
%
%
%
%
%
%
%
%
%
%
%
%
%
%
%
%
%
%
%
%
%
%
%
%
%
%
%
%
%
%
%
%
%
%
%
%
%
%
%
%
%
%
%
%
%
%
%
%
%
%
%
%
%
%
%
%
%
%
%
%
%
%
%
%
%
%
%
%
%
%
%
%
%
%
%
%
%
%
%
%
%
%
%
%
%
%
%
%
%
%
%
%
%
%
%
%
%
%
%
%
%
%
%
%
%
%
%
%
%
%
%
%
%
%
%
%
%
%
%
%
%
%
%
%
%
%
%
%
%
%
%
%
%
%
%
%
%
%
%
%
%
%
%
%
%
%
%
%
%
%
%
%
%
%
%
%
%
%
%
%
%
%
%
%
%
%
%
%
%
%
%
%
%
%
%
%
%
%
%
%
\addplot [semithick, darkviolet1910191, mark=*, mark size=2, mark options={solid}, only marks]
table {%
0.996246337890625 0.989500245558773
0.1885986328125 0.112787684804146
0.216583251953125 0.903994987214009
0.926422119140625 0.250588494301343
0.49468994140625 0.498602855255804
0.01495361328125 0.520245897474978
0.58428955078125 0.0354112686074277
0.61822509765625 0.876746430930244
0.9354248046875 0.624650713813951
0.2484130859375 0.374536401971244
0.27093505859375 0.640247929685515
0.4156494140625 0.205151653711324
0.679443359375 0.648529187623838
0.06268310546875 0.736066656505614
0.8076171875 0.436671239140375
0.4246826171875 0.971701468272113
0.0450439453125 0.304831580551745
0.4803466796875 0.704669003708784
0.7752685546875 0.11436264797033
};
\addplot [semithick, gray, mark=*, mark size=0.5, mark options={solid}, only marks]
table {%
0.996246337890625 0.989500245558773
0.931427001953125 0.993158224525394
0.994781494140625 0.92609527680401
0.965576171875 0.956713915561652
0.96392822265625 0.989229284153838
0.9981689453125 0.959152568206066
0.93719482421875 0.960778336635675
0.9805908203125 0.975003810394757
0.98333740234375 0.941269115480364
0.950103759765625 0.973649003370082
0.9783935546875 0.997764568409287
0.9476318359375 0.995732357872276
0.961090087890625 0.934630561059459
};
\addplot [semithick, gray, mark=*, mark size=0.5, mark options={solid}, only marks]
table {%
0.24957275390625 0.124574505918813
0.18218994140625 0.128232484885434
0.22119140625 0.15546410608139
0.21368408203125 0.0977493268302596
0.1556396484375 0.155057663973988
0.248291015625 0.092059137326627
0.19171142578125 0.157496316618402
0.208740234375 0.128638926992836
0.2314453125 0.111568358481939
0.23583984375 0.140832190214906
0.172027587890625 0.165083235956578
0.170654296875 0.14327084285932
0.230621337890625 0.0882656776575386
};
\addplot [semithick, gray, mark=*, mark size=0.5, mark options={solid}, only marks]
table {%
0.24871826171875 0.875273078290911
0.186920166015625 0.873511829158834
0.216583251953125 0.903994987214009
0.27783203125 0.904435299497028
0.2196044921875 0.846415688665346
0.249176025390625 0.90765296618063
0.2481689453125 0.848041457094955
0.21392822265625 0.878118173042727
0.2305908203125 0.863892699283646
0.20440673828125 0.859828278209622
0.2628173828125 0.861758878219783
0.260162353515625 0.890887229250284
0.237091064453125 0.893021050314146
};
\addplot [semithick, gray, mark=*, mark size=0.5, mark options={solid}, only marks]
table {%
0.865264892578125 0.248082101305695
0.931640625 0.314128943758573
0.926422119140625 0.250588494301343
0.8968505859375 0.222273027485648
0.89556884765625 0.275923385662755
0.92926025390625 0.283239343595997
0.8671875 0.275720164609053
0.89410400390625 0.250317532896408
0.8865966796875 0.295432606818066
0.91058349609375 0.291368185744043
0.878997802734375 0.233517925790445
0.939605712890625 0.265220410167827
0.95050048828125 0.296651933140273
};
\addplot [semithick, gray, mark=*, mark size=0.5, mark options={solid}, only marks]
table {%
0.49468994140625 0.498602855255804
0.49798583984375 0.560382055580958
0.43463134765625 0.502260834222425
0.499267578125 0.436823654930651
0.560333251953125 0.496587579806601
0.5281982421875 0.52314179749022
0.466217041015625 0.468390658605565
0.464752197265625 0.536672932649156
0.530670166015625 0.466104421751427
0.49615478515625 0.529086013310979
0.465118408203125 0.504293044759437
0.497802734375 0.466087486663618
0.529296875 0.55281207133059
};
\addplot [semithick, gray, mark=*, mark size=0.5, mark options={solid}, only marks]
table {%
0.0008544921875 0.500736676319667
0.06201171875 0.563328760859625
0.05853271484375 0.502768886856678
0.027496337890625 0.532710122101983
0.032989501953125 0.470524479669427
0.05908203125 0.534064929126657
0.030426025390625 0.561973953834951
0.00396728515625 0.471472844586699
0.0294189453125 0.502362444749276
0.061553955078125 0.477163034090332
0.018341064453125 0.580263848668055
0.041046142578125 0.579586445155718
0.043975830078125 0.488136970990195
};
\addplot [semithick, gray, mark=*, mark size=0.5, mark options={solid}, only marks]
table {%
0.551788330078125 0.0642686582329929
0.62127685546875 0.0671137529848092
0.58428955078125 0.0354112686074277
0.58941650390625 0.0939389320733628
0.523223876953125 0.0356822300123626
0.5537109375 0.0333790580704161
0.613037109375 0.039475689681451
0.59307861328125 0.0634557740181883
0.5712890625 0.0809327846364883
0.599029541015625 0.0198309878236719
0.63226318359375 0.0490270792054057
0.633453369140625 0.0228793036291893
0.569366455078125 0.0532947213331301
};
\addplot [semithick, gray, mark=*, mark size=0.5, mark options={solid}, only marks]
table {%
0.61822509765625 0.876746430930244
0.6810302734375 0.87247878880252
0.616119384765625 0.937983708445528
0.55743408203125 0.875527104608037
0.646148681640625 0.841047265830073
0.59130859375 0.841792409693644
0.653564453125 0.908245694253925
0.58270263671875 0.906823146878016
0.6510009765625 0.870446578265508
0.5869140625 0.871056241426612
0.56072998046875 0.84626327287507
0.61968994140625 0.906010262663212
0.643585205078125 0.936696641772087
};
\addplot [semithick, gray, mark=*, mark size=0.5, mark options={solid}, only marks]
table {%
0.874267578125 0.622008840115836
0.99420166015625 0.62099273484733
0.93499755859375 0.682771935172484
0.93267822265625 0.561245745059188
0.9354248046875 0.624650713813951
0.966461181640625 0.585903232908263
0.8980712890625 0.589696692577351
0.964263916015625 0.655404833274061
0.90576171875 0.657978966620942
0.929931640625 0.592135345221765
0.875274658203125 0.651069450795102
0.965362548828125 0.621805619062135
0.907958984375 0.616521871665905
};
\addplot [semithick, gray, mark=*, mark size=0.5, mark options={solid}, only marks]
table {%
0.2484130859375 0.374536401971244
0.2447509765625 0.43834781283341
0.247833251953125 0.311402394621416
0.18743896484375 0.376568612508256
0.30865478515625 0.380937865162831
0.279541015625 0.413046791647615
0.214630126953125 0.344324205321005
0.2158203125 0.407864654778235
0.27587890625 0.344764517604024
0.2451171875 0.342021033379058
0.250885009765625 0.404376026689698
0.21929931640625 0.373317075649037
0.275146484375 0.383782959914647
};
\addplot [semithick, gray, mark=*, mark size=0.5, mark options={solid}, only marks]
table {%
0.249053955078125 0.62531118223848
0.30657958984375 0.690240308896002
0.309844970703125 0.626835340141239
0.24578857421875 0.685058172026622
0.28045654296875 0.65853782451862
0.277435302734375 0.59337160663178
0.21966552734375 0.654575013971447
0.3128662109375 0.656912056089011
0.279815673828125 0.623177361174618
0.24493408203125 0.591576487324087
0.220489501953125 0.618672627817575
0.2774658203125 0.691053193110806
0.25189208984375 0.657318498196413
};
\addplot [semithick, gray, mark=*, mark size=0.5, mark options={solid}, only marks]
table {%
0.368682861328125 0.250063506579282
0.431304931640625 0.186726278175752
0.432037353515625 0.248911920608308
0.370147705078125 0.189097190468933
0.404205322265625 0.218428762553134
0.3995361328125 0.153939948178631
0.462982177734375 0.214770783586513
0.4013671875 0.25148605395519
0.4332275390625 0.157597927145252
0.43304443359375 0.217751359040797
0.40264892578125 0.187674643093024
0.46051025390625 0.246202306558959
0.375 0.222222222222222
};
\addplot [semithick, gray, mark=*, mark size=0.5, mark options={solid}, only marks]
table {%
0.680206298828125 0.679418787786415
0.68621826171875 0.616013819031652
0.647857666015625 0.647174380599163
0.718994140625 0.649748513946045
0.679443359375 0.648529187623838
0.688507080078125 0.585530660976477
0.70977783203125 0.679825229893817
0.6832275390625 0.713153482700808
0.65142822265625 0.618858913783468
0.647613525390625 0.681315517620959
0.712677001953125 0.622787854155024
0.61956787109375 0.64710664024793
0.74371337890625 0.648122745516435
};
\addplot [semithick, gray, mark=*, mark size=0.5, mark options={solid}, only marks]
table {%
0.12420654296875 0.749072803942488
0.062042236328125 0.688106487832139
0.094635009765625 0.78023336551
0.091064453125 0.72306050906874
0.09454345703125 0.684854950972921
0.08740234375 0.749885688157293
0.0650634765625 0.718183203779912
0.11761474609375 0.723466951176142
0.105987548828125 0.757608088197937
0.10882568359375 0.70395773002083
0.0511474609375 0.670629477213839
0.078887939453125 0.695422445765381
0.0797119140625 0.67428745618046
};
\addplot [semithick, gray, mark=*, mark size=0.5, mark options={solid}, only marks]
table {%
0.8076171875 0.436671239140375
0.801300048828125 0.374892038815221
0.744903564453125 0.436230926857356
0.8094482421875 0.498450439465529
0.86676025390625 0.431387491744145
0.840484619140625 0.464580263848668
0.77008056640625 0.402123660011177
0.775390625 0.462277091906722
0.835723876953125 0.406052600382733
0.80303955078125 0.405781638977798
0.778228759765625 0.430439126826873
0.80615234375 0.465935070873342
0.781982421875 0.490321597317482
};
\addplot [semithick, gray, mark=*, mark size=0.5, mark options={solid}, only marks]
table {%
0.432373046875 0.941624752324341
0.371490478515625 0.934782976849735
0.46240234375 0.972107910379515
0.394561767578125 0.966417720875883
0.4613037109375 0.936747447035513
0.40032958984375 0.937153889142915
0.4246826171875 0.971701468272113
0.437042236328125 0.910328710054362
0.4678955078125 0.998933089468069
0.36444091796875 0.971091805111009
0.44464111328125 0.983894731494183
0.44482421875 0.961133973479652
0.41131591796875 0.954630899761215
};
\addplot [semithick, gray, mark=*, mark size=0.5, mark options={solid}, only marks]
table {%
0.0640869140625 0.312147538484987
0.0926513671875 0.277599959355789
0.0927734375 0.347508001828989
0.036651611328125 0.34059848600315
0.058868408203125 0.282070822537215
0.093017578125 0.313366864807194
0.063018798828125 0.345475791291978
0.104736328125 0.295076969974089
0.073883056640625 0.293722162949415
0.04681396484375 0.360107707158462
0.0450439453125 0.304831580551745
0.072784423828125 0.329624549103287
0.06903076171875 0.266219580348524
};
\addplot [semithick, gray, mark=*, mark size=0.5, mark options={solid}, only marks]
table {%
0.4957275390625 0.750190519737845
0.436676025390625 0.685430743958407
0.43255615234375 0.750596961845247
0.49462890625 0.690443529949703
0.460113525390625 0.718352554657996
0.5308837890625 0.718640451150739
0.469970703125 0.659960371894528
0.503570556640625 0.659977306982337
0.525390625 0.684499314128944
0.5264892578125 0.747904282883707
0.49169921875 0.71970736168267
0.4620361328125 0.746532540771224
0.43994140625 0.657521719250114
};
\addplot [semithick, gray, mark=*, mark size=0.5, mark options={solid}, only marks]
table {%
0.741485595703125 0.134379921759894
0.802734375 0.133058984910837
0.740478515625 0.0632525529644871
0.766021728515625 0.0975630408643669
0.71630859375 0.101051668952904
0.7686767578125 0.163542142966011
0.80035400390625 0.10216938474826
0.772247314453125 0.131704177886162
0.800811767578125 0.163948585073414
0.71209716796875 0.127470405934055
0.828277587890625 0.128046198919541
0.7364501953125 0.160798658741046
0.74102783203125 0.0872326373012244
};
\end{axis}

\end{tikzpicture}

%
\begin{tikzpicture}

\definecolor{darkgray176}{RGB}{176,176,176}
\definecolor{darkviolet1910191}{RGB}{191,0,191}
\definecolor{gray}{RGB}{128,128,128}

\begin{axis}[
width=0.75\fwidth,
height=0.75\fwidth,
at={(0\fwidth,0\fwidth)},
tick align=outside,
tick pos=left,
x grid style={darkgray176},
xmin=-0.02, xmax=1.02,
xtick style={color=black},
y grid style={darkgray176},
ymin=-0.02, ymax=1.02,
ytick style={color=black},
legend style={at={(0.03,0.03)}, anchor=south west, legend cell align=left, align=left, draw=white!15!black}
]
\addplot [semithick, red, mark=x, mark size=2.5, mark options={solid}, only marks]
table {%
0.961090087890625 0.934630561059459
0.230621337890625 0.0882656776575386
0.237091064453125 0.893021050314146
0.95050048828125 0.296651933140273
0.529296875 0.55281207133059
0.043975830078125 0.488136970990195
0.569366455078125 0.0532947213331301
0.643585205078125 0.936696641772087
0.907958984375 0.616521871665905
0.275146484375 0.383782959914647
0.25189208984375 0.657318498196413
0.375 0.222222222222222
0.74371337890625 0.648122745516435
0.0797119140625 0.67428745618046
0.781982421875 0.490321597317482
0.41131591796875 0.954630899761215
0.06903076171875 0.266219580348524
0.43994140625 0.657521719250114
0.74102783203125 0.0872326373012244
};
\addlegendentry{$Y_1$}
\addplot [line width=0.24pt, black]
table {%
0.5 1
0.476209042088129 0.999433669591504
0.452471978347909 0.997735961286542
0.428842580863357 0.994910720940466
0.405374377819795 0.990964348631353
0.382120532245286 0.985905784161771
0.359133721579285 0.979746486807249
0.336466018341289 0.972500409357334
0.314168772169836 0.964183966508036
0.292292493499057 0.954815997677259
0.270886739136295 0.944417724327462
0.25 0.933012701892219
0.229679591272201 0.92062676641559
0.209971545214401 0.907287976025168
0.190920506889697 0.893026547371394
0.172569633027358 0.877874787177129
0.154960494258944 0.861867019052535
0.138132980947465 0.845039505741056
0.122125212822871 0.827430366972643
0.106973452628606 0.809079493110303
0.0927120239748321 0.790028454785599
0.0793732335844095 0.770320408727799
0.0669872981077808 0.75
0.0555822756725383 0.729113260863705
0.0451840023227408 0.707707506500943
0.0358160334919637 0.685831227830164
0.0274995906426658 0.663533981658711
0.0202535131927513 0.640866278420715
0.0140942158382292 0.617879467754714
0.0090356513686467 0.594625622180205
0.00508927905953366 0.571157419136643
0.00226403871345771 0.547528021652091
0.000566330408496019 0.523790957911871
0 0.5
0.000566330408496019 0.476209042088129
0.00226403871345765 0.452471978347909
0.0050892790595336 0.428842580863358
0.00903565136864665 0.405374377819795
0.0140942158382291 0.382120532245286
0.0202535131927513 0.359133721579285
0.0274995906426658 0.336466018341289
0.0358160334919637 0.314168772169836
0.0451840023227407 0.292292493499057
0.0555822756725383 0.270886739136295
0.0669872981077806 0.25
0.0793732335844094 0.229679591272201
0.0927120239748321 0.209971545214401
0.106973452628606 0.190920506889698
0.122125212822871 0.172569633027358
0.138132980947465 0.154960494258944
0.154960494258944 0.138132980947465
0.172569633027357 0.122125212822871
0.190920506889697 0.106973452628606
0.209971545214401 0.0927120239748322
0.229679591272201 0.0793732335844094
0.25 0.0669872981077808
0.270886739136295 0.0555822756725383
0.292292493499056 0.045184002322741
0.314168772169836 0.0358160334919638
0.336466018341289 0.0274995906426658
0.359133721579285 0.0202535131927513
0.382120532245287 0.0140942158382291
0.405374377819795 0.0090356513686467
0.428842580863357 0.00508927905953366
0.452471978347908 0.00226403871345771
0.476209042088129 0.000566330408496019
0.5 0
0.523790957911871 0.000566330408496019
0.547528021652091 0.00226403871345771
0.571157419136642 0.0050892790595336
0.594625622180205 0.00903565136864659
0.617879467754713 0.0140942158382291
0.640866278420715 0.0202535131927513
0.663533981658711 0.0274995906426657
0.685831227830164 0.0358160334919637
0.707707506500943 0.0451840023227407
0.729113260863705 0.0555822756725382
0.75 0.0669872981077805
0.770320408727799 0.0793732335844093
0.790028454785599 0.0927120239748321
0.809079493110303 0.106973452628606
0.827430366972643 0.122125212822871
0.845039505741056 0.138132980947465
0.861867019052535 0.154960494258944
0.877874787177129 0.172569633027357
0.893026547371394 0.190920506889697
0.907287976025168 0.209971545214401
0.92062676641559 0.229679591272201
0.933012701892219 0.25
0.944417724327462 0.270886739136295
0.954815997677259 0.292292493499057
0.964183966508036 0.314168772169836
0.972500409357334 0.336466018341289
0.979746486807249 0.359133721579285
0.985905784161771 0.382120532245286
0.990964348631353 0.405374377819795
0.994910720940466 0.428842580863357
0.997735961286542 0.452471978347908
0.999433669591504 0.476209042088129
1 0.5
1 1
0.5 1
};
\addplot [semithick, red, mark=*, mark size=2.5, mark options={solid}, only marks]
table {%
0.16 0.6
0.7 0.3
0.85 0.85
};
%
%
%
%
%
%
%
%
%
%
%
%
%
%
%
%
%
%
%
%
%
%
%
%
%
%
%
%
%
%
%
%
%
%
%
%
%
%
%
%
%
%
%
%
%
%
%
%
%
%
%
%
%
%
%
%
%
%
%
%
%
%
%
%
%
%
%
%
%
%
%
%
%
%
%
%
%
%
%
%
%
%
%
%
%
%
%
%
%
%
%
%
%
%
%
%
%
%
%
%
%
%
%
%
%
%
%
%
%
%
%
%
%
%
%
%
%
%
%
%
%
%
%
%
%
%
%
%
%
%
%
%
%
%
%
%
%
%
%
%
%
%
%
%
%
%
%
%
%
%
%
%
%
%
%
%
%
%
%
%
%
%
%
%
%
%
%
%
%
%
%
%
%
%
%
%
%
%
%
%
%
%
%
%
%
%
%
%
%
%
%
%
%
%
%
%
%
%
%
%
%
%
%
%
%
%
%
%
%
%
%
%
%
%
%
%
%
%
%
%
%
%
%
%
%
%
%
%
%
%
%
%
%
%
%
%
%
%
%
%
%
%
%
%
%
%
%
%
%
%
%
%
%
%
%
%
%
%
%
%
%
%
%
%
%
%
%
%
%
%
%
%
%
%
%
%
%
%
%
%
%
%
%
%
%
%
%
%
%
%
%
%
%
%
%
%
%
%
%
%
%
%
%
%
%
%
%
%
%
\addplot [semithick, darkviolet1910191, mark=*, mark size=2, mark options={solid}, only marks]
table {%
0.996246337890625 0.989500245558773
0.1885986328125 0.112787684804146
0.216583251953125 0.903994987214009
0.926422119140625 0.250588494301343
0.49468994140625 0.498602855255804
0.01495361328125 0.520245897474978
0.58428955078125 0.0354112686074277
0.61822509765625 0.876746430930244
0.9354248046875 0.624650713813951
0.2484130859375 0.374536401971244
0.27093505859375 0.640247929685515
0.4156494140625 0.205151653711324
0.679443359375 0.648529187623838
0.06268310546875 0.736066656505614
0.8076171875 0.436671239140375
0.4246826171875 0.971701468272113
0.0450439453125 0.304831580551745
0.4803466796875 0.704669003708784
0.7752685546875 0.11436264797033
};
\addplot [semithick, gray, mark=*, mark size=0.5, mark options={solid}, only marks]
table {%
0.996246337890625 0.989500245558773
0.931427001953125 0.993158224525394
0.994781494140625 0.92609527680401
0.965576171875 0.956713915561652
0.96392822265625 0.989229284153838
0.9981689453125 0.959152568206066
0.93719482421875 0.960778336635675
0.9805908203125 0.975003810394757
0.98333740234375 0.941269115480364
0.950103759765625 0.973649003370082
0.9783935546875 0.997764568409287
0.9476318359375 0.995732357872276
};
\addplot [semithick, gray, mark=*, mark size=0.5, mark options={solid}, only marks]
table {%
0.24957275390625 0.124574505918813
0.18218994140625 0.128232484885434
0.22119140625 0.15546410608139
0.21368408203125 0.0977493268302596
0.1556396484375 0.155057663973988
0.248291015625 0.092059137326627
0.19171142578125 0.157496316618402
0.208740234375 0.128638926992836
0.2314453125 0.111568358481939
0.23583984375 0.140832190214906
0.172027587890625 0.165083235956578
0.170654296875 0.14327084285932
};
\addplot [semithick, gray, mark=*, mark size=0.5, mark options={solid}, only marks]
table {%
0.24871826171875 0.875273078290911
0.186920166015625 0.873511829158834
0.216583251953125 0.903994987214009
0.27783203125 0.904435299497028
0.2196044921875 0.846415688665346
0.249176025390625 0.90765296618063
0.2481689453125 0.848041457094955
0.21392822265625 0.878118173042727
0.2305908203125 0.863892699283646
0.20440673828125 0.859828278209622
0.2628173828125 0.861758878219783
0.260162353515625 0.890887229250284
};
\addplot [semithick, gray, mark=*, mark size=0.5, mark options={solid}, only marks]
table {%
0.865264892578125 0.248082101305695
0.931640625 0.314128943758573
0.926422119140625 0.250588494301343
0.8968505859375 0.222273027485648
0.89556884765625 0.275923385662755
0.92926025390625 0.283239343595997
0.8671875 0.275720164609053
0.89410400390625 0.250317532896408
0.8865966796875 0.295432606818066
0.91058349609375 0.291368185744043
0.878997802734375 0.233517925790445
0.939605712890625 0.265220410167827
};
\addplot [semithick, gray, mark=*, mark size=0.5, mark options={solid}, only marks]
table {%
0.49468994140625 0.498602855255804
0.49798583984375 0.560382055580958
0.43463134765625 0.502260834222425
0.499267578125 0.436823654930651
0.560333251953125 0.496587579806601
0.5281982421875 0.52314179749022
0.466217041015625 0.468390658605565
0.464752197265625 0.536672932649156
0.530670166015625 0.466104421751427
0.49615478515625 0.529086013310979
0.465118408203125 0.504293044759437
0.497802734375 0.466087486663618
};
\addplot [semithick, gray, mark=*, mark size=0.5, mark options={solid}, only marks]
table {%
0.0008544921875 0.500736676319667
0.06201171875 0.563328760859625
0.05853271484375 0.502768886856678
0.027496337890625 0.532710122101983
0.032989501953125 0.470524479669427
0.05908203125 0.534064929126657
0.030426025390625 0.561973953834951
0.00396728515625 0.471472844586699
0.0294189453125 0.502362444749276
0.061553955078125 0.477163034090332
0.018341064453125 0.580263848668055
0.041046142578125 0.579586445155718
};
\addplot [semithick, gray, mark=*, mark size=0.5, mark options={solid}, only marks]
table {%
0.551788330078125 0.0642686582329929
0.62127685546875 0.0671137529848092
0.58428955078125 0.0354112686074277
0.58941650390625 0.0939389320733628
0.523223876953125 0.0356822300123626
0.5537109375 0.0333790580704161
0.613037109375 0.039475689681451
0.59307861328125 0.0634557740181883
0.5712890625 0.0809327846364883
0.599029541015625 0.0198309878236719
0.63226318359375 0.0490270792054057
0.633453369140625 0.0228793036291893
};
\addplot [semithick, gray, mark=*, mark size=0.5, mark options={solid}, only marks]
table {%
0.61822509765625 0.876746430930244
0.6810302734375 0.87247878880252
0.616119384765625 0.937983708445528
0.55743408203125 0.875527104608037
0.646148681640625 0.841047265830073
0.59130859375 0.841792409693644
0.653564453125 0.908245694253925
0.58270263671875 0.906823146878016
0.6510009765625 0.870446578265508
0.5869140625 0.871056241426612
0.56072998046875 0.84626327287507
0.61968994140625 0.906010262663212
};
\addplot [semithick, gray, mark=*, mark size=0.5, mark options={solid}, only marks]
table {%
0.874267578125 0.622008840115836
0.99420166015625 0.62099273484733
0.93499755859375 0.682771935172484
0.93267822265625 0.561245745059188
0.9354248046875 0.624650713813951
0.966461181640625 0.585903232908263
0.8980712890625 0.589696692577351
0.964263916015625 0.655404833274061
0.90576171875 0.657978966620942
0.929931640625 0.592135345221765
0.875274658203125 0.651069450795102
0.965362548828125 0.621805619062135
};
\addplot [semithick, gray, mark=*, mark size=0.5, mark options={solid}, only marks]
table {%
0.2484130859375 0.374536401971244
0.2447509765625 0.43834781283341
0.247833251953125 0.311402394621416
0.18743896484375 0.376568612508256
0.30865478515625 0.380937865162831
0.279541015625 0.413046791647615
0.214630126953125 0.344324205321005
0.2158203125 0.407864654778235
0.27587890625 0.344764517604024
0.2451171875 0.342021033379058
0.250885009765625 0.404376026689698
0.21929931640625 0.373317075649037
};
\addplot [semithick, gray, mark=*, mark size=0.5, mark options={solid}, only marks]
table {%
0.249053955078125 0.62531118223848
0.30657958984375 0.690240308896002
0.309844970703125 0.626835340141239
0.24578857421875 0.685058172026622
0.28045654296875 0.65853782451862
0.277435302734375 0.59337160663178
0.21966552734375 0.654575013971447
0.3128662109375 0.656912056089011
0.279815673828125 0.623177361174618
0.24493408203125 0.591576487324087
0.220489501953125 0.618672627817575
0.2774658203125 0.691053193110806
};
\addplot [semithick, gray, mark=*, mark size=0.5, mark options={solid}, only marks]
table {%
0.368682861328125 0.250063506579282
0.431304931640625 0.186726278175752
0.432037353515625 0.248911920608308
0.370147705078125 0.189097190468933
0.404205322265625 0.218428762553134
0.3995361328125 0.153939948178631
0.462982177734375 0.214770783586513
0.4013671875 0.25148605395519
0.4332275390625 0.157597927145252
0.43304443359375 0.217751359040797
0.40264892578125 0.187674643093024
0.46051025390625 0.246202306558959
};
\addplot [semithick, gray, mark=*, mark size=0.5, mark options={solid}, only marks]
table {%
0.680206298828125 0.679418787786415
0.68621826171875 0.616013819031652
0.647857666015625 0.647174380599163
0.718994140625 0.649748513946045
0.679443359375 0.648529187623838
0.688507080078125 0.585530660976477
0.70977783203125 0.679825229893817
0.6832275390625 0.713153482700808
0.65142822265625 0.618858913783468
0.647613525390625 0.681315517620959
0.712677001953125 0.622787854155024
0.61956787109375 0.64710664024793
};
\addplot [semithick, gray, mark=*, mark size=0.5, mark options={solid}, only marks]
table {%
0.12420654296875 0.749072803942488
0.062042236328125 0.688106487832139
0.094635009765625 0.78023336551
0.091064453125 0.72306050906874
0.09454345703125 0.684854950972921
0.08740234375 0.749885688157293
0.0650634765625 0.718183203779912
0.11761474609375 0.723466951176142
0.105987548828125 0.757608088197937
0.10882568359375 0.70395773002083
0.0511474609375 0.670629477213839
0.078887939453125 0.695422445765381
};
\addplot [semithick, gray, mark=*, mark size=0.5, mark options={solid}, only marks]
table {%
0.8076171875 0.436671239140375
0.801300048828125 0.374892038815221
0.744903564453125 0.436230926857356
0.8094482421875 0.498450439465529
0.86676025390625 0.431387491744145
0.840484619140625 0.464580263848668
0.77008056640625 0.402123660011177
0.775390625 0.462277091906722
0.835723876953125 0.406052600382733
0.80303955078125 0.405781638977798
0.778228759765625 0.430439126826873
0.80615234375 0.465935070873342
};
\addplot [semithick, gray, mark=*, mark size=0.5, mark options={solid}, only marks]
table {%
0.432373046875 0.941624752324341
0.371490478515625 0.934782976849735
0.46240234375 0.972107910379515
0.394561767578125 0.966417720875883
0.4613037109375 0.936747447035513
0.40032958984375 0.937153889142915
0.4246826171875 0.971701468272113
0.437042236328125 0.910328710054362
0.4678955078125 0.998933089468069
0.36444091796875 0.971091805111009
0.44464111328125 0.983894731494183
0.44482421875 0.961133973479652
};
\addplot [semithick, gray, mark=*, mark size=0.5, mark options={solid}, only marks]
table {%
0.0640869140625 0.312147538484987
0.0926513671875 0.277599959355789
0.0927734375 0.347508001828989
0.036651611328125 0.34059848600315
0.058868408203125 0.282070822537215
0.093017578125 0.313366864807194
0.063018798828125 0.345475791291978
0.104736328125 0.295076969974089
0.073883056640625 0.293722162949415
0.04681396484375 0.360107707158462
0.0450439453125 0.304831580551745
0.072784423828125 0.329624549103287
};
\addplot [semithick, gray, mark=*, mark size=0.5, mark options={solid}, only marks]
table {%
0.4957275390625 0.750190519737845
0.436676025390625 0.685430743958407
0.43255615234375 0.750596961845247
0.49462890625 0.690443529949703
0.460113525390625 0.718352554657996
0.5308837890625 0.718640451150739
0.469970703125 0.659960371894528
0.503570556640625 0.659977306982337
0.525390625 0.684499314128944
0.5264892578125 0.747904282883707
0.49169921875 0.71970736168267
0.4620361328125 0.746532540771224
};
\addplot [semithick, gray, mark=*, mark size=0.5, mark options={solid}, only marks]
table {%
0.741485595703125 0.134379921759894
0.802734375 0.133058984910837
0.740478515625 0.0632525529644871
0.766021728515625 0.0975630408643669
0.71630859375 0.101051668952904
0.7686767578125 0.163542142966011
0.80035400390625 0.10216938474826
0.772247314453125 0.131704177886162
0.800811767578125 0.163948585073414
0.71209716796875 0.127470405934055
0.828277587890625 0.128046198919541
0.7364501953125 0.160798658741046
};
\end{axis}

\end{tikzpicture}

\caption{Visualization of the key steps of the proof of \Cref{th:inverse_statement_generalized}.
For the sake of a helpful visualization, $q=0.2$ was chosen, despite it does not match the assumptions. \\
\textbf{Top left}: Initial set $Y_N$ (gray), reference set $Y_0$ (red) with circles of radius $\frac{1}{2} q = 0.1$ and set $Y_1'$ (purple, see Eq.~\eqref{eq:definition_y'}). 
\textbf{Top right}: Set $Y_1''$ (purple, see Eq.~\eqref{eq:definition_y''}) and clusters of points refering to the sets $T_i^{(N)}$ (gray, see Eq.~\eqref{eq:proof:definition_TiN}).
%
\textbf{Bottom left}: Restriction of the sets $T_i^{(N)}$ to $\tilde{T}_i^{(N)}$, see Eq.~\eqref{eq:proof:definition_TiN_tilde}. 
\textbf{Bottom right}: Additional visualization of $Y_1$ (red crosses, see Eq.~\eqref{eq:set_Y1}).}
\label{fig:vis_proof}
\end{figure}


\begin{proof}[Proof of \Cref{th:utility_theorem_new}]
%
As a first step, we approximate the integral $\int_\Omega |g(y)|^2 ~ \mathrm{d}y = \Vert g \Vert_{L^2(\Omega)}$ with a quadrature:
Let $(y_j)_{j \in \N}$ be the sequence of geometric greedy points in $\Omega \subset \R^d$, $Y_N := \{y_1, ..., y_N\}$. 
Let $N \in \N$:
\begin{itemize}
\item We cover $\Omega$ with patches $\{ E_i^{(N)} \}_{i=1}^N$, 
where $E_j^{(N)} \subset \Omega$ consists of all the points that are closer to $y_j$ than to any other $y_i, i=1, ..., N, i \neq j$:
\begin{align*}
E_j^{(N)} := \{y \in \Omega ~:~ \Vert y - y_j \Vert_2 \leq \Vert y - y_i \Vert_2 \quad \forall i \neq j \}.
\end{align*}
Like this it follows that $\bigcup_{i=1}^N E_i^{(N)} = \Omega$ and $\mu(E_i^{(N)} \cap E_j^{(N)}) = 0$ for all $i \neq j$.
Furthermore all the patches $E_i^{(N)}, i=1, ..., N$ are asymptotically of the same size:
On the one hand, using that for any $i=1, ..., N$, 
the patch $E_i^{(N)}$ is included in a ball of radius $h_{Y_N, \Omega}$ around $y_i$, it holds
\begin{align}
\label{eq:patch_size_upper_bound}
\mu(E_i^{(N)}) \leq C_d h_{Y_N, \Omega}^d \stackrel{\text{Eq.\ } \eqref{eq:geometric_uniformity}}{\leq} C_d 2^d q_{Y_N}^d \stackrel{\text{Eq.\ } \eqref{eq:geometric_bounds}}{\leq} 2^d C_d C_\Omega^d N^{-1}.
\end{align}
On the other hand, for any $i=1, ..., N$, 
the cone $C(y_i, \xi(y_i), \alpha, q_{Y_N})$ associated to $y_i$ is contained in $E_i^{(N)}$, thus
\begin{align}
\label{eq:patch_size_lower_bound}
\mu(E_i^{(N)}) \geq C_{d, \alpha} q_{Y_N}^d \stackrel{\text{Eq.\ } \eqref{eq:geometric_uniformity}}{\geq} 
C_{d, \alpha} 2^{-d} h_{Y_N, \Omega}^d 
\stackrel{\text{Eq.\ } \eqref{eq:geometric_bounds}}{\geq}  2^{-d} c_\Omega^{d} C_{d, \alpha} N^{-1}.
\end{align}
%
%
\item For $N \rightarrow \infty$ it holds $\max_{i=1, ..., N} \sup_{a, b \in \overline{E_i^{(N)}}}( \Vert a - b \Vert_2) \rightarrow 0$, 
therefore for any integrable function $g \in \mathcal{C}(\Omega)$ ($g \in \mathcal{C}(\Omega)$ bounded provides integrability) the definition of the Rieman integral yields
\begin{align*}
\int_{\Omega} |g(y)|^2 ~ \mathrm{d}y = \lim_{N \rightarrow \infty} \sum_{i=1}^N |g(y_i)|^2 \cdot \mu(E_i^{(N)}).
\end{align*}
%
%
%
In particular we have
\begin{alignat*}{2}
&\forall \varepsilon > 0 ~ \exists N_\varepsilon \in \N ~ \forall N \geq N_\epsilon && \quad \left| \int_{\Omega} |g(y)|^2 ~ \mathrm{d}y - \sum_{i=1}^N |g(y_i)|^2 \cdot \mu(E_i^{(n)}) \right| < \varepsilon \\
\Rightarrow \quad &\exists \tilde{N} \in \N ~ \forall N \geq \tilde{N}&& \quad \sum_{i=1}^N |g(y_i)|^2 \cdot \mu(E_i^{(n)}) < 2 \cdot \int_{\Omega} |g(y)|^2 ~ \mathrm{d}y.
\end{alignat*}
Using Eq.\ \eqref{eq:patch_size_lower_bound} we obtain finally
\begin{align}
\label{eq:desired_estimate}
\forall N > \tilde{N}& \quad \frac{c_\Omega^{d} C_{d, \alpha}}{2^{d}} \cdot \frac{1}{N} \cdot \sum_{i=1}^N |g(y_i)|^2 < 2 \cdot \int_{\Omega} |g(y)|^2 ~ \mathrm{d}y \notag \\
\Leftrightarrow \forall N > \tilde{N}& \quad \frac{1}{N} \cdot \sum_{i=1}^N |g(y_i)|^2 < \frac{2^{d+1}}{c_\Omega^{d} C_{d, \alpha}} \cdot \int_{\Omega} |g(y)|^2 ~ \mathrm{d}y.
\end{align}
\end{itemize}

The set $Y_N$ of points can be quite huge as $N$ is not further specified.
Thus, as a next step we want to construct smaller sets of points that still satisfy an inequality as in Eq.\ \eqref{eq:desired_estimate}.
This is done by selecting a suitable subset of $Y_N$ (namely the set $Y_1$, see Eq.~\eqref{eq:set_Y1}): \\
For $q \leq \min \left( \frac{2^{-1/d} c_\Omega}{7 C_\Omega} \cdot q_{Y_0}, \frac{2}{5}r \right) \leq \frac{c_\Omega}{7 C_\Omega} \cdot q_{Y_0} < \frac{1}{7} \cdot q_{Y_0}$, we define 
\begin{align}
\label{eq:definition_y'}
Y'_1 := \{ y_1, y_2, ..., y_{n^*}\}
\end{align}
with $n^* \geq 2$ the smallest integer such that $h_{Y_1' \cup \{ y_{n* + 1} \}, \Omega} < q$.
%
From this definition we conclude
\begin{align}
\label{eq:proof:geom_props_1}
\frac{1}{2} q \leq \frac{1}{2} h_{Y'_1, \Omega} \leq h_{Y'_1 \cup \{ y_{n^*+1} \}, \Omega} < q
\end{align}
where the middle inequality is because of Eq.~\eqref{eq:geometric_uniformity}.
For the separation distance $q_{Y_1'}$ we have due to Eq.~\eqref{eq:geometric_uniformity}
\begin{align}
\label{eq:proof:geom_props_2}
\frac{1}{2} q \leq \frac{1}{2} h_{Y_1', \Omega} \leq q_{Y_1'} \leq h_{Y_1', \Omega} < 2q.
\end{align}
With this, we define $Y''_1 \subset Y'_1$ as the subset of all the points which are not too close to $Y_0$, i.e.\
\begin{align}
\label{eq:definition_y''}
Y''_1 := Y'_1 \cap \left( \Omega \setminus \bigcup_{y \in Y_0} B_{\frac{1}{2} q}(y) \right),
\end{align}
where $B_r(y) := \{ z \in \R^d ~ | ~ \Vert z - y \Vert_2 < r \}$, i.e.\ open balls. 

The set $Y''_1$ is non-empty and well distributed:
%
\begin{itemize}
\item The set $Y_1''$ consists of at least two points:
Let $y \neq y' \in Y_0$ be the closest neighboring points of $Y_0$, i.e.\ $\Vert y - y' \Vert_2 = 2q_{Y_0} > 14q$, 
and thus $\mathrm{dist}(B_{\frac{1}{2}q}(y), B_{\frac{1}{2}q}(y')) > 13q$.
On the other hand, $h_{Y_1', \Omega} < 2q$, i.e.\ there must be at least two points in $Y_1''$ (namely ``in between'' $y$ and $y'$ as $\Omega$ is connected).
%
\item For the separation distance $q_{Y''_1}$ we simply have by Eq.~\eqref{eq:proof:geom_props_2}
\begin{align*}
q_{Y''_1} \geq q_{Y'_1} \geq \frac{1}{2} q.
\end{align*}
\item For the fill distance $h_{Y''_1, \Omega}$ we calculate:
\begin{align*}
h_{Y''_1, \Omega} \equiv& \sup_{y \in \Omega} \min_{y'' \in Y''_1} \Vert y - y'' \Vert_2 \\
=& \sup_{y \in \Omega} \min_{y'' \in Y''_1} \min_{y' \in Y'_1} \Vert y - y' + y' - y'' \Vert_2 \\
\leq& \sup_{y \in \Omega} \min_{y'' \in Y''_1} \min_{y' \in Y'_1} \Vert y - y' \Vert_2 + \Vert y' - y'' \Vert_2.
\end{align*}
Now pick $\tilde{y}' \in Y'_1$ as the closest $y'$ to $y$, i.e.\ $\tilde{y}' := \tilde{y}'(y) := \argmin_{y' \in Y'_1} \Vert y - y' \Vert_2$.
Then we continue to estimate:
\begin{align*}
h_{Y''_1, \Omega} \leq& \sup_{y \in \Omega} \left( \min_{y'' \in Y''_1} \Vert y - \tilde{y}'(y) \Vert_2 + \Vert \tilde{y}'(y) - y'' \Vert_2 \right) \\
\leq& \sup_{y \in \Omega} \left( \Vert y - \tilde{y}'(y) \Vert_2 + \min_{y'' \in Y''_1}  \Vert \tilde{y}'(y) - y'' \Vert_2 \right) \\
\leq& \sup_{y \in \Omega} \Vert y - \tilde{y}'(y) \Vert_2 + \sup_{y \in \Omega} \min_{y'' \in Y''_1}  \Vert \tilde{y}'(y) - y'' \Vert_2 \\
\leq& h_{Y'_1, \Omega} + \sup_{y \in \Omega} \min_{y'' \in Y''_1}  \Vert \tilde{y}'(y) - y'' \Vert_2.
%
%
%
%
%
%
\end{align*}
If $\tilde{y}'(y) \in Y_1''$, then the second summand is obviously zero.
If $\tilde{y}'(y) \in Y_1' \setminus Y_1''$, then by Eq.\ \eqref{eq:definition_y''} there is a point $y_0(y) \in Y_0$ such that $\Vert \tilde{y}'(y) - y_0(y) \Vert_2 < \frac{1}{2} q$. 
Then we choose any point $\tilde{y} := \tilde{y}(y) \in \Omega$ with $\Vert y_0(y) - \tilde{y}(y) \Vert_2 = \frac{5}{2} q$.
%
Such a point exists as $\frac{5}{2} q \leq r$, i.e. the distance is smaller than the radius of the cone condition of $\Omega$.
Adding $0 = - y_0(y) + y_0(y) - \tilde{y}(y) + \tilde{y}(y)$ gives
\begin{align*}
\sup_{y \in \Omega} \min_{y'' \in Y''_1}  \Vert \tilde{y}'(y) - y'' \Vert_2 
&= \sup_{y \in \Omega} \min_{y'' \in Y''_1}  \Vert \tilde{y}'(y) - y_0(y) + y_0(y) - \tilde{y}(y) + \tilde{y}(y) - y'' \Vert_2 \\
&\leq \sup_{y \in \Omega} \min_{y'' \in Y''_1}  \Vert \tilde{y}'(y) - y_0(y) \Vert_2 + \Vert y_0(y) - \tilde{y}(y) \Vert_2+ \Vert \tilde{y}(y) - y'' \Vert_2 \\
&\leq \sup_{y \in \Omega} \Vert \tilde{y}'(y) - y_0(y) \Vert_2 + \sup_{y \in \Omega} \Vert y_0(y) - \tilde{y}(y) \Vert_2 + \sup_{y \in \Omega} \min_{y'' \in Y''_1} \Vert \tilde{y}(y) - y'' \Vert_2 \\
&\leq \frac{1}{2} q + \frac{5}{2} q + 2q = 5q.
\end{align*}
In the last step we used for the last term that the closest point $y^*(y) \in Y'_1$ to $\tilde{y}(y)$ has a distance of at most $2q$ (due to $h_{Y'_1, \Omega} < 2q$) and that this point must actually belong to $Y_1''$ 
(because $\Vert y_0(y) - y^*(y) \Vert_2 = \Vert y_0(y) - \tilde{y}(y) + \tilde{y}(y) - y^*(y) \Vert_2 \geq \left| \Vert y_0(y) - \tilde{y}(y) \Vert - \Vert \tilde{y}(y) - y^*(y) \right| > |\frac{5}{2} q - 2q | = \frac{1}{2} q$). \\
Thus, finally we obtain
\begin{align*}
h_{Y_1'', \Omega} \leq 5q + 2q = 7q.
\end{align*}
\end{itemize}
For ease of notation, we set $Y_1'' =: \{y''_1, ..., y''_{|Y_1''|} \}$.
Now we create sets $T_i^{(N)}$ by assigning nearby points from $Y_N$ (for $N$ large enough) to the points from $Y_1''$:
\begin{align}
\label{eq:proof:definition_TiN}
T_i^{(N)} :=& \{ z \in Y_N ~ | ~ \Vert y''_i - z \Vert_2 \leq \frac{1}{3} q \}, \qquad &&i=1, ..., |Y_1''|.
\end{align}
The distance $\frac{1}{3}q$ was chosen such that $T_i^{(N)} \cap T_j^{(N)} = \emptyset$ for any pairwise distinct $i, j = 1, ..., |Y_1''|$. 
For convenience we introduce
\begin{align*}
T_i^{(N)} =: \left\{ y_1^{(i, N)}, ..., y_{|T_i^{(N)}|}^{(i, N)} \right\}.
\end{align*}
Now we prove that for $N \rightarrow \infty$, all the sets $T_i^{(N)}$, $i=1, ..., |Y_1''|$ are of asymptotically equal size:
\begin{itemize}
\item The number of points in $T_i^{(N)}$ must be large enough, to at least allow to cover the surrounding of $y_i'' \in Y_1''$ in terms of the volume:
Due to the assumption on the cone condition, 
the surrounding of $y_i''$ consist at least of the associated cone $C(y_i'', \xi(y_i''), \alpha, \frac{1}{3} q)$:
%
\begin{align*}
\bigcup_{y \in T_i^{(N)}} B_{h_{Y_N, \Omega}}(y) &\supset C(y_i'', \xi(y_i''), \alpha, \frac{1}{3}q) \\
\Rightarrow \quad \left|T_i^{(N)} \right| \cdot C_d h_{Y_N, \Omega}^{d} &\geq \mu(C(y_i'', \xi(y_i''), \alpha, \frac{1}{3}q)) = 3^{-d} C_{d, \alpha} \cdot q^d.
\end{align*}
Rearranging and using the geometric bounds from Eq.\ \eqref{eq:geometric_bounds} and Eq.\ \eqref{eq:geometric_uniformity} gives
\begin{align}
\label{eq:subset_size_lower_bound}
\left| T_i^{(N)} \right| &\geq \frac{3^{-d} C_{d, \alpha}}{C_d} \cdot \frac{q^d}{h_{Y_N, \Omega}^d}
\geq \frac{3^{-d} C_{d, \alpha}}{C_d} \cdot \frac{q^d}{2^d C_\Omega^d N^{-1}} \notag \\
&= \frac{C_{d, \alpha} q^d}{6^d C_d C_\Omega^d} \cdot N.
\end{align}
\item Using again a volume comparison argument, we have (for $N > n$ large enough such that $q_{Y_N} \leq \frac{2}{3} q$):
\begin{align*}
%
\bigcup_{y \in T_i^{(N)}} B_{q_{Y_N}}(y) &\subset B_{\frac{1}{3}q + q_{Y_N}}(y''_i) \\
\Rightarrow \quad \left| T_i^{(N)} \right| \cdot C_d q_{Y_N}^d &\leq C_d \cdot \left( \frac{1}{3}q + q_{Y_N} \right)^d \leq C_d q^d.
\end{align*}
Using the geometric bounds from Eq.\ \eqref{eq:geometric_bounds} and the asymptotical uniform distribution of Eq.\ \eqref{eq:geometric_uniformity} we arrive at
\begin{align}
\label{eq:subset_size_upper_bound}
\left| T_i^{(N)} \right| &\leq q^d \cdot q_{Y_N}^{-d} \leq q^d 2^d c_\Omega^d N.
\end{align}
\item Combining Eq.\ \eqref{eq:subset_size_lower_bound} and Eq.\ \eqref{eq:subset_size_upper_bound} formalizes our claim that the sets $T_j^{(N)}$ are of asymptotical equal size, 
because the lower and upper bound do actually not depend on $i = 1, ..., |Y_1''|$:
\begin{align}
\label{eq:subset_size_asymp}
\frac{C_{d, \alpha}}{6^d C_d C_\Omega^d} \cdot q^d N \leq 
\left| T_j^{(N)} \right| 
\leq 2^d c_\Omega^d \cdot q^d N.
\end{align}
\end{itemize}
Now we define $\tilde{T}_i^{(N)} \subset T_i^{(N)}$ for $i=1, ..., |Y_1''|$ be restricting the sets to the same amount of points, $m := \min_{i=1, ..., |Y_1''|} |T_i^{(N)}|$:
\begin{align}
\label{eq:proof:definition_TiN_tilde}
\tilde{T}_i^{(N)} := \left\{ y_1^{(i, N)}, ..., y_m^{(i, N)} \right\}.
\end{align}
We remark that the sets $\tilde{T}_i^{(N)}$ depend on the ordering of the points $\{ y_1^{(i, N)}, ..., y_{|T_i^{(N)}|}^{(i, N)} \}$, however it does not matter for the following which one is chosen. \\
Now we want to show that the number of points contained in 
\begin{align*}
\tilde{T}^{(N)} := \bigcup_{i=1}^{|Y_1''|} \tilde{T}_i^{(N)}
\end{align*}
is comparable to the number of points in the set $Y_N$ (which consists of $N$ points):
\begin{itemize}
\item For this we need a lower bound on $|Y_1''|$.
From Eq.\ \eqref{eq:proof:geom_props_2} we have $q_{Y_1'} \geq \frac{1}{2}q$, i.e.\ in every (open!) ball $B_{\frac{1}{2} q}(y)$ for $y \in Y_0$, there can be at most one point from $Y_1'$.
Therefore, we obtain $|Y_1''| \geq |Y_1'| - |Y_0|$.
Further using the geometric bound Eq.~\eqref{eq:geometric_bounds} and Eq.~\eqref{eq:proof:geom_props_1}, we obtain
\begin{align*}
|Y_1'| \geq c_\Omega^{d} h_{Y_1', \Omega}^{-d} \geq c_\Omega^d (2q)^{-d} = 2^{-d} c_\Omega^d q^{-d} 
%
\end{align*}
as well as
\begin{align*}
|Y_0| \leq C_\Omega^d q_{Y_0}^{-d} \leq C_\Omega^d \frac{1}{2} \frac{c_\Omega^d}{7^d C_\Omega^d} q^{-d} 
= \frac{1}{2} \frac{c_\Omega^d}{7^d} q^{-d}.
\end{align*}
Thus we have
\begin{align}
\label{eq:lower_bound_Y1''}
|Y_1''| \geq |Y_1'| - |Y_0| \geq& ~ \frac{c_\Omega^d}{2^d} q^{-d} - \frac{1}{2} \frac{c_\Omega^d}{7^d} q^{-d} \notag \\
>& \left( \frac{1}{2} \frac{c_\Omega^d}{2^d} \right) q^{-d}.
\end{align}
Now we can finally estimate the ratio of points contained in $\tilde{T}^{(N)}$ by putting together Eq.~\eqref{eq:subset_size_lower_bound} and \eqref{eq:lower_bound_Y1''}:
\begin{align}
\label{eq:subset_size_ratio}
\left| \tilde{T}^{(N)} \right| = \left| \bigcup_{i=1}^{|Y_1''|} \tilde{T}_i^{(N)} \right| \geq& ~ |Y_1''| \cdot \min_{i=1, ..., |Y_1''|} | \tilde{T}_i^{(N)} | \notag \\
\geq& \left( \frac{1}{2} \frac{c_\Omega^d}{2^d} \right) q^{-d} \cdot \frac{C_{d, \alpha}}{6^d C_d C_\Omega^d} \cdot q^d N \notag \\
=& \left( \frac{1}{2} \cdot \frac{c_\Omega^d C_{d, \alpha}}{12^d C_d C_\Omega^d} \right)  N, \notag \\
\Leftrightarrow \qquad \frac{N}{\left| \tilde{T}^{(N)} \right|} \leq& ~ 2 \cdot \frac{12^d C_d C_\Omega^d}{c_\Omega^d C_{d, \alpha}}.
\end{align}
\end{itemize}
Now we are ready to derive our main statement:
Recalling Eq.~\eqref{eq:proof:definition_TiN_tilde},
we consider the sum $\sum_{y \in \tilde{T}^{(N)}} |g(y)|$ (consisting of $|\tilde{T}^{(N)}| = |Y_1''| \cdot \min_{i=1, ..., |Y_1''|} |T_i^{(N)}|$ summands)
and split it into $m \equiv \min_{i=1, ..., |Y_1''|} |T_i^{(N)}|$ summands, such that every summand uses exactly one point from every $\tilde{T}_i^{(N)}$:
\begin{align*}
\frac{1}{|\tilde{T}^{(N)}|} \sum_{y \in \tilde{T}^{(N)}} |g(y)|^2
= \frac{1}{m} \sum_{\ell = 1}^m \left( \frac{1}{|Y_1''|} \sum_{i=1}^{|Y_1''|} |g(y_{\ell}^{(i, N)})|^2 \right).
\end{align*}
Now we have
\begin{align*}
\frac{1}{m} \sum_{\ell = 1}^m \left( \frac{1}{|Y_1''|} \sum_{i=1}^{|Y_1''|} |g(y_{\ell}^{(i, N)})|^2 \right) 
&~~~=~~~ \frac{1}{|\tilde{T}^{(N)}|} \sum_{y \in \tilde{T}^{(N)}} |g(y)|^2 \\
%
&~~~\leq~~~ \frac{N}{|\tilde{T}^{(N)}|} \cdot \frac{1}{N} \cdot \sum_{y \in Y_N} |g(y)|^2 \\
&\stackrel{\text{Eq.\ }\eqref{eq:desired_estimate}}{\leq} \frac{N}{|\tilde{T}^{(N)}|} \cdot \frac{2^{d+1}}{c_\Omega^{d} C_{d, \alpha}} \cdot \int_{\Omega} |g(y)|^2 ~ \mathrm{d}y \\ 
&\stackrel{\text{Eq.\ }\eqref{eq:subset_size_ratio}}{\leq} 2 \cdot \frac{12^d C_d C_\Omega^d}{c_\Omega^d C_{d, \alpha}} \cdot \frac{2^{d+1}}{c_\Omega^{d} C_{d, \alpha}} \cdot \int_{\Omega} |g(y)|^2 ~ \mathrm{d}y \\
&~~~=:~~~ \tilde{C}_{d, \alpha} \cdot \int_{\Omega} |g(y)|^2 ~ \mathrm{d}y.
\end{align*}
From here we can conclude that there exists at least one $\ell^* \in \{ 1, ..., m \}$ such that
\begin{align}
\label{eq:desired_estimate_2}
\frac{1}{|Y_1''|} \sum_{i=1}^{|Y_1''|} |g(y_{\ell^*}^{(i, N)})|^2 \leq \tilde{C}_{d, \alpha} \cdot \int_{\Omega} |g(y)|^2 ~ \mathrm{d}y,
\end{align}
which is exactly the inequality Eq~\eqref{eq:utility_theorem_new} to be proven.
Finally we are left with analyzing the geometric properties of this set 
\begin{align}
\label{eq:set_Y1}
Y_1 := \{ y_{\ell^*}^{(i,N)} \}_{i=1}^{|Y_1''|} = \{ y_{\ell^*}^{(i,N)} \}_{i=1}^{|Y_1|}.
\end{align}

%

\begin{itemize}
\item For the fill distance $h_{Y_1, \Omega}$ we calculate
\begin{align}
\label{eq:proof:subset_upper_bound_fill_dist}
h_{Y_1, \Omega} &\equiv \sup_{y \in \Omega} \min_{i=1, ..., |Y_1|} \Vert y - y_{\ell^*}^{(i, N)} \Vert_2 \notag \\
&\leq \sup_{y \in \Omega} \min_{i=1, ..., |Y_1|} \Vert y - \argmin_{y_j'' \in Y_1''}(\Vert y - y_j'' \Vert_2) + \argmin_{y_j'' \in Y_1''}(\Vert y - y_j'' \Vert_2) - y_{\ell^*}^{(i,N)} \Vert_2 \notag \\
&\leq \sup_{y \in \Omega} \min_{i=1, ..., |Y_1|} \Vert y - \argmin_{y_j'' \in Y_1''}(\Vert y - y_j'' \Vert_2) \Vert_2 + \Vert \argmin_{y_j'' \in Y_n}(\Vert y - y_j'' \Vert_2) - y_{\ell^*}^{(i,N)} \Vert_2 \notag \\
&\leq \sup_{y \in \Omega} \Vert y - \argmin_{y_j'' \in Y_1''}(\Vert y - y_j'' \Vert_2) \Vert_2 + \sup_{y \in \Omega} \min_{i=1, ..., |Y_1''|} \Vert \argmin_{y_j'' \in Y_1''}(\Vert y - y_j'' \Vert_2) - y_{\ell^*}^{(i,N)} \Vert_2 \notag \\
&\leq h_{Y_1'', \Omega} + \frac{1}{3} q 
\leq 7q + \frac{1}{3}q = \frac{22}{3}q,
\end{align}
where the definition of the sets $T_i^{(n, N)}$ from Eq.\ \eqref{eq:proof:definition_TiN} was used for the final step.
%
%
%
%
%
\item For the separation distance $q_{Y_1}$ it holds
\begin{align*}
q_{Y_1} &\equiv \min_{\substack{i \neq j \\ i, j = 1, ..., |Y_1|}} \Vert y_{\ell^*}^{(i,N)} - y_{\ell^*}^{(j,N)} \Vert_2 \\
&= \min_{\substack{i \neq j \\ i, j = 1, ..., |Y_1|}} \left \Vert y_{\ell^*}^{(i,N)} - y''_i + y''_i - y''_j + y''_j - y_{\ell^*}^{(j,N)} \right \Vert_2 \\
&\geq \min_{\substack{i \neq j \\ i, j = 1, ..., |Y_1|}} \left| \left \Vert y_i'' - y_j'' \right \Vert_2 - \left \Vert y_{\ell^*}^{(i,N)} - y''_i + y''_j - y_{\ell^*}^{(j,N)} \right \Vert_2 \right|.
\end{align*}
The first term can be lower bounded by $2 \cdot q_{Y_1''} \geq q$ due to the definition of the separation distance of $Y_1''$, 
while the latter term can be upper bounded (after applying the triangle inequality) by $\frac{1}{3} q + \frac{1}{3} q$ due to the definition of the sets $T_i^{(N)}$ in Eq.\ \eqref{eq:proof:definition_TiN}.
Thus we finally obtain
\begin{align}
\label{eq:proof:subset_lower_bound_sep_dist}
q_{Y_1} \geq 2q_{Y_1''} - \frac{2}{3} q \geq q - \frac{2}{3}q = \frac{1}{3}q.
\end{align}
\item Combining Eq.\ \eqref{eq:proof:subset_upper_bound_fill_dist} and Eq.\ \eqref{eq:proof:subset_lower_bound_sep_dist} gives
\begin{align*}
\frac{1}{3} q \leq q_{Y_1} \leq h_{Y_1, \Omega} \leq \frac{22}{3}q,
\end{align*}
i.e.\ also the set $Y_1$ is well distributed.
\end{itemize}
Finally we want to show that also the union $Y_0 \cup Y_1$ is well distributed:
\begin{align*}
q_{Y_0 \cup Y_1} =& \min (q_{Y_0}, q_{Y_1}, \min_{y0 \in Y_0, y_1 \in Y_1} \Vert y_0 - y_1 \Vert_2) \\
\geq& \min \left(q_{Y_0}, \frac{1}{3}q, \frac{1}{2}q - \frac{1}{3}q \right) = \frac{1}{6}q,
\end{align*}
where the last quantity was estimated with Eq.~\eqref{eq:definition_y''} and Eq.~\eqref{eq:proof:definition_TiN}.
\end{proof}






\subsection{Proof of \Cref{prop:disc_L2_from_cont_L2_new}}

\begin{proof}[Proof of \Cref{prop:disc_L2_from_cont_L2_new}]
Let $(y_j)_{j \in \N}$ be the sequence of geometric greedy points in $\Omega \subset \R^d$, $Y_m := \{ y_1, ..., y_m \}$ for all $m \in \N$ as specified in \Cref{subsec:geometric_properties}.
We define $m_0 \in \{2, 3, ...\} \subset \N$ as the smallest integer such that $Y_{m_0}$ satisfies 
\begin{align*}
h_{Y_{m_0}, \Omega} < h_0 ~~ \text{and} ~~ \frac{2^{-1/d}c_\Omega}{7C_\Omega} \cdot q_{Y_{m_0}} < \frac{2}{5}r.
\end{align*}
We then define $X_0 := Y_{m_0}$. 
Then the assumption of this \Cref{prop:disc_L2_from_cont_L2_new} yields
\begin{align*}
\Vert f - s_{f, X_{0}} \Vert_{L^2(\Omega)} \leq c_f \cdot h_{X_{0}}^\mu.
\end{align*}
Define 
\begin{align}
\label{eq:definition_a}
a := \frac{3}{22} \cdot \frac{2^{-1/d} c_\Omega}{7C_\Omega} \in \left( 0, \frac{3}{154} \right) \subset (0, 1).
\end{align}
%
%
Now \Cref{th:utility_theorem_new} can be applied for $g := f - s_{f, X_{0}}$ (which is continuous and bounded),
the reference set $Y_0 := X_0$ and $q := a \cdot q_{X_{0}}$ (which satisfies $q \leq \min \left( \frac{2^{-1/d} c_\Omega}{7 C_\Omega} \cdot q_{Y_0}, \frac{2}{5}r \right)$ ), 
which yields a set $Y_1 \subset \Omega$ that satisfies
\begin{align}
\begin{aligned}
\label{eq:proof:properties_initial}
\frac{1}{3} a q_{X_{0}} \leq& ~ q_{Y_{1}} \leq h_{Y_{1}, \Omega} \leq \frac{22}{3} a  q_{X_{0}}, \\
\frac{1}{6} aq_{X_{0}} \leq& ~ q_{X_{0} \cup Y_{1}}, \\
\Vert f - s_{f, X_{0}} \Vert_{L^2(Y_{1})} \leq& ~ \sqrt{\tilde{C}_{d, \alpha}} \cdot \Vert f - s_{f, X_{0}} \Vert_{L^2(\Omega)}.
\end{aligned}
\end{align}
Now we simply define $X_1 := X_0 \cup Y_1$. Due to $s_{f, X_0}|_{X_0} = f|_{X_0}$,
Eq.~\eqref{eq:proof:properties_initial} turns into
\begin{align}
\begin{aligned}
\label{eq:proof:properties_initial2}
\frac{1}{6} a q_{X_{0}} \leq& ~ q_{X_{1}} \leq h_{X_{1}, \Omega} \leq \frac{22}{3} a  q_{X_{0}}, \\
\Vert f - s_{f, X_{0}} \Vert_{L^2(X_{1})} \leq& ~ \sqrt{\tilde{C}_{d, \alpha}} \cdot \Vert f - s_{f, X_{0}} \Vert_{L^2(\Omega)}.
\end{aligned}
\end{align}
Now the idea is to iteratively apply \Cref{th:utility_theorem_new} to build the sequence $(X_m)_{m \geq 0}$ of sets. \\ 
We formalize this strategy via induction: 
\begin{enumerate}
\item The start of the induction is given by the sets $X_{0}$ and $X_{1}$ from above, 
which satisfy the desired geometric estimates on the distribution of the points as well as the estimate between discrete $L^2(X)$ norm and continuous $L^2(\Omega)$ norm, 
see Eq.~\eqref{eq:proof:properties_initial2}.
\item Assume for $m \geq 0$ we have two sets $X_{m}, X_{m+1}$ such that 
\begin{align}
\frac{1}{6} a^{m+1} q_{X_{0}} \leq& ~ q_{X_{m+1}} \leq h_{X_{m + 1}, \Omega} \leq \frac{22}{3} a^{m+1} q_{X_{0}}, \label{eq:proof:properties_assump1} \\
%
\Vert f - s_{f, X_{m}} \Vert_{L^2(X_{m + 1})} \leq& ~ \sqrt{\tilde{C}_{d, \alpha}} \cdot \Vert f - s_{f, X_{m}} \Vert_{L^2(\Omega)}. \notag %
\end{align}
Now we can leverage again \Cref{th:utility_theorem_new} for $g := f - s_{f, X_{m+1}}$ and $q: = a^{m+2} q_{X_0}$ and reference set $Y_0 := X_{m+1}$, as the choice of $q$ satisfies the requirements:
\begin{align*}
q \equiv 6a \cdot \frac{1}{6} a^{m+1} q_{X_0} \stackrel{\text{Eq.}~\eqref{eq:proof:properties_assump1}}{\leq} 6a \cdot q_{X_{m+1}} \stackrel{\text{Eq.}~\eqref{eq:definition_a}}{\leq} \frac{18}{22} \cdot \frac{2^{-1/d} c_\Omega}{7C_\Omega} \cdot q_{X_{m+1}}.
\end{align*}
Thus, \Cref{th:utility_theorem_new} gives a set $Y_1$ satisfying
\begin{align}
\begin{aligned}
\label{eq:proof:properties_initial3}
\frac{1}{3} a^{m+2} q_{X_0} \leq& ~ q_{Y_1} \leq h_{Y_1, \Omega} \leq \frac{22}{3} a ^{m+2} q_{X_0}, \\
\frac{1}{6} a^{m+2} q_{X_0} \leq& ~ q_{X_{m+1} \cup Y_1}, \\
\Vert f - s_{f, X_{m+1}} \Vert_{L^2(Y_1)} \leq& ~ \sqrt{\tilde{C}_{d, \alpha}} \cdot \Vert f - s_{f, X_{m+1}} \Vert_{L^2(\Omega)}. 
\end{aligned}
\end{align}
Now we define $X_{m+2} := X_{m+1} \cup Y_1$.
Due to $s_{f, X_{m+1}}|_{X_{m+1}} = f|_{X_{m+1}}$ Eq.~\eqref{eq:proof:properties_initial3} turns into
\begin{align*}
\frac{1}{6} a^{m+2} q_{X_0} \leq& ~ q_{X_{m+2}} \leq h_{X_{m+2}, \Omega} \leq \frac{22}{3} a^{m+2} q_{X_0}, \\
\Vert f - s_{f, X_{m+1}} \Vert_{L^2(X_{m+2})} \leq& ~ \sqrt{\tilde{C}_{d, \alpha}} \cdot \Vert f - s_{f, X_{m+1}} \Vert_{L^2(\Omega)}. 
\end{align*}
Thus, we have found a new pair of sets $X_{m+1}, X_{m+2}$ with corresponding properties and the induction step is finished.
\end{enumerate}
Overall we have found a nested sequence of sets $(X_m)_{m \geq 0}$ that satisfy the geometric properties for $m \geq 0$.
To obtain the desired estimate on the discrete $L^2$ norm, we just need to leverage the assumption on the decay of $\Vert f - s_{f, X} \Vert_{L^2(\Omega)}$.
The constant $c_0$ is given as $c_0 := q_{X_0}$ and only depends on $h_0$ and the region $\Omega$.
\end{proof}



%
%
%
%
%
%
%
%
%
%
%
%
%
%
%
%
%
%
%
%
%
%
%
%
%
%
%
%
%




%
%


%
%

%

%

%
%
%
%
%
%
%
%
%



\ifappendix

~ \newpage
~ \newpage

\section{Future research}

\begin{itemize}
\item \tw{I could do some numerical experiments on the analytic gap via mpmath, plotting stuff as in Eq.~\eqref{eq:discrete_Linfty_error}.}
\end{itemize}

\section{Discussion} \label{sec:disucssion}


\begin{idea}
Can \Cref{th:one-to-one_correspondence} help to answer the search for the maximal $\theta > 1$ such that $k \in \mathcal{H}_\theta$?
Especially we could consider $f = k(\cdot, x^*)$ for $x^* \in \Omega$ and interpolate it with \textit{any} well distributed points $X \subset \Omega$.
But not sure if this yields anything.
\end{idea}

\subsubsection{Further ideas}

\tw{Few more bullet points I came up with and which needs to be worked on:}

\begin{itemize}
\item In the setting of \Cref{th:inverse_statement_generalized}, assume $\vartheta > 1$.
Then we know $f \in \ns$.
This allows us to write (standard power function estimate + Duchons localization principle)
\begin{align*}
\Vert f - s_{f, X} \Vert_{L^\infty(\Omega)} &\leq c h_X^{\frac{s_\infty}{2}} \cdot \Vert f - s_{f, X} \Vert_{\ns} \\
\Rightarrow \qquad \Vert f - s_{f, X} \Vert_{L^2(\Omega)} &\leq c h_X^{\frac{s_\infty + d}{2}} \cdot \Vert f - s_{f, X} \Vert_{\ns}
\end{align*}
On the other hand we have
\begin{align*}
\Vert f - s_{f, X} \Vert_{L^2(\Omega)} &\leq c_f \cdot h_X^{\vartheta \cdot \frac{s_\infty + d}{2}},
\end{align*}
which limits the decay of $\Vert f - s_{f, X} \Vert_{\ns}$, we obtain a statement like
\begin{align*}
\sup_{\substack{X \subset \Omega \\ h_X \leq h}} \Vert f - s_{f, X} \Vert_{\ns} \geq c \cdot h^{(\vartheta_1 - 1) \cdot \frac{s_\infty + d}{2}}.
\end{align*}
See also 25.01.23\textbackslash A1 (also took a picture).
\item So far I assumed a faster polynomial convergence rate, i.e.\ e.g.\ $\vartheta' < \vartheta$.
However this can be refined using also logarithmic decay rates.
Working in the setting of \Cref{th:inverse_statement_generalized} and assuming
\begin{align*}
\Vert f - s_{f, X_m} \Vert_{L^2(X_{m+1})} \leq h_{X_m, \Omega}^{\frac{s_\infty + d}{2}} \cdot \log(h_{X_m, \Omega})^\beta
\end{align*}
we can calculate (similar to the proof of \Cref{th:inverse_statement_generalized} using $\vartheta' = 1$):
\begin{align*}
\Vert s_{f, X_{m+1}} - s_{f, X_m} \Vert_{\ns}^2 &\leq \Vert A_{X_{m+1} \cup X_m}^{-1} \Vert \cdot |X_{m+1}| \cdot \Vert s_{f, X_{m+1}} - s_{f, X_m} \Vert_{L^2(X_m+1)}^2 \\
&\leq C \cdot q_{X_m \cup X_{m+1}}^{-s_\infty} \cdot q_{X_{m+1}}^{-d} \cdot h_{X_m, \Omega}^{s_\infty + d} \cdot \log(h_{X_m, \Omega})^\beta \\
&\leq C \cdot a^{-s_\infty m} \cdot a^{-dm} \cdot a^{m (s_\infty + d)} (m \cdot \log(a))^\beta \\
&= C \cdot m^\beta.
\end{align*}
In order to have $\sum_{m=m_0}^\infty \Vert s_{f, X_{m+1}} - s_{f, X_m} \Vert_{\ns}^2 \leq C \cdot \sum_{m=m_0}^\infty m^\beta \stackrel{!}{=} \infty$, we need $\beta > 1$. \\
Like this we reduced $\vartheta' < \vartheta$ to $\vartheta' = \vartheta$ if we have additionally a logarithmic converging term as $\log(h_{X_m, \Omega})^\beta$ with $\beta > 1$. \\
It's pretty clear that this can be generalized beyond the $\vartheta' = 1$ case.
\end{itemize}


\subsection{Case $\vartheta \in (1, 2]$}

Aboves theorem and proof above does not give better results like $f \in \Ht$ for $\vartheta > 1$, 
which does not seem to be reasonable (why should we have such a break in this nice scale of convergence?):

\begin{itemize}
\item The approach of constructing a sequence of interpolants $s_n$ based on the kernel $k$, which converges to $f \in \Ht$ is intrinsically limited to $\vartheta \leq 1$:
\begin{itemize}
\item We estimated $\Vert \cdot \Vert_{\Ht}$ by $\Vert \cdot \Vert_{L^2(\Omega)}$ and $\Vert \cdot \Vert_{\ns}$.
If $\vartheta > 1$, we cannot use this anymore, because as an upper bound we would need something bigger, e.g.\ $\Vert \cdot \Vert_{\mathcal{H}_{\theta = 2}}$.
\item Furthermore the maximal $\vartheta^* > 1$ such that $k(\cdot, x) \in \Ht$ is unclear. 
If $f \in TL^2(\Omega)$, then it has a very fast rate of convergence. 
However as $k(\cdot, x) \notin TL^2(\Omega)$, we can never construct a sequence of interpolants based on the kernel $k$ that converges in $\Vert \cdot \Vert_{\mathcal{H}_{\theta = 2}}$ norm to $f$.
\end{itemize}
\item Therefore the idea is to use a power kernel $k^{(\theta)}(x, y)$ with $\vartheta > 1$ for the approximation:
\begin{align}
\label{eq:power_kernel}
k^{(\theta)}(x,y) = \sum_{n=1}^\infty \lambda_n^{\theta} \varphi_n(x) \varphi_n(y).
\end{align}
\begin{itemize}
\item Using a power kernel $k^{(\theta)}(x, y)$, one can construct a sequence of interpolants $s_n^{(\vartheta)}$ based on the power kernel $k^{(\theta)}$ and reuse similar ideas as before: \\
Maybe one could immediately use the power kernel $s_n^{(\theta=2)}$ and measure the approximation rate in some weaker norm, i.e.\
\begin{align}
\label{eq:idea_how_to_estimate}
&\Vert s_{X_n}^{(\theta = 2)} - s_{X_{n+1}}^{(\theta = 2)} \Vert_{\Ht} \notag \\
\leq~& \Vert s_{X_n}^{(\theta = 2)} - s_{X_{n+1}}^{(\theta = 2)} \Vert_{L^2(\Omega)}^{\textcolor{red}{?}} \cdot \Vert s_{X_n}^{(\theta = 2)} - s_{X_{n+1}}^{(\theta = 2)} \Vert_{\mathcal{H}_{\theta = 2}}^{\textcolor{red}{?}}.
\end{align}
However then the question is how to estimate the two quantities on the right hand side.
I hoped to be able to use arguments similar to \textit{escaping the native space}, because this is what is actually happening here: 
$k^{(\theta = 2)}$ is a very smooth kernel, and we want to show that $f \in \Ht$ for some $\vartheta \in (1, 2)$ -- i.e.\ this approach sounds promising.
However I was not able to work out estimates on this.
See 02.01.23\textbackslash A and 18.01.23\textbackslash A for some approaches. \\
Latest numerical examples indicate that a convergence order using a kernel $k$ does not transfer to the same convergence order when using the convolutional kernel $k^{(2)}$, see (Section 2.5 of) kernel\_approx/ 56\_example\_wendland.
\item However even though it might be possible to derive convergence rates on Eq.\ \eqref{eq:idea_how_to_estimate} based on estimates on
\begin{align*}
\Vert f  - s_{X_n} \Vert_{L^2(\Omega)} = \Vert f  - s_{X_n}^{(1)} \Vert_{L^2(\Omega)} \leq C h_{X_n}^\mu,
\end{align*}
it will be challenging to conclude the remaining proof:
Stability bounds (i.e.\ bounds on $\Vert A^{-1} \Vert_{2,2}$ via lower bounds on the smallest eigenvalue of the kernel matrix $A$) seem to be only only available for RBF kernels, however the power kernel is likely not RBF anymore!
Therefore as I'm unaware of any results that lower bound the smallest eigenvalue of the kernel matrix of non-RBF kernels, I have no good clue how to proceed here. 
Deriving lower bounds on the smalles eigenvalue of the kernel matrix for power kernels might be a research topic in its own.
\item Furthermore I'm not even sure whether the whole proof strategy really works:
I'm not sure whether asymptotically uniformly distributed points are the correct choise of point to approximate function in the power spaces $\Ht$, $\vartheta > 1$ 
-- or whether it might be more beneficial in this case to place more points at the boundary (oversampling near the boundary).
\end{itemize}
\item As a workaround, one could argument to just use a smoother kernel right from the beginning -- but not the power kernel.
This is nice from the practical point of view because we can still derive results:
\begin{itemize}
\item Assume a kernel $k$ such that $\ns \asymp H^\tau(\Omega)$, i.e.\ we are in the Sobolev setting.
Then $TL^2(\Omega)$ will not be $H^{2\tau}(\Omega)$, but some subset due to unclear boundary conditions. \\
Instead we could pick a kernel $k_2$ which is smoother, such that $\mathcal{H}_{k_2} \asymp H^{2\tau}(\Omega)$.
Using the same approach and theory above, we can show that a sufficient fast approximability of $f$ results in $f \in H^{\tilde{\tau}}$ for $\tau < \tilde{\tau} < 2\tau$.
This statement then still reveals the smoothness of $f$, but no longer the boundary conditions (which are anyway unclear). \\
The nice part is, that $f \in H^{\tilde{\tau}}$ for $\tau < \tilde{\tau} < 2\tau$ has some classical Sobolev smoothness which directly translates to $\vartheta < 1$ power spaces smoothness of $k_2$, but to not necessarily to $\vartheta > 1$ power space smoothness of $k$ (because there also boundary conditions contribute).
\item We already had to assume an interior cone condition for our technical statements to hold (see \Cref{th:utility_theorem_new} and \Cref{prop:disc_L2_from_cont_L2_new}),
 which is slightly more general than Lipschitz boundary (we could have interior cone condition, but violate an outerior cone condition, i.e.\ no Lipschitz boundary). \\
But if we additionally assume Lipschitz boundary, then Sobolev kernels yield RKHS that are norm equivalent to Sobolev spaces.
\end{itemize}
but a bit unsatisfacory because we do not stick to the general RKHS theory anymore and make use of Sobolev equivalences.
\end{itemize}


\subsection{On $L^\infty(\Omega)$ convergence rates}
\label{subsec:L2_conv_rates_sharp}

\begin{center}
\textbf{\textcolor{red}{\tw{I DID NOT CONTINUE FROM HERE ON!}}}
\end{center}

Based on Duchons localization principle \cite{duchon1978erreur} it is in general possible to show that the $\Vert \cdot \Vert_{L^2(\Omega)}$ error decays faster by a rate of $d/2$ than the $\Vert \cdot \Vert_{L^\infty(\Omega)}$ error \cite{light1998power, wendland1997sobolev}. \\

\subsection{Connect $L^\infty$ and $L^2$ error}
\label{subsec:connect_Linfty_L2}

\tw{
\begin{itemize}
\item I did some experiments using the Wendland $k=0$ kernel and $0.5 + x - x^2$ as target function, which allows for superconvergence.
I was quite surprised to observe a convergence rate as $n^{-2}$, because this is the maximal possible one (lower bound for linear interpolation, see \cite{wenzel2023analysis} and thus as fast as the $\Vert \cdot \Vert_{L^2(\Omega)}$ error rate). \\
One possible explanation might be that I squeezed out all the remaining convergence terms from the $\Vert f - s_{f, X_n} \Vert_{\ns}$ part, which is the required one for Duchon's localization principle.
\end{itemize}
}

Now I want to assume a lower bound on the $L^2$ error and then deduce a slower lower bound on the $L^\infty$ error (i.e.\ a slower rate of convergence).

\begin{lemma}
\label{lem:faster_L2_convergence}
Consider $f \in \ns$. Assume
\begin{align*}
\exists_{C>0} \forall_{X \subset \Omega, h_X/q_X \leq c_0} \Vert f - \Pi_X(f) \Vert_{L^\infty(\Omega)} \leq C h_X^\mu \cdot \Vert f - \Pi_X(f) \Vert_{\ns}.
\end{align*}
Then
\begin{align*}
\exists_{C > 0} \forall_{X \subset \Omega, h_X/q_X \leq c_0} \Vert f - \Pi_{X}(f) \Vert_{L^2(\Omega)} \leq C h_X^{\mu + d/2} \Vert f - \Pi_{X}(f) \Vert_{\ns}.
\end{align*}
\end{lemma}

\begin{proof}
This is pretty straight forward by using \cite{wendland1999meshless}.
\end{proof}

\noindent \Cref{lem:faster_L2_convergence} can also be stated in the reverse way:

\begin{lemma}
\label{lem:bounded_Linfty_convergence}
Consider $f \in \ns$.
Assume
\begin{align}
\label{eq:lower_L2_bound}
\exists_{c>0} \exists_{\mu > 0} \forall_{X \subset \Omega, h_X / q_X \leq c_0} \Vert f - \Pi_{X}(f) \Vert_{L^2(\Omega)} \geq c h_X^{\mu} \Vert f - \Pi_{X}(f) \Vert_{\ns}.
\end{align}
Then
\begin{align}
\label{eq:no_upper_Linfty_bound}
\forall_{\epsilon > 0} \nexists_{C > 0} \forall_{X \subset \Omega, h_X / q_X \leq c_0} \Vert f - \Pi_{X}(f) \Vert_{L^\infty(\Omega)} \leq C h_X^{\mu - d/2 + \epsilon} \Vert f - \Pi_{X}(f) \Vert_{\ns}.
\end{align}
%
%
%
%
%
%
%
%
\end{lemma}

\begin{proof}
Proof via contradiction:
Assume there exists an $\epsilon > 0$ and a $C > 0$ such that Eq.\ \eqref{eq:no_upper_Linfty_bound} holds.
Then leveraging \Cref{lem:faster_L2_convergence} we can conclude an additional $L^2$ decay rate based on that $L^\infty$ decay rate, i.e.\ $\mu - d/2 + \epsilon + d/2 = \mu + \epsilon$.
However this contradicts the assumption from \Cref{eq:lower_L2_bound}.
\end{proof}

Some remarks on both \Cref{lem:faster_L2_convergence} and \Cref{lem:bounded_Linfty_convergence}:
\begin{itemize}
\item We do not need the $\ns$-norm on the right hand side withing \Cref{eq:lower_L2_bound} and \Cref{eq:no_upper_Linfty_bound}, also any power space norm (as long as finite) should do the job! 
The only crucial ingredient is that the proof of the additional $L^2(\Omega)$ convergence rate works, see \cite{wendland1999meshless}.
\item Maybe we can formulate Lemma \ref{lem:bounded_Linfty_convergence} also in a positive way by considering $\sup_{X \subset \Omega, h_X / q_X \geq c_0} \Vert f - \Pi_X(f) \Vert_{L^\infty}$ or something!
I think I did this by also havin gLemma \Cref{lem:faster_L2_convergence}.
\item \textcolor{red}{Interesting observation:}
If we take the ratio of Eq.\ \eqref{eq:lower_L2_bound} and \eqref{eq:no_upper_Linfty_bound}, then the $\Vert \cdot \Vert_{\ns}$ part on the right hand side cancels.
\end{itemize}



\subsection{Approaching the $d/2$ gap}

Now the main idea is to connect the results of the previous subsection:
\begin{enumerate}
\item Assume we are given an $L^\infty(\Omega)$ decay rate.
Then we can leverage \Cref{lem:faster_L2_convergence} to obtain a faster $L^2(\Omega)$ decay rate. 
\item Using the results of \Cref{subsec:L2_conv_rates_sharp} we can conclude in which particular power space $f$ is contained.
\end{enumerate}

\tw{The main obstacle is that the previous subsection frequently makes use of $\Vert f - \Pi_X(f) \Vert_{\ns}$ expressions, which however do not appear originally in the assumption of Theorem 6.1 from \cite{schaback2002inverse}.}






~ \newpage
~ \newpage










\section{$d/2$ gap}

\subsection{Introduction} \label{subsec:introduction}

Robert Schaback mentioned the $d/2$ gap from time to time, but I never really understood what it is about.
In October 2022, Gabriele suggested \cite{karvonen2022error} to me, which mentioned the $d/2$ gap and referenced \cite{schaback2002inverse}, which introduces the notion $d/2$ gap after the proof of its Theorem 6.1: \\
\cite{schaback2002inverse} considers kernels $k$ with RBF $\Phi$ with Fourier transform $\hat{\Phi}$ satisfying
\begin{align}
\label{eq:decay_condition}
c_1 (1+\Vert \omega \Vert_2)^{-d-s_\infty} \leq \hat{\Phi}(\omega) \leq c_2(1+\Vert \omega \Vert_2)^{-d-s_\infty}, \quad 0 < c_1 \leq c_2
\end{align}
\textit{Direct estimates} bound the error for functions for functions $f \in \ns$  in the RKHS as
\begin{align}
\label{eq:direct_estimate}
\Vert f - s_{f, X} \Vert_{L^\infty(\Omega)} \leq C_\Phi h^{s_\infty/2} \Vert f - s_{f, X} \Vert_{\ns}.
\end{align}
\textit{Inverse estimates} state that a sufficiently fast error decay implies that the function is included in the RKHS. 
\cite[Theorem 6.1]{schaback2002inverse} states (with notation adopted to here):
\begin{theorem}
%
Let $\Omega \subset \R^d$ be a bounded and open domain satisfying an interior cone condition. 
The positive definite basis function $\Phi$ should satisfy the decay condition \ref{eq:decay_condition}. 
Suppose further that for some $f \in C(\Omega)$ there exists constants $\mu > 0$ and $c_f > 0$ such that $\Vert f - s_{f, X}\Vert_{L^\infty(\Omega)} \leq c_f h_X^\mu$ for all $X \subset \Omega$ with $h_X$ sufficiently small. 
If $2\mu > s_\infty + d$, then $f$ must belong to the native space $\ns$.
\end{theorem}
The $d/2$ gap occurs, because Theorem \ref{th:inverse_statement} requires a decay of $\mu > s_\infty/2 + d/2$, while Eq.\ \ref{eq:direct_estimate} only gives $s_\infty/2$.
We note it holds that $h^{d/2} = n^{-1/2}$ for asymptotically uniformly distributed points.

\subsection{Analysis of proof}

When analyzing the proof of Theorem \ref{th:inverse_statement}, one observe that one estimate within the second block of estimates is not necessarily sharp.
For convenience, we include this block of estimates here:
\begin{align}
\label{eq:nonsharp_estimate}
\Vert s_{f, Y} - s_{f, X} \Vert_{\ns}^2 &\leq \Vert A_{Y, \Phi}^{-1} \Vert_{2,2} \Vert s_{f, Y} - s_{f, X} \Vert_{L^2(Y)}^2 \notag \\
&\leq \frac{1}{\gamma_Y} \sum_{y \in Y} |f(y) - s_{f, X}(y)|^2 \notag \\
&\leq c_\Phi^{-1} q_Y^{-s_\infty} |Y| c_f^2 h_X^{2\mu},
\end{align}
where $\frac{1}{\gamma_Y}$ is a sharp upper bound for the norm of the inverse matrix.
Running several numerical experiments, we observe that the step in Eq.\ \eqref{eq:nonsharp_estimate} is not necessarily sharp (namely for $f$ with low smoothness, e.g.\ $x^\alpha \in H^1([0, 1])$, $\alpha = 0.51$):
Here a (discrete) $L^2$ norm was estimated with help of an $L^\infty$ estimate (namely $\Vert f - s_{f, X}\Vert_{L^\infty(\Omega)} \leq c_f h_X^\mu$ as stated as assumption in Theorem \ref{th:inverse_statement}.
This is a straightforward estimate, but therefore not necessarily sharp\footnote{Interestingly, this estimate seems to be sharp for functions $f$ with higher smoothness, e.g.\ $x^{2} \in H^1([0, 1])$}: 
When recalling sampling inequalities, $L^2(\Omega)$ estimates give a faster rate of decay than $L^\infty(\Omega)$ estimates --- namely an additional factor of $h^{d/2}$, which coincides exactly with the observed gap of $d/2$.

\subsection{Overview of approaches how to close the gap}

\noindent I tried several approaches to close this gap, i.e.\ estimating Eq.\ \eqref{eq:nonsharp_estimate} in a better way:
\begin{enumerate}
\item One could try to modify the estimates of the first block of inequalities within the proof of \cite[Theorem 6.1]{schaback2002inverse} (which is not stated here) in order to obtain $\Vert s_{f, Y} - s_{f, X} \Vert_{L^\infty(Y)}^2$ instead of $\Vert s_{f, Y} - s_{f, X} \Vert_{L^2(Y)}^2$.
The difficulty here is that this yields to different operator norms like $\Vert A_{Y, \Phi}^{-1} \Vert_{\mathcal{L}(L^\infty(Y) \rightarrow L^1(Y))}$ (03.11.22\textbackslash A3) or even different operators (03.11.22\textbackslash A2, 02.11.22\textbackslash C1) which are difficult to estimate.
\item I tried to estimate $\Vert s_{f, Y} - s_{f, X} \Vert_{L^2(Y)}^2$ via $\Vert s_{f, Y} - s_{f, X} \Vert_{L^2(\Omega)}^2$, but could not find a connection.
\item I tried to use the additional decay of $n^{-1/2}$ due to $f$-greedy, however then possibly we loose the asymptotically uniform distirbution of points, which is required e.g.\ to upper bound $\Vert A_{Y, \Phi}^{-1} \Vert_{2,2}$
\item I tried to estimate $\Vert s_{f, Y} - s_{f, X} \Vert_{L^2(Y)}^2$ via $\Vert P_X \Vert_{L^2(Y)}^2 \cdot \Vert s_{f, Y} - s_{f, X} \Vert_{\ns}^2$, i.e.\ by using the standard power function estimate. 
However $\Vert P_X \Vert_{L^2(Y)}^2$ seems to decay only with the same rate as $\Vert P_X \Vert_{L^\infty(Y)}^2$, thus not providing any benefit.
\item I tried to understand why $\Vert f - s_{f, X} \Vert_{L^2(\Omega)}^2$ decays faster than $\Vert f - s_{f, X} \Vert_{L^2(\Omega)}^2$: 
Therefore I analyzed \cite[Theorem 5]{wendland1997sobolev}, which extends ideas from \textit{Duchon's localization trick} \cite{duchon1978erreur} (french paper). 
This idea worked out best and will be elaborated on in Subsection \ref{subsec:on_closing_the_gap}.
\end{enumerate}


\subsection{On closing the $d/2$ gap} \label{subsec:on_closing_the_gap}

I analyzed \cite[Theorem 5]{wendland1997sobolev} and its proof in detail, and I think it can be modified to obtain the following theorem:
\tw{The following thoughts can probably just found in Subsection \textit{2.2.2. Sobolev bounds on discrete q norms} of \cite{narcowich2005sobolev}}

\begin{theorem}[Version of Theorem 5 from Wendland paper]
\label{th:modified_theorem}
Let $\Omega$ be an open and bounded subset of $\R^d$ having the cone property and a Lipschitz boundary.
Denote by $s_f$ the interpolant to $f \in W_2^m(\R^d)$ with $m = \lceil s \rceil$.
Let $s$ and $h$ be defined as in \cite[Theorem 1]{wendland1997sobolev}.
Then we have for sufficiently small $h$ and $p \geq 2$:
\begin{align}
\label{eq:modified_theorem}
\frac{1}{|T_h|} \sum_{t \in T_h} | (f-s_f)(t)|^p \leq \frac{1}{|T_h|} \Vert P_{X} \Vert_{L^\infty(\Omega)}^p \Vert f - s_{f, X} \Vert_{W_2^m(\Omega)}^p \\
\stackrel{p=2}{\Rightarrow} \frac{1}{|T_h|} \sum_{t \in T_h} | (f-s_f)(t)|^2 \leq \frac{1}{|T_h|} \Vert P_{X} \Vert_{L^\infty(\Omega)}^2 \Vert f - s_{f, X} \Vert_{W_2^m(\Omega)}^2 \notag
\end{align}
\end{theorem}

\tw{I don't want to have the Power function here, modify this!}

\begin{proof}
The proof is quite close to the original proof: 
Instead of considering the integral and decomposing it into the sum $\sum_{t \in T_h}$, we leave the integral out and directly consider the sum:
\begin{align*}
\sum_{t \in T_h} |(f - s_{f, X})(t)|^p &\leq \sum_{t \in T_h} C P_X^p(t) \cdot \Vert (f-s_f)^{B(t, Mh)} \Vert_{W_2^m(\R^d)}^p \\
&\leq C \Vert P_X \Vert_{L^\infty(\Omega)}^p \cdot \sum_{t \in T_h} \Vert (f-s_f)^{B(t, Mh)} \Vert_{W_2^m(\R^d)}^p.
\end{align*}
For the $\sum_{t \in T_h}$ part we use exactly the same calculation as in \cite{wendland1997sobolev} which shows
\begin{align*}
\sum_{t \in T_h} \Vert (f-s_f)^{B(t, Mh)} \Vert_{W_2^m(\R^d)}^p \leq \left( \sum_{|\alpha| \leq m} \int_{\R^d} M_1 |D^\alpha(f-s_{f, X})|^2 \mathrm{d}x \right)^{p/2},
\end{align*}
such that we obtain
\begin{align*}
\sum_{t \in T_h} |(f - s_{f, X})(t)|^p &\leq C \Vert P_X \Vert_{L^\infty(\Omega)}^p \Vert f - s_{f, X} \Vert_{W_2^m(\R^d)}^p \\
&\leq C' \Vert P_X \Vert_{L^\infty(\Omega)}^p \Vert f - s_{f, X} \Vert_{W_2^m(\Omega)}^p.
\end{align*}
Dividing this result by $|T_h|$ gives the final statement.
\end{proof}

\begin{rem}
Using\footnote{\cite[Lemma 2]{wendland1997sobolev} only mentions $|T_h| \leq C h^{-d}$. 
However \cite[Lemma 2]{wendland1997sobolev} (ii) states $\Omega \subset \cup_{t \in T_h} B(t, Mh)$, such that a volume comparison gives $|\Omega| \leq |T_h| \cdot c_{M,d} h^d$, which gives $|T_h| \geq |\Omega| c_{M,d}^{-1} h^{-d}$, i.e.\ we have the equivalence $|T_h| \asymp h^{-d}$} $|T_h| \asymp h^{-d}$
turns Eq.\ \eqref{eq:modified_theorem} into (after taking the $p$th-root)
\begin{align}
\label{eq:result_with_prefactor}
\left( \frac{1}{|T_h|} \sum_{t \in T_h} | (f-s_f)(t)|^p \right)^{1/p} \leq h_X^{d/p} \Vert P_{X} \Vert_{L^\infty(\Omega)} \Vert f - s_{f, X} \Vert_{W_2^m(\Omega)}.
\end{align}
The left hand side is kind of a discretization of the integral $\Vert f - s_f \Vert_{L^p(\Omega)}$, i.e.\ the Eq.\ \eqref{eq:result_with_prefactor} looks like a discrete version of the statement \cite[Theorem 5]{wendland1997sobolev} (after replacing $\Vert P_X \Vert_{L^\infty(\Omega)}$ by $h^{k + 1/2}$).
\end{rem}

\begin{rem}
\label{rem:how_to_use_theorem}
Some more remarks:
\begin{itemize}
\item Theorem \ref{th:modified_theorem} is not directly applicable to solve the $d/2$ gap (i.e.\ to better estimate the step within Eq.\ \eqref{eq:nonsharp_estimate}) but close:
\begin{itemize}
\item Theorem \ref{th:modified_theorem} requires the assumption that our function $f$ is Sobolev. 
However this is exactly that the inverse Theorem \ref{th:inverse_statement} wants to show.
One relief might be that we want to apply Theorem \ref{th:modified_theorem} to $s_{f,Y} - s_{f,X}$, which are both sums of kernels and therefore Sobolev.
\item One further problem might be, that the function which we consider changes "all the time", because $s_{f,Y}$ is updated.
The original proof within \cite{wendland1997sobolev} circumvented this by using $s_{f, Y}(y) = f(y)$ for all $y \in Y$.
In order to use this trick as well, we need to remove all Sobolev related assumptions on $f$ within Theorem \ref{th:modified_theorem}, which might be difficult because we heavily rely on restriction and extension statements.
\item The relation between $t \in T_h$ and $y \in Y$ is a bit unclear.
I hoped that we could use $t = y \in Y$, which might indeed work because $Y$ is asymptotically uniformly distributed (and this is what I expect $T_h$ to be) - but this needs to be checked.
\end{itemize}
\item I do not really understand \textit{why} Duchons localization principle works.
In my opinion it should also be applicable to the power function, thus deriving a faster rate of decay for $\Vert P_n \Vert_{L^2(\Omega)}$ than for $\Vert P_n \Vert_{L^\infty(\Omega)}$ (which however does not hold numerically).
The key step needs to be somewhere hidden in the use of the restriction and extension statements.
\end{itemize}
\end{rem}

\noindent Based on these Remarks \ref{rem:how_to_use_theorem} we pose the following conjecture as a follow up to Theorem \ref{th:modified_theorem}:

\begin{conj}
\label{conj:conjecture1}
Using the same assumptions as Theorem \ref{th:inverse_statement}, we have for uniformly distributed points $Y \subset \Omega$ with $h_Y \leq h_X / 2$:
\begin{align}
\label{eq:conjectured_eq}
\frac{1}{|Y|} \sum_{y \in Y} | (f-s_{f,X})(y)|^p \leq \frac{1}{|X|} \Vert P_{X} \Vert_{L^\infty(\Omega)}^p \Vert f - s_{f, X} \Vert_{W_2^m(\Omega)}^p
\end{align}
\end{conj}

\begin{rem}
Note that we are interested in the case $|Y| \approx 2^d|X|$, i.e.\ $|Y| \asymp |X|$.
This could be seen as a kind of \textit{average error} (in contrast to a worst case error) which Robert Schaback sometimes mentioned.
\end{rem}

\noindent The following Lemma shows how to close the $d/2$ gap based on Conjecture \ref{conj:conjecture1}:

\begin{lemma}
\label{lem:d_half_gap_closed}
Same setting as Theorem \ref{th:inverse_statement}.
Suppose further that for some $f \in C(\Omega)$ there exists constants $\mu > 0$ and $c_f > 0$ such that $\Vert f - s_{f, X}\Vert_{L^\infty(\Omega)} \leq c_f h_X^\mu$ for all $X \subset \Omega$ with $h_X$ sufficiently small. 
If $\mu > s_\infty / 2$, then $f$ must belong to the native space $\ns$.
\end{lemma}

\begin{proof}
We use mainly the same proof as in \cite[Theorem 6.1]{wendland1997sobolev}, but we make use of Eq.\ \eqref{eq:conjectured_eq} (for $p=2$) to better estimate the step from Eq.\ \eqref{eq:nonsharp_estimate}:
\begin{align*}
%
\Vert s_{f, Y} - s_{f, X} \Vert_{\ns}^2 &\leq c_\Phi^{-1} q_Y^{-s_\infty} \sum_{y \in Y} |f(y) - s_{f, X}(y)|^2 \notag \\
&\leq c_\Phi^{-1} q_Y^{-s_\infty} \Vert P_{X} \Vert_{L^\infty(\Omega)}^2 \Vert f - s_{f, X} \Vert_{W_2^m(\Omega)}^2
\end{align*}
%
%
Now we quickly have to cheat a bit because Theorem \ref{th:modified_theorem} and Conjecture \ref{conj:conjecture1} were stated using the power function. 
However I should state them using the assumption $\Vert f - s_{f, X}\Vert_{L^\infty(\Omega)} \leq c_f h_X^\mu$, thus replacing $\Vert P_{X} \Vert_{L^\infty(\Omega)}^2$ by $c_f^2 h_X^{2\mu}$:
\begin{align*}
%
\Vert s_{f, Y} - s_{f, X} \Vert_{\ns}^2 &\leq c_\Phi^{-1} c_f^2 q_Y^{-s_\infty} h_X^{2\mu}
\end{align*}
Now we use (following again \cite[Theorem 6.1]{wendland1997sobolev}) $Y = X_{n+1}, X = X_n$ with asymptotically uniformly distributed points, such that $q_Y \asymp h_X$.
Since $\mu > s_\infty/2$ we have $2\mu - s_\infty > 0$ and thus obtain a decay. 
Now we can reuse the remaining arguments on \cite[Proof of Theorem 6.1]{wendland1997sobolev} to conclude convergence of $s_{f, X}$ to an element which must be $f$.
\end{proof}



\section{Further remarks}

\begin{itemize}
\item Crucial inside: We fix a function, and then we interpolate (with denser set of points). 
If we consider for example the power function, we do \textit{not} observe a faster rate of decay for $\Vert P_n \Vert_{L^2(\Omega)}$ than for $\Vert P_n \Vert_{L^\infty(\Omega)}$.
\item The $f$-greedy analysis provides also an additional factor of $n^{-1/2} = h^{d/2}$, which is however in my opinion not related at all to this $d/2$ gap here:
The $d/2$ gap here is about asymptotically uniformly distributed points. 
The $f$-greedy however usually gives non-uniformly distributed points.
\item I did not further check the implications for the formulas and thoughts presented in \cite{karvonen2022error}.
\end{itemize}

\section{Numerical experiments}

For simplicity I did only numerical experiments in 1D so far: 
I used the Brownian bridge kernel $k(x, y) = \min(x, y) - xy$ and considered the family of functions
\begin{align*}
f_\alpha(x) = x^\alpha \cdot (1-x)^\alpha,
\end{align*}
which satisfies $f_\alpha(0) = f_\alpha(1) = 0$ for $\alpha > 0$. \\

\noindent From kernel\_approx/37\_brownian\_bridge\_example we recall:

\begin{itemize}
\item $\alpha > -1/2$: $f_\alpha \in L^2(\Omega)$, 
\item $\alpha > 1/2$: $f_\alpha \in \ns$, 
\item $\alpha > 3/2$: $f_\alpha \in TL^2(\Omega)$. 
\item Special case: $\alpha = 1$: $f \in TL^2(\Omega)$, especially $f_1 = \int_\Omega k(\cdot, y) 1 ~ \mathrm{d}y$
\end{itemize}

\noindent For the following approximation experiments I used $2^n + 1$ equidistant points $X_n$ for $n \in \{2, 3, .., 10\}$. \\
I investigated the approximation of the functions $f_\alpha$ for several $\alpha > 0$ values on the sequence of denser sets $\{ X_n \}_{n \in \{2, 3, .., 10\}}$ using several metrics.
I plotted the following quantities:
\begin{enumerate}
\item \textit{Discrete $L^\infty$ error}:
\begin{align}
\label{eq:discrete_Linfty_error}
\max_{y \in X_{n+1}} | s_{f,X_{n+1}}(y) - s_{f,X_n}(y) |
\end{align}
\item \textit{Discrete $L^2$ error}:
\begin{align}
\label{eq:discrete_L2_error}
\Vert s_{f, X_{n+1}} - s_{f, X_n} \Vert_{L^2(X_{n+1})} = \sqrt{\sum_{y \in X_{n+1}} | s_{f,X_{n+1}}(y) - s_{f,X_n}(y) |^2}
\end{align}
\item 
Squared RKHS norm of subsequent differences
\begin{align}
\label{eq:squared_rkhs_norm_diff}
\Vert s_{f, X_{n+1}} - s_{f, X_n} \Vert_{\ns}^2,
\end{align}
which is upper estimated according to Eq.\ \eqref{eq:nonsharp_estimate}
\begin{align}
\label{eq:upper_estimate}
\Vert A_{X_{n+1}}^{-1} \Vert \cdot \Vert s_{f, X_{n+1}} - s_{f, X_n} \Vert_{L^2(X_{n+1})}^2
\end{align}

\item Ratio of previous quantities:
\begin{align}
\label{eq:ratio_quantities}
\frac{\sum_{y \in X_{n+1}} | s_{f,X_{n+1}}(y) - s_{f,X_n}(y) |^2}{|X_{n+1}| \cdot \max_{y \in X_{n+1}} | s_{f,X_{n+1}}(y) - s_{f,X_n}(y) |} 
%
\end{align}
\end{enumerate}
The results are depicted in Figure \ref{fig:brownian_bridge_example}: 
We make the following observations:
\begin{enumerate}
\item Figure \ref{fig:brownian_bridge_example} bottom left:
Estimating Eq.\ \eqref{eq:squared_rkhs_norm_diff} by Eq.\ \eqref{eq:upper_estimate} seems to be a sharp estimate, the dashed line behaves all the time according to the solid line.
\item Figure \ref{fig:brownian_bridge_example} bottom right:
For rougher functions (i.e.\ smaller values of $\alpha$, i.e. blue colors) the ratio from Eq.\ \eqref{eq:ratio_quantities} approaches zero, which means that the numerator and the denominator do not follow the same rate of convergence.
This elucidates that the last step within Eq.\ \eqref{eq:nonsharp_estimate} is not sharp for rougher functions!
\end{enumerate}






\section{Gaussian kernel}

Talking to Jens Wirth, he pointed me to \cite[Chapter 2.1]{butzer2013semi} and I try to approximately summarize some ideas:
Similar to Subsection \ref{subsec:introduction}, where the notions of \textit{direct estimates} and \textit{inverse estimates} were introduced,
\cite{butzer2013semi} discusses \textit{direct theorems} and \textit{converse theorems} for family of operators $\{ T(t); 0 < t < \infty \}$ converging to the identity operator for $t \rightarrow 0+$.
While there is no obvious direct connections, the notions are very similar and its probable that the concepts are related:
Especially \cite[Definition 2.1.1]{butzer2013semi} introduces \textit{saturation or Favard classes}:
Those are classes of functions, where we have an equivalence of of direct and converse theorems. \\
Their main theorem is a saturation theorem for semi-groups of operators and Jens remarked a connection to my setting:
If we consider the Laplace operator and the corresponding semi-group, we end up with something like $\sum_{k=0}^\infty (-\Delta)^k / k!$, which is related to the Gaussian kernel. \\
Finally due to some results within \cite{butzer2013semi} he presumed that while the might not be a $d/2$ gap for Sobolev kernels, the situation might be different for the Gaussian (or other analytic kernels). \\
And here I agree: If we check the proof of \cite[Theorem 6.1]{schaback2002inverse} and try to run it for the Gaussian kernel, 
we end up looking at Eq.\ \eqref{eq:nonsharp_estimate}:
\begin{align*}
\Vert s_{f, Y} - s_{f, X} \Vert_{\ns}^2 &\leq \Vert A_{Y, \Phi}^{-1} \Vert_{2,2} \Vert s_{f, Y} - s_{f, X} \Vert_{L^2(Y)}^2 \notag \\
&\leq \frac{1}{\gamma_Y} \sum_{y \in Y} |f(y) - s_{f, X}(y)|^2 \notag
\end{align*}
However for the Gaussian kernel we encounter a quite big gap here:
\begin{itemize}
\item According to \cite[Corollary 12.4]{wendland2005scattered}, the smallest eigenvalue for the kernel matrix of the Gaussian kernel $e^{-\alpha \Vert x - y \Vert_2^2}$ is bounded as
\begin{align*}
\lambda_{\min}(A_X) \geq C_d(2\alpha)^{-d/2} e^{-40.61d^2 / (q_X^2 \alpha)} q_X^{-d},
\end{align*}
\item while the error decays slower: 
Instead of having $q_X^{-2}$ in the exponential for the lower bound of the smallest eigenvalue, 
we likely only have $h_X^{-1}$ (and probably some $\log(h_X)$ term) in the exponential of the upper bound on the error (see Theorem 3 within kernel\_approx/32\_conv\_speed\_gaussian or \cite[Theorem 1.2]{karvonen2022approximation} for 1D for $k(x, y) = \exp(-0.5\epsilon^2 (x-y)^2)$):
\begin{align*}
c_{1,2} \left( \frac{\epsilon}{2} \right)^n (n!)^{-1/2} &\leq \inf_{A_n} \sup_{0 \neq f \in \ns} \frac{ \Vert f - A_n f \Vert_{L^2([-1,1])}}{\Vert f \Vert_{\ns}} \\
&\leq 2 \sqrt{2c_L(\epsilon^2)} n^{-1/8} e^{\epsilon\sqrt{n}} \left( \frac{\epsilon}{2} \right)^n (n!)^{-1/2},
\end{align*}
where $A_n$ is a linear approximation using standard information.
Recalling
\begin{align*}
\sqrt{2\pi} n^{n + 1/2} e^{-n} < n! < \sqrt{2\pi} n^{n + 1/2} e^{-n+1}
\end{align*}
we can see that the error decay has only the rate $e^{-\log(n)n}$ (corresponding to $e^{\log(h)h}$ instead of $e^{-n^2}$ (corresponding to $e^{h^{-2}}$).
\end{itemize}
Alltogether this means that the proof strategy of \cite[Theorem 6.1]{schaback2002inverse} via Eq.\ \eqref{eq:nonsharp_estimate} does not yield a convergence of $\Vert s_{f,Y} - s_{f, X} \Vert_{\ns}$ for functions with a worst case optimal decay rate.
I.e.\ either the proof strategy of \cite[Theorem 6.1]{schaback2002inverse} does not work for Gaussian kernels (and probably other analytic kernels) or there is indeed a (huge) gap for the Gaussian kernel. \\
It might be interesting to check whether \cite{karvonen2022approximation} also has some statements on the lowest eigenvalue of the kernel matrix (or whether it is possible to derive some based on those techniques).

Maybe we can see here some connection between \textit{closing the $d/2$ gap} but having an \textit{analytic gap}.

\subsection{Numerical results}

In order to get a better feeling for this, I wanted to run similar experiments as done in Figure \ref{fig:brownian_bridge_example} -- however due to numerical issues this did not work out properly.




\section{Ideas}

\begin{itemize}
\item We could replace $f$ by $s_N, \Vert f - s_N \Vert_{L^\infty(\Omega)} < \epsilon$ and then split
\begin{align*}
\Vert f - s_n \Vert_{L^\infty(\Omega)} &\leq \Vert f - s_N \Vert_{L^\infty(\Omega)} + \Vert s_N - s_n \Vert_{L^\infty(\Omega)} \\
&\leq \epsilon + \Vert s_N - s_n \Vert_{L^\infty(\Omega)} 
\end{align*}
However this does not give us much, as we cannot apply the assumption on the approximability of $f$ anymore.
\item Does the discrete $L^2$ error also decay faster (by a rate of $h^{d/2}$) than the $L^\infty$ error for functions outside the RKHS?
This might shed light onto where to look for the missing term to come from.
\item asdf
\item What happens, if we directly assume a bound on the $\Vert \cdot \Vert_{L^2(\Omega)}$ error instead of on the $\Vert \cdot \Vert_{L^\infty(\Omega)}$ error?
Then we only somehow need to go for the discrete $L^2(Y)$ error, but then the proof should work out!
\end{itemize}

\begin{align*}
\left \Vert \sum_{i=1}^n \alpha_i k(\cdot, x_i) \right \Vert_{\Ht}, \qquad 0 < \vartheta < 1
\end{align*}

~ \newpage 

\section{Idea from 24.11.22}

\begin{theorem}
\label{th:inverse_statement_improved}
Let $\Omega \subset \R^d$ be a bounded and open domain satisfying an interior cone condition. 
The positive definite basis function $\Phi$ should satisfy the decay condition \ref{eq:decay_condition}. 
Suppose further that for some $f \in C(\Omega)$ there exists constants $\mu > 0$ and $c_f > 0$ such that $\Vert f - s_{f, X}\Vert_{L^\infty(\Omega)} \leq c_f h_X^\mu$ for all $X \subset \Omega$ with $h_X$ sufficiently small. 
If $2\mu > s_\infty$, then $f$ must belong to the native space $\ns$.
\end{theorem}

From the assumption it follows directly that
\begin{align*}
\Vert f - s_{f, X}\Vert_{L^\infty(\Omega)} &\leq \frac{c_f}{\Vert f \Vert_{L^2(\Omega)}} h_X^\mu \cdot \Vert f \Vert_{L^2(\Omega)} \\ 
&=: c' h_X^\mu \cdot \Vert f \Vert_{L^2(\Omega)}
\end{align*}



\section{Idea from 18.11.22}


We need the following Corollary, which is a simplification of e.g.\ Corrollary 2 or 3 from \textit{Continuous superconvergence from Gabriele}:
\begin{cor}
%
Consider $\theta_1 < \theta < \theta_2$. Then under some assumptions it holds 
\begin{align}
%
\Vert g - \Pi(g) \Vert_{\mathcal{H}_{\theta}} &\leq \Vert g - \Pi(g) \Vert_{\mathcal{H}_{\theta_1}}^{\frac{\theta_2 - \theta}{\theta_2 - \theta_1}} \cdot \Vert g - \Pi(g) \Vert_{\mathcal{H}_{\theta_2}}^{\frac{\theta - \theta_1}{\theta_2 - \theta_1}} \notag \\
\Rightarrow \Vert s_{f, Y} - s_{f, X} \Vert_{\mathcal{H}_{\theta}} &\leq \Vert s_{f, Y} - s_{f, X} \Vert_{\mathcal{H}_{\theta_1}}^{\frac{\theta_2 - \theta}{\theta_2 - \theta_1}} \cdot \Vert s_{f, Y} - s_{f, X} \Vert_{\mathcal{H}_{\theta_2}}^{\frac{\theta - \theta_1}{\theta_2 - \theta_1}} \notag \\
&\leq \Vert s_{f, Y} - s_{f, X} \Vert_{L^2(\Omega)}^{1 - \theta} \cdot \Vert s_{f, Y} - s_{f, X} \Vert_{\ns}^{\theta}
\end{align}
\end{cor}

\begin{proof}
Can be proven elementary via Hölder inequality.
The conclusion is done by inserting $g = s_{f, Y}, \Pi(g) = s_{f, X}$ and $\theta_1 = 0, \theta_2 = 1$.
\end{proof}

Now we state and prove the following improvement of Theorem \ref{th:inverse_statement}, which replaced the assumption $2\mu > s_\infty + d$ by $2\mu > s_\infty$ and thus removed the $d/2$ gap.

\begin{theorem}
\label{th:inverse_statement_improved}
Let $\Omega \subset \R^d$ be a bounded and open domain satisfying an interior cone condition. 
The positive definite basis function $\Phi$ should satisfy the decay condition \ref{eq:decay_condition}. 
Suppose further that for some $f \in C(\Omega)$ there exists constants $\mu > 0$ and $c_f > 0$ such that $\Vert f - s_{f, X}\Vert_{L^\infty(\Omega)} \leq c_f h_X^\mu$ for all $X \subset \Omega$ with $h_X$ sufficiently small. 
If $2\mu > s_\infty$, then $f$ must belong to the native space $\ns$.
\end{theorem}

\tw{The proof idea is good in my opinion, however it does not work out :(}

\cite{schaback2002inverse} assumed additional convergence orders to show that the function is included in the original space.
Here we don't want to assume additional convergence orders, but instead show (as a first step) that the function in included in a less smooth space (in terms of interpolation spaces).
This then enables us to obtain faster error rates for the $L^2(X)$ error, which are faster by a rate of $h_X^{d/2}$ compared to the straightoforward derived rates based on the $L^\infty(\Omega)$ error.

\begin{proof}
We start similar to the proof in \cite{schaback2002inverse}, but consider a weaker norm on the left hand side, to which we apply Corollary \ref{cor:hoelder_interpol}:
\begin{align*}
\Vert s_{f, Y} - s_{f, X} \Vert_{\Ht}^2 &\leq \Vert s_{f, Y} - s_{f, X} \Vert_{L^2(\Omega)}^{2 - 2\vartheta} \cdot \Vert s_{f, Y} - s_{f, X} \Vert_{\ns}^{2\vartheta} \\
&\leq \Vert s_{f, Y} - s_{f, X} \Vert_{L^2(\Omega)}^{2 - 2\vartheta} \cdot \Vert A_{Y, \Phi}^{-1} \Vert_{2,2}^{\vartheta} \cdot \Vert s_{f, Y} - s_{f, X} \Vert_{L^2(Y)}^{2\vartheta}
\end{align*}
where we estimated the $\Vert s_{f, Y} - s_{f, X} \Vert_{\ns}$ norm just in the standard way as done in \cite{schaback2002inverse}.
Observe that for $\vartheta = 1$ we reobtain the standard statement.
The introduction of $0 < \vartheta < 1$ thus allowed us to slow down the growing factor $\Vert A_{Y, \Phi}^{-1} \Vert_{2,2}$,
while the two $L^2$ errors approximately stay the same because their exponents cancel: $(2 - 2\vartheta) + 2\vartheta = 2$.

Now we could try to estimate the $\Vert s_{f, Y} - s_{f, X} \Vert_{L^2(\Omega)}$ and $\Vert s_{f, Y} - s_{f, X} \Vert_{L^2(Y)}$ factor in a sophisticated way, but as a first step we estimate them in the straightforward way by using the $\Vert \cdot \Vert_{L^\infty(\Omega)}$ bound:
\begin{align*}
\Vert s_{f, Y} - s_{f, X} \Vert_{L^2(\Omega)} &\leq \sqrt{|\Omega|} \cdot \Vert s_{f, Y} - s_{f, X} \Vert_{L^\infty(\Omega)} \\
&\leq \sqrt{|\Omega|} \cdot \left( \Vert f - s_{f, Y} \Vert_{L^\infty(\Omega)} + \Vert f - s_{f, X} \Vert_{L^\infty(\Omega)} \right) \\ 
&\leq 2C \sqrt{|\Omega|} \cdot h_X^\mu \\
\Vert s_{f, Y} - s_{f, X} \Vert_{L^2(X)} &\leq \sqrt{|Y|} \cdot \Vert s_{f, Y} - s_{f, X} \Vert_{L^\infty(\Omega)} \\
&\leq \sqrt{|Y|} \cdot \left( \Vert f - s_{f, Y} \Vert_{L^\infty(\Omega)} + \Vert f - s_{f, X} \Vert_{L^\infty(\Omega)} \right) \\ 
&\leq 2C \sqrt{|Y|} \cdot h_X^\mu
\end{align*}
Additionally using $\Vert A_{Y, \Phi}^{-1} \Vert_{2,2} \leq c q_Y^{-s_\infty}$ we obtain alltogether (merging all constants as $C$):
\begin{align}
\begin{aligned}
\label{eq:important_estimate}
\Vert s_{f, Y} - s_{f, X} \Vert_{\Ht}^2 &\leq \Vert s_{f, Y} - s_{f, X} \Vert_{L^2(\Omega)}^{2 - 2\vartheta} \cdot \Vert A_{Y, \Phi}^{-1} \Vert_{2,2}^{\vartheta} \cdot \Vert s_{f, Y} - s_{f, X} \Vert_{L^2(Y)}^{2\vartheta} \\
&\leq C q_Y^{-s_\infty \vartheta} \cdot h_X^{(2-2\vartheta)\mu} \cdot h_X^{2\vartheta \mu} |Y|^\vartheta. 
\end{aligned}
\end{align}
Now --- as in \cite{schaback2002inverse} --- we consider $X_n$ such that $|X_n| \leq c 2^{nd}$ and
\begin{align*}
c_1 2^{-n} \leq q_{X_n} \leq h_{X_n} \leq c_2 2^{-n}.
\end{align*}
Using $X_{n+1} = Y \supset X = X_{n+1}$ we obtain:
\begin{align*}
\Vert s_{f, X_{n+1}} - s_{f, X_n} \Vert_{\Ht}^2 &\leq C q_{X_{n+1}}^{-s_\infty \vartheta} \cdot h_{X_{n}}^{(2-2\vartheta)\mu} \cdot h_{X_n}^{2\vartheta \mu} |X_{n+1}|^\vartheta \\
&\leq C 2^{s_\infty \vartheta (n+1)} \cdot 2^{-(2-2\vartheta)\mu n} \cdot 2^{-2\vartheta \mu n} 2^{(n+1)d \vartheta} \\
&\leq C 2^{s_\infty \vartheta (n+1) -(2-2\vartheta)\mu n -2\vartheta \mu n + (n+1)d \vartheta} \\
&\leq C 2^{s_\infty \vartheta n -(2-2\vartheta)\mu n -2\vartheta \mu n + n d \vartheta} \cdot 2^{s_\infty \vartheta + d \vartheta} \\
&\leq C 2^{-n(-s_\infty \vartheta + (2-2\vartheta)\mu + 2\vartheta \mu  - d \vartheta)} \cdot 2^{s_\infty \vartheta + d \vartheta} \\
&\leq C 2^{-n(-s_\infty \vartheta + 2\mu - d \vartheta)} \cdot 2^{s_\infty \vartheta + d \vartheta} \\
&\leq C 2^{-n(2\mu -(s_\infty + d) \vartheta)} \cdot 2^{s_\infty \vartheta + d \vartheta} \\
\end{align*}
Not that for $\vartheta = 1$ we reobtain the same expression as in \cite{schaback2002inverse}, i.e.\ we didn't make too many mistakes.
In order to use the telescoping sum argument from \cite{schaback2002inverse}, we need
\begin{align*}
0 \stackrel{!}{<} 2\mu -(s_\infty + d) \vartheta \Leftrightarrow (s_\infty + d) \vartheta \stackrel{!}{<} 2\mu 
\Leftrightarrow \vartheta &\stackrel{!}{<} \frac{2\mu}{(s_\infty + d)}
\end{align*}
Thus for any $\vartheta < 2\mu / (s_\infty + d)$ we obtain a limiting element $f_\vartheta \in H_\vartheta$,
such that $s_{f, X_n} \rightarrow f_\vartheta$ in $\Vert \cdot \Vert_{\Ht}$ for $n \rightarrow \infty$ by following the argumentation used in \cite{schaback2002inverse}.
\begin{itemize}
\item If $2\mu > s_{\infty} + d$, then $ 2\mu / (s_\infty + d) > 1$ and we can choose $\vartheta = 1$ to conclude $f_\vartheta = f \in \ns$. 
This was the way it was done in \cite{schaback2002inverse}.
\item Here we have $2\mu > s_{\infty}$, thus $ 2\mu / (s_\infty + d) > s_\infty / (s_\infty + d)$. This means $f_\vartheta \in H_{s_\infty / (s_\infty + d) - \varepsilon} \supset H_{s_\infty / (s_\infty + d)} \supset \mathcal{H}_1 = \ns$ for all $\varepsilon > 0$.
\item \tw{It might be interesting to check whether we the smoothness $\tau$ of $\Ht$ is sufficiently hight, i.e.\ $\tau > d/2$ for Sobolev embedding.
On the other hand $f \in \mathcal{C}(\Omega)$ was stated as an assumption.} 
%
%
%
%
%
%
%
%
%
%
\end{itemize}

This limiting element $f_\vartheta \in \Ht$ allows us to use sampling inequalities or anything, which gives us an additional $n^{-1/2} = h_X^{d/2}$ convergence order for $\Vert f_\vartheta - s_{f, X_n} \Vert_{L^2(\Omega)}$ and hopefully also $\Vert f_\vartheta - s_{f, X_n} \Vert_{L^2(X)}$. 
\tw{This needs to be done more thoroughly, check e.g.\ Subsection \textit{2.2.2. Sobolev bounds on discrete q norms} from \cite{narcowich2005sobolev}.}.
However the convergence order is now determined by the smoothness of $f_\vartheta$, 
which is given as $\tau = \vartheta \cdot \frac{s_\infty + d}{2}$, 
thus we can expect:
\begin{align}
\begin{aligned}
\label{eq:additional_order}
\Vert f_\vartheta - s_{f, X_n} \Vert_{L^2(\Omega)} \leq C h_{X_n}^{\vartheta \cdot \frac{s_\infty + d}{2}}, \\
\Vert f_\vartheta - s_{f, X_n} \Vert_{L^2(X)} \leq C h_{X_n}^{\vartheta \cdot \frac{s_\infty + d}{2}},
\end{aligned}
\end{align}
where I'm not that sure about the second one (discrete error).
Using $\vartheta = s_\infty / (s_\infty + d)$ we obtain 
\begin{align*}
\vartheta \cdot \frac{s_\infty + d}{2} = \frac{s_\infty}{2},
\end{align*}
such that Eq.\ \eqref{eq:additional_order} turns into
\begin{align}
\begin{aligned}
\label{eq:additional_order}
\Vert f_\vartheta - s_{f, X_n} \Vert_{L^2(\Omega)} \leq C h_{X_n}^{s_\infty / 2},
%
\end{aligned}
\end{align}
which is however exactly the same rate of convergence which we derive by directly using the $L^\infty(\Omega)$ bound.
The goal is to derive a convergence rate as
\begin{align*}
\frac{s_\infty + d}{2}.
\end{align*}
\tw{I think here we need to dig into more details:
We know already the $\Vert f - s_{f, X} \Vert_{L^\infty(\Omega)} \leq c_f h_X^\mu, \mu > s_\infty/2$ error bound.
Therefore we should be able to expect a faster error decay for $\Vert f - s_{f, X} \Vert_{L^2(\Omega)}$.
However, using the argumentation which follows, we do not obtain any faster rate of convergence!}

~ \\
\hrule
~ \\

\textcolor{red}{The remaining text is not important (it's just how to conclude the proof if we had obtained a faster convergence rate!))}

Now we can go back to the original expression
\begin{align*}
\Vert s_{f, Y} - s_{f, X} \Vert_{\ns}^2 &\leq \Vert A_{Y, \Phi}^{-1} \Vert_{2,2} \Vert s_{f, Y} - s_{f, X} \Vert_{L^2(Y)}^2
\end{align*}
and estimate the right hand side $\Vert \cdot \Vert_{L^2(Y)}^2$ expression in a better way by exploiting the additional $h_X^{d/2}$ convergence order from Eq.\ \eqref{eq:additional_order}:
\begin{align*}
\Vert s_{f, Y} - s_{f, X} \Vert_{L^2(Y)} &\leq \Vert f_\vartheta - s_{f, Y} \Vert_{L^2(Y)} + \Vert f_\vartheta - s_{f, X} \Vert_{L^2(Y)} \\
&\leq C h_Y^{\vartheta \cdot \frac{s_\infty + d}{2}} + C h_X^{\vartheta \cdot \frac{s_\infty + d}{2}} \\
&\leq 2C h_X^{\vartheta \cdot \frac{s_\infty + d}{2}}.
\end{align*}
Using $\vartheta = s_\infty / (s_\infty + d)$ we obtain 
\begin{align*}
\vartheta \cdot \frac{s_\infty + d}{2} = \frac{s_\infty}{2},
\end{align*}
which is however exactly the same rate of convergence which we derive by directly using the $L^\infty(\Omega)$ bound.
The goal is to derive a convergence rate as
\begin{align*}
\frac{s_\infty + d}{2}.
\end{align*}
\end{proof}


\section{Ideas/thoughts collected on 04.12.22}



\begin{enumerate}
\item So far I considered $L^\infty$ estimates and tried to derive statements for the $L^2$ error.
However maybe the reverse direction might work as well and even be easier: 
We start with $L^2$ statements and derive statements on the $L^\infty$ error.
In particular I thought whether it is possible to derivate something like:
Assume the $L^2$ error decays with a given rate. 
Then it follows that the $L^\infty$ error decays with a slower rate.
This means: 
Is it possible to do the proof of \cite{wendland1999meshless} the other way round?
With such statements it might be possible to derive something via contraposition.
\begin{enumerate}
\item The best possible $L^2$ error is lower bounded by the eigenvalues of the kernel. 
Look up this statement and try to get something out of it.
\end{enumerate}
\item I thought whether some flat limit statements could help us.
However they can be only / mostly used to derive estimates on the power function, which are however not helpful.
\item I think a crucial question is: What is the difference between data adapted points and uniformly distributed points?
Because also uniformly distributed points can be data adapted (i.e.\ when $f$-greedy yields uniformly distributed points), 
i.e.\ if $P$-greedy and $f$-greedy provide the same rate of convergence. \\
However in my experience this mostly happens, if the function is quite smoooth.
\begin{enumerate}
\item Can we show that: "If a function is at the boundary of the RKHS, then the corresponding convergence rate (for $P$-greedy points) is given by the convergence rate of the Power function decay?
This makes sense if I think about my Wendland $k=0$ example.
And furthermore it would allow us to use the $d/2$ gap: 
Assume $f_\vartheta \in \Ht$ for some $0 < \vartheta < 1$, but not in a smoother space.
Then the previous conjecture (extended by some escaping the native space and superconvergence statements) provides that we cannot converge as fast as given by the $\mu > s_\infty/2$ assumption, i.e.\ a contradiction.
\item Therefore check again the $L^2$ convergence rates for the Wendland $k=0$ example.
I expect them to nicely follow my assumed conjecture, i.e.\ that the smoothness directly scales with the convergence rate!
It think so far I only analyzed the $L^\infty$ convergence rates.
\end{enumerate}
\item I was wondering whether the $L^\infty$ norm is the intrinsic norm for error estimate statements or whether the $L^2$ space is the correct norm to deal with. 
This is also related to the above points of trying to start with $L^2$ estimates and then derive $L^\infty$ estimates.
\item Lower bound for the $L^2$ approximation rate for some function:
Use power space arguments in conjunction with escaping the native space/continuous superconvergence statements: 
Consider $f \in H_\theta, 0 < \theta < 2$, then consider kernel interpolation using the $k_\theta$ kernel and use the lower bound on the $L^2$ approximation error due to the decay of the kernel eigenvalues - which should be deriveable from the decay of the $k$ eigenvalues.
Then move to the $k$ interpolant and obtain hopefully the same bound there!
\item For any kind of required localization argument, we can maybe make use of my \textit{stability of convergence rates} results:
Here we showed that the convergence rates stay the same when going to subdomains.
Only challenging thing: We do not have that basic assumption of $f \in \ns$.
\end{enumerate}

\subsection{Worst-case bounds on $L^2$ error}


We start by citing an $L^2$ bound on the interpolation error, here taken from \cite[Section 3]{santin2016approximation} (see especially 28\_lower\_bound\_approx):

\begin{theorem}
\label{th:bounds_kolmogorov_width}
Let $\lambda_n$ be the eigenvalues of the kernel integral operator.
Let $S(\ns)$ be the unit ball of $\ns$.
Then for the Kolmogorov $n$-width $d_n$ and the quantity $\kappa_n$ it holds
\begin{align*}
d_n &:= \inf_{V_n \subset L^2, \dim(V_n)=n} \sup_{f \in S(\ns)} \Vert f - \Pi_{L^2, V_n}(f) \Vert_{L^2} = \sqrt{\lambda_{n+1}} \\
\kappa_n &:= \inf_{V_n \subset L^2, \dim(V_n)=n} \sup_{f \in S(\ns)} \Vert f - \Pi_{\ns, V_n}(f) \Vert_{L^2} = \sqrt{\lambda_{n+1}}
\end{align*}
\end{theorem}

Therefore we know that the error of best approximation is related / limited by the eigenvalue decay.

The initial idea was to combine Theorem \ref{th:bounds_kolmogorov_width} with Lemma \ref{lem:bounded_Linfty_convergence}.
However 
\begin{enumerate}
\item Theorem \ref{th:bounds_kolmogorov_width} does not involve the $\ns$-norm of the residual!
Maybe this could be found my closely inspecting the original proofs (mind that we already consider $f \in S(\ns)$, i.e.\ in the unit ball).
\item Theorem \ref{th:bounds_kolmogorov_width} has statements on $\sup_{f \in S(\ns)}$, i.e.\ no statement for a fixed function! 
\end{enumerate}

\subsection{Idea: Escaping native space / superconvergence using lower bounds}

So far, \cite{santin2023continuous} derive or improved upper bounds on several quantities using interpolation spaces and result from interpolation theory.
Can we also use the same arguments and results from interpolation theory, to lower bound some intermediate quantities? \\
We recall
\begin{align*}
\Vert g - \Pi(g) \Vert_{\mathcal{H}_{\theta}} \leq \Vert g - \Pi(g) \Vert_{\mathcal{H}_{\theta_1}}^{\frac{\theta_2 - \theta}{\theta_2 - \theta_1}} \cdot \Vert g - \Pi(g) \Vert_{\mathcal{H}_{\theta_2}}^{\frac{\theta - \theta_1}{\theta_2 - \theta_1}}
\end{align*}

\begin{enumerate}
\item Either we use this standard estimate and lower bound quantities before the interval $(\theta, \theta_2)$ or behind the interval $(\theta_1, \theta)$.
\item Or we use directly lower estimates, provided by considering inverse operators. 
However this might be challenging, as the operator in consideration is a projection and therefore not invertable.
\end{enumerate}

~ \newpage



\section{Removed stuff}

\noindent The following two quantities are of interest:
\begin{enumerate}
\item Bound on continuous $L^2(\Omega)$ norm:
\begin{align}
\label{eq:quantity_cont_L2}
\sup_{X \subset \Omega, h_X / q_X \leq c_0} \Vert f - \Pi_X(f) \Vert_{L^2(\Omega)} \leq c_f h_X^{\mu}
\end{align}
\item Bound on discrete $L^2(\Omega)$ norm:
\begin{align}
\label{eq:quantity_discr_L2}
\sup_{X \subset \Omega, h_X / q_X \leq c_0} \sup_{Y \subset \Omega, h_Y \leq h_X} \Vert f - \Pi_X(f) \Vert_{\ell^2(Y)} \leq c_f h_X^{\mu}
\end{align}
\end{enumerate}
We would like to show:
\begin{align*}
\text{Eq.\ } \eqref{eq:quantity_cont_L2} \Leftrightarrow \text{Eq.\ } \eqref{eq:quantity_discr_L2}.
\end{align*}

\begin{itemize}
\item["$\Leftarrow$"] One should be able to prove this by using a fine grained $Y \subset \Omega$, such that we have a fine discretization of $\Omega$, then applying Monte Carlo like arguments.
\item["$\Rightarrow$"] No clue. 
Some argument would be required that the considered function cannot be to spiky.
This should work due to bounded RKHS norm (or at least some interpolation space norm).
\end{itemize}


~ \newpage 


\section{Miscellaneous}

\subsection{Comments on the requirements within \Cref{sec:new_approach_via_L2}}

In \Cref{sec:new_approach_via_L2} we used the assumption
\begin{align*}
\forall_{X \subset \Omega, h_X/q_X \leq c_0}
\end{align*}
frequently, and here I want to comment on two points regarding this:
\begin{enumerate}
\item We used (asymptotically) uniformly distributed points, measured in terms of the ratio $h_X/q_X \leq c_0$: \\
If we lift that assumption, we could also consider $f$-greedy points.
However, simple examples using the Brownian bridge kernel on the scale of test functions $f_\alpha(x) = x^\alpha \cdot (1-x)^\alpha$ reveals,
that $f$-greedy always provides an $L^\infty(\Omega)$ convergence rate of $n^{-2}$, independent of the value of $\alpha$, 
and therefore also for superconvergence and escaping the native space functions.
In particular this $L^\infty(\Omega)$ convergence rates translates into a $L^2(\Omega)$ convergence rate, which is already the best possible convergence rate \cite[Section 6.2]{wenzel2023analysis}.
This reveals that for $f$-greedy selected points it is not possible to conclude the smoothness of the underlying function for either 
I assume that the convergence rate stays the same also for weak $f$-greedy points, i.e.\ it's more about the asymptotic distribution.
\item We used that the convergence rate needs to hold \textit{for all} such (asymptotically) uniform point distributions: \\
Also this assumption seems to be crucial when we think about functions $f^* = \sum_{j=1}^N \alpha_j k(\cdot, x_j^*)$ for a fixed set of centers $X^* := \{ x_1^*, ..., x_N^*\} \subset \Omega$.
If we pick $X \supset X^*$, then our interpolant is exact and we even obtain an error of zero, i.e.\ we have even any convergence rate.
However, as we take the supremum over any set of (asymptotically) uniformly distributed points, we obtain also not-exact approximations and therefore some convergence rate. \\
This exact convergence for $f^*$ likely depends on the discussion of the next subsection, i.e.\ on the power space $\HT$ where $k(\cdot, x)$ belongs to (for all $x \in X^*$).
\end{enumerate}


\subsection{$k(\cdot, x) \in \mathcal{H}^\Theta$ for some $\Theta > 0$?}
\label{subsec:kernel_transl_smoother}

From time to time I noticed that the kernel function $k(\cdot, x) \in \ns$ for fixed $x \in \Omega$ is usually smoother than the RKHS $\ns$. 
This raises the question how much additional smoothness we usually have. \\
Under some assumptions, we can decompose the kernel as
\begin{align*}
k(x, y) = \sum_{j=1}^\infty \lambda_j \varphi_j(x) \varphi_j(y),
\end{align*}
where $\{ \varphi_j \}_{j=1}^\infty$ is a orthonormal basis in $L^2(\Omega)$ and it holds $\Vert \varphi_j \Vert_{\ns}^2 = \frac{1}{\lambda_j}$ for all $j = 1, 2, ...$.
Then the power spaces are given by
\begin{align*}
\mathcal{H}^\Theta = \left\{ f \in L^2(\Omega) ~ | ~ \sum_{j=1}^\infty \frac{|\langle f, \varphi_j \rangle_{L^2(\Omega)}|^2}{\lambda_j^\Theta} < \infty \right\} 
 = \left\{
\begin{array}{ll}
L^2(\Omega), & \Theta = 0 \\
\ns, & \Theta = 1 \\
TL^2(\Omega), & \Theta = 2
\end{array}
\right.
\end{align*}
Now we consider $k(\cdot, x) = \sum_{j=1}^\infty \lambda_j \varphi_j(\cdot) \varphi_j(x)$ and calculate by using the $L^2(\Omega)$ orthonormality
\begin{align*}
\sum_{j=1}^\infty \frac{|\langle k(\cdot, x), \varphi_j \rangle_{L^2(\Omega)}|^2}{\lambda_j^\Theta}
&= \sum_{j=1}^\infty \frac{|\langle \sum_{i=1}^\infty \lambda_i \varphi_i(\cdot) \varphi_i(x), \varphi_j \rangle_{L^2(\Omega)}|^2}{\lambda_j^\Theta} \\
&= \sum_{j=1}^\infty \frac{\lambda_j^2 |\varphi_j(x)|^2}{\lambda_j^\Theta} = \sum_{j=1}^\infty \lambda_j^{2 - \Theta} |\varphi_j(x)|^2 \stackrel{!}{<} \infty.
\end{align*}
\tw{This seems to be equivalent to the question for the minimal $\theta > 0$ such that the power kernel from Eq.~\eqref{eq:power_kernel} is actually a kernel (which is apparently an unsolved question in itself).}
\textcolor{gray}{
This shows that the question for which $\Theta > 1$ we have $k(\cdot, x) \in \mathcal{H}^\Theta$ solely depends on the decay rate of the eigenvalues. \\
Assuming $\lambda_j \asymp j^{-\alpha}$ for some $\alpha > 1$ we have 
\begin{align*}
\sum_{j=1}^\infty \lambda_j^{2 - \Theta} \asymp \sum_{j=1}^\infty j^{-\alpha(2 - \Theta)} \stackrel{!}{<} \infty
\Leftrightarrow \alpha (2-\Theta) \stackrel{!}{>} 1 
\Leftrightarrow \Theta < 2 - \frac{1}{\alpha}.
\end{align*}
Ths shows that for very smooth kernels (i.e.\ large $\alpha \gg 1$), the function $k(\cdot, x)$ is almost contained in the image of the kernel integral operator due to $\Theta \lesssim 2$.
For $1 < \alpha \rightarrow 1$, i.e.\ very unsmooth kernels, the function $k(\cdot, x)$ is is only contained in a very small range of interpolation spaces because  due to $\Theta \gtrsim 1$. \\
I discussed this once with Daniel Winkle, see 12.08.22\textbackslash A2.}

\subsection{Interesting references}

\begin{itemize}
\item "STABILITY ON INTERPOLATION OF SCATTERED DATA VIA KERNELS" might be worth citing, despite it is not too much related. $\rightarrow$ If I recall correctly, this paper has several flaws, mistakes, and non-sharp statements.
\item Check papers from Toni Karvonen: \url{https://arxiv.org/pdf/2203.05400.pdf} or \url{https://arxiv.org/pdf/2001.10965.pdf}
\end{itemize}


\section{Removed text snippets}

I might be possible to leverage results from escaping the native space. Idea:

\begin{enumerate}
\item When we assume in Theorem \ref{th:inverse_statement} only a decay rate of $h_X^\mu$ with $2\mu > s_\infty$ (instead of $> s_\infty + d$),
then it might be possible to show that our function is included in some interpolation space $\Ht$ for some $0 < \vartheta < 1$. \\
The only tricky part hereby might be the change from the kernel $k$ to the kernel $k_\vartheta$, which however might work by inserting a zero
\item Now working in the space $\Ht$ we should be able to conclude that discrete $L^2$ error from Eq.\ \eqref{eq:discrete_L2_error} decays actually faster.
Here we should be able to leverage statements from Sobolev spaces, because for $0 < \vartheta < 1$ the interpolation space $\Ht$ is norm equivalent to a Sobolev space with reduced smoothness.
\item This faster convergence rate for the discrete $L^2$ error might then help us to conclude that we are even contained in the RKHS!
\end{enumerate}

~ \\
\hrule
~ \\

\noindent More details: 

\begin{itemize}
\item First approach:
\begin{align*}
\Vert s_{f, Y} - s_{f, X} \Vert_{\Ht}^2 \leq (?) \cdot \Vert s_{f, Y} - s_{f, X} \Vert_{L^2(X)}^2,
\end{align*}
where $(?)$ needs to be investigated: 
I guess this is some kernel matrix of the interpolation kernel $k_\theta$.
In that case one could use known lower bounds on the smallest eigenvalue, because everything should only depend on the smoothness on the corresponding RKHS, which should be norm equivalent to a Sobolev space of smaller smoothness. \\
\end{itemize}


\section{Some figures}

\begin{figure}[h]
\centering
\setlength\fwidth{.75\textwidth}
\input{Figures/d-2-gap_fig8.tex}
\input{Figures/d-2-gap_fig9.tex}
\setlength\fwidth{.41\textwidth}
\input{Figures/d-2-gap_fig5.tex}
%
\input{Figures/d-2-gap_fig7.tex}
\caption{
Visualization of different quantities ($y$-axis) over the number of points in $X_n$ ($x$-axis): \newline
Top: Visualization of the \textit{discrete $L^\infty$ error} according to Eq.\ \eqref{eq:discrete_Linfty_error}.
Middle: Visualization of the \textit{discrete $L^2$ error} according to Eq.\ \eqref{eq:discrete_L2_error}.
Bottom left: Visualization of the estimate from Eq.\ \eqref{eq:nonsharp_estimate}. The dashed line shows Eq.\ \eqref{eq:squared_rkhs_norm_diff}, and the solid line shows the upper bound \eqref{eq:upper_estimate}.
Bottom right: Visualization of the ratio of Eq.\ \eqref{eq:ratio_quantities}.
}
\label{fig:brownian_bridge_example}
\end{figure}


\section{Removed stuff 2}


\begin{theorem}
\label{th:lower_bound_1}
Let $f \in \ns$ but $f \notin \Ht$ for any $\vartheta > 1$.
Then for any $c_0 > 2$ and $\epsilon > 0$ there exists a $C > 0$ such that it holds 
\begin{align*}
\sup_{X \subset \Omega, h_X / q_X \leq c_0} \Vert f - \Pi_X(f) \Vert_{L^2(\Omega)} \geq C h_X^{\frac{s_\infty + d}{2} + \epsilon}
\end{align*}
\end{theorem}

\noindent Note that we do not need the $\Vert f - \Pi_X(f) \Vert_{\ns}$ norm on the right hand side!

\begin{proof}
We distinguish two cases:
\begin{enumerate}
\item Let $f$ be a finite sum of kernel translates with centers $x_j \in \overline{\Omega}$, 
i.e.\ $f = \sum_{j=1}^n \alpha_j k(\cdot, x_j)$. 
Then, due to \Cref{subsec:kernel_transl_smoother} we have that $k(\cdot, x) \in \Ht$ for some $\vartheta > 1$, and thus $f \in \Ht$.
However then \Cref{th:L2_conv_rate} provides a convergence rate of $\vartheta \cdot \left( \frac{s_\infty + d}{2} \right)$, which is a contradiction to the assumption "for any $\epsilon > 0$".
\tw{I had to fix something in \Cref{subsec:kernel_transl_smoother}, i.e.\ it does not work like this here! 
I.e.\ we might need another approach here!}
\item Let $f$ not be a finite sum of kernel translates with centers from $\overline{\Omega}$.
Proof via contradiction: Assume the opposite:
\tw{This is not the opposite! The opposite is $\forall X \subset \Omega$!}
\begin{align*}
\exists_{\epsilon > 0} \exists_{c_0 > 2} \forall_{C>0} \exists_{X \subset \Omega, h_X/q_X \leq c_0} \Vert f - \Pi_X(f) \Vert_{L^2(\Omega)} \leq C h_X^{\frac{s_\infty + d}{2} + \epsilon}.
\end{align*}
Now we pick $C = C_j = 1/j$ and each obtain a corrensponding set $X = X_j \subset \Omega$.
Then $\{ |X_j| \}_{j \in \N}$ must be unbounded:
\begin{itemize}
\item Assume $\{ |X_j| \}_{j \in \N}$ is bounded, then (due to $C_j \rightarrow 0$) there is an arbitrary accurate approximation of $f$ via $\Pi_{X_{j^*}}(f)$ in the $L^2(\Omega)$ norm with finitely many centers $X_{j^*}$, which means that $f$ itself is a finite combination of kernel translates centered in $\overline{\Omega}$. 
%
This however disagrees with our case distinction.
\end{itemize}
Now (possibly by considering a subsequence) we can assume that $|X_j|$ is monotonically increasing.
Furthermore (possibly by considering another subsequence) we can assume that $h_{X_j}  \asymp 2^{-j}$ (or even faster).
Especially we recall that $q_{X_j} \asymp h_{X_j}$ because the existence of the sets $X_j$ came along with the property of $h_X / q_X \leq c_0$. \\
By our contradiction we have a $\Vert f - \Pi_{X_j} \Vert_{L^2(\Omega)}$ bound, which turns into a $\Vert f - \Pi_{X_j} \Vert_{L^2(Y)}$ bound via \Cref{prop:disc_L2_from_cont_L2}.
Now we can put those sets into the improved proof of Theorem 6.1 of \cite{schaback2002inverse} (see Theorem \ref{th:inverse_statement_improved}),
to show that we are included in a interpolation space for $\vartheta > 1$.
This is a contradiction to the assumption! \\
More details (on the last part) at 05.12.22\textbackslash B2.
\end{enumerate}
\end{proof}

\noindent Now we want to modify the assumption of \Cref{th:lower_bound_1} by modifying $f \in \ns$ to $f \in \Ht$ for some $0 < \vartheta < 2$:

\begin{theorem}
\label{th:lower_bound_2}
Consider $0 < \vartheta < 2$.
Let $f \in \Ht$ but $f \notin \mathcal{H}_{\vartheta_2}$ for any $\vartheta_2 > \vartheta$.
Then for any $c_0 > 2$ and $\epsilon > 0$ there exists a $C > 0$ such that it holds 
\begin{align*}
\sup_{X \subset \Omega, h_X / q_X \leq c_0} \Vert f - \Pi_X(f) \Vert_{L^2(\Omega)} \geq C h_X^{\vartheta \cdot \frac{s_\infty + d}{2} + \epsilon}
\end{align*}
\end{theorem}

Mind that $\Pi_X$ is not actually a projection for $0 < \vartheta < 1$, but we still use it to denote the kernel interpolation based on the interpolation points $X$.

\begin{proof}
\textcolor{gray}{I think this should be possible by leveraging \Cref{th:lower_bound_1} in conjunction with theory of power spaces: \\
We start by assuming the contradiction:
\begin{align*}
\exists_{\epsilon > 0} \exists_{c_0 > 2} \forall_{C>0} \exists_{X \subset \Omega, h_X/q_X \leq c_0} \Vert f - \Pi_X(f) \Vert_{L^2(\Omega)} \leq C h_X^{\Theta \cdot \frac{s_\infty + d}{2} + \epsilon}.
\end{align*}
\tw{Elaborate here! Estimate on discrete $L^2$ norm is now provided via \Cref{prop:disc_L2_from_cont_L2}.}
Now with the same arguments as in the proof of \Cref{th:lower_bound_1}, we can consider an increasing set of interpolation points $\{ X_j \}_{j \in \N}$. Then we recall \Cref{eq:important_estimate}:
\begin{align*}
\Vert s_{f, Y} - s_{f, X} \Vert_{\Ht}^2 &\leq \Vert s_{f, Y} - s_{f, X} \Vert_{L^2(\Omega)}^{2 - 2\vartheta} \cdot \Vert A_{Y, \Phi}^{-1} \Vert_{2,2}^{\vartheta} \cdot \Vert s_{f, Y} - s_{f, X} \Vert_{L^2(Y)}^{2\vartheta}.
\end{align*}
Applying the assumed convergence rates from before and following the remaining parts of the proof of Theorem 6.1 from \cite{schaback2002inverse}, 
we should be able to obtain a limiting element in $\mathcal{H}_\vartheta$ with $\vartheta > \Theta$, 
which is the desired contradiction.}
\end{proof}



\section{Removed stuff 3}

\subsubsection{Some extra calculation}

\tw{Some extra calculation to check for a lower bound}

\begin{align*}
\Vert s_X \Vert_{\ns}^2 &= \alpha^\top A_{X} \alpha = \alpha^\top A_{X} A_{X}^{-1} A_{X} \alpha 
\geq \lambda_{\min}(A_X^{-1}) \cdot \Vert A_X \alpha \Vert_{2}^2 \\
&= \lambda_{\max}(A_X)^{-1} \cdot \Vert A_X \alpha \Vert_{2}^2 \\
& = \lambda_{\max}(A_X)^{-1} \cdot |X| \cdot \Vert s_X \Vert_{L^2(X)}^2
\end{align*}

Together with the statement from \Cref{prop:upper_bound_ns_norm_via_L2} we obtain the bounds

\begin{align*}
\lambda_{\max}(A_X)^{-1} \cdot |X| \cdot \Vert s_X \Vert_{L^2(X)}^2 \leq \Vert s_X \Vert_{\ns}^2 \leq \lambda_{\min}(A_X)^{-1} \cdot |X| \cdot \Vert s_X \Vert_{L^2(X)}^2
\end{align*}

However there is quite a gap between the lower and the upper bound.
As the analysis \Cref{th:inverse_statement_generalized} seems to be quite well, the upper bound rather seems to be frequently sharp: 
Likely the coefficient vector is (at least a bit) aligned with the eigenvector to the smallest eigenvalue.

\else
\fi





\newpage



\end{document} 
